% ============================================================================
% Fichier     : SNSD_main.tex
% Titre       : Stratégie numérique et schéma directeur SI
% UE          : MSN-2S4-MSI01 (3 ECTS, 30 h)
% Auteur      : MINKA MI NGUIDJOÏ Thierry Emmanuel
% Version     : 1.0 -- novembre 2026
% Compilation : pdflatex, biber, makeindex, pdflatex (deux fois)
% ============================================================================

\documentclass[12pt, twoside, openright, a4paper]{book}

\usepackage{Style_SNSD}
\raggedbottom

\addbibresource{bibliographie/SNSD_references.bib}

\hypersetup{
    pdftitle={Stratégie numérique et schéma directeur SI},
    pdfauthor={MINKA MI NGUIDJOÏ Thierry Emmanuel},
    pdfsubject={MSN-2S4-MSI01 · Master Professionnel en Management des Systèmes d'Information · Chaire UNESCO d'Afrique Centrale Accès TIC},
    pdfkeywords={stratégie numérique, schéma directeur, alignement stratégique, portefeuille de projets, ISO/IEC 38500, COBIT 2019, SND30, marchés publics, NEK}
}

\begin{document}
\emergencystretch=3em

\frontmatter

    % ==============================================================================
% Fichier : front/00_page_titre.tex
% Titre   : Page de titre - Stratégie numérique et schéma directeur SI
% Auteur  : MINKA MI NGUIDJOÏ Thierry Emmanuel
% Version : 1.0 -- novembre 2026
% ==============================================================================

\begin{titlepage}
\thispagestyle{empty}

% --- Fond décoratif : bandeaux bleu marine et filets or ---
\begin{tikzpicture}[remember picture, overlay]
    \fill[CouleurPrimaire]
        (current page.north west) rectangle
        ([yshift=-3.6cm] current page.north east);
    \fill[CouleurAccent]
        ([yshift=-3.6cm] current page.north west) rectangle
        ([yshift=-4.0cm] current page.north east);
    \fill[CouleurPrimaire]
        (current page.south west) rectangle
        ([yshift=2.5cm] current page.south east);
    \fill[CouleurAccent]
        ([yshift=2.5cm] current page.south west) rectangle
        ([yshift=2.9cm] current page.south east);
    \node[align=center, text=white] at ([yshift=-1.8cm] current page.north) {%
        {\fontsize{11}{13}\selectfont
        \textsc{République du Cameroun}}\\[0.3ex]
        {\color{CouleurAccent!60!white}\rule{0.4\paperwidth}{0.4pt}}\\[0.6ex]
        {\fontsize{10}{12}\selectfont\textbf{\institutionprimaire}}\\[0.2ex]
        {\fontsize{8}{10}\selectfont \tutelle}\\[0.2ex]
        {\fontsize{9}{11}\selectfont \unitepedago}};
    \node[fill=white, rounded corners=4pt, inner sep=2pt]
        at ([xshift=1.75cm, yshift=-1.8cm] current page.north west)
        {\includegraphics[height=2.6cm]{figures/logos/enspy_logo.jpg}};
    \node[fill=white, rounded corners=4pt, inner sep=2pt]
        at ([xshift=-1.75cm, yshift=-1.8cm] current page.north east)
        {\includegraphics[height=2.6cm]{figures/logos/unesco_logo.jpeg}};
    \node[align=center] at ([yshift=1.25cm] current page.south) {%
        {\color{white}\bfseries\large Année académique \anneescolaire}\\[0.3ex]
        {\color{CouleurAccent!70!white}\small
        Version 1.0 | Document pédagogique de l'UE \codeUE}};
\end{tikzpicture}

\vspace*{0.6cm}

\begin{center}
    {\color{CouleurPrimaire}\rule{0.85\textwidth}{1.2pt}}\\[1.2ex]
    {\fontsize{26}{32}\selectfont\bfseries\color{CouleurPrimaire}
        Stratégie numérique\\[0.4ex]
        et schéma directeur SI}\\[0.8ex]
    {\color{CouleurAccent}\rule{0.85\textwidth}{2pt}}\\[0.5ex]
    {\color{CouleurPrimaire}\rule{0.85\textwidth}{0.5pt}}
\end{center}

\vspace{0.8cm}

\begin{center}
    {\fontsize{13}{16}\selectfont\itshape\color{CouleurSecondaire}
        Aligner, financer, gouverner~:\\
        du diagnostic à la feuille de route d'une organisation camerounaise}
\end{center}

\vspace{0.5cm}

\begin{center}
    \begin{tikzpicture}
        \node[draw=CouleurPrimaire, fill=CouleurPrimaire, rounded corners=6pt,
              inner xsep=14pt, inner ysep=6pt]{%
            \color{white}\bfseries\large \sigleCours{} - UE~\codeUE
        };
    \end{tikzpicture}
\end{center}

\vspace{1.0cm}

% --- Schéma de la discipline : le cycle d'un schéma directeur ---
\begin{center}
\begin{tikzpicture}[
    etape/.style = {draw=CouleurBordDef, fill=CouleurDef, rounded corners=3pt,
        minimum width=2.5cm, minimum height=0.8cm, align=center,
        font=\small\bfseries\color{CouleurPrimaire}, inner sep=2pt},
    arrow/.style = {-{Stealth[length=4pt]}, thick, color=CouleurSecondaire}
]
\node[etape] (n1) at (0, 0)      {Cadrer};
\node[etape] (n2) at (3.2, 0)    {Diagnostiquer};
\node[etape] (n3) at (6.4, 0)    {Fixer\\la cible};
\node[etape] (n4) at (6.4, -1.4) {Arbitrer};
\node[etape] (n5) at (3.2, -1.4) {Financer};
\node[etape] (n6) at (0, -1.4)   {Piloter};
\draw[arrow] (n1) -- (n2); \draw[arrow] (n2) -- (n3); \draw[arrow] (n3) -- (n4);
\draw[arrow] (n4) -- (n5); \draw[arrow] (n5) -- (n6); \draw[arrow] (n6) -- (n1);
\node[draw=CouleurPrimaire, fill=CouleurPrimaire, rounded corners=4pt,
      align=center, font=\footnotesize\bfseries\color{white}, inner xsep=8pt]
      at (3.2, -2.8)
      {Trois preuves~: aligné, finançable, gouverné. Chaque constat~: \OUI{} / \NON{} / \NEK{}};
\end{tikzpicture}
\end{center}

\vspace{0.6cm}

\begin{center}
    {\large\bfseries\color{CouleurPrimaire}
        \auteurcours}\\[0.5ex]
    {\normalsize\color{CouleurSecondaire}
       Enseignant Chercheur | \laboratoire{} | ENSPY}\\[0.2ex]
    {\small\color{CouleurBordInfo}
        \institutionprimaire}\\[0.2ex]
    {\small\color{CouleurBordInfo}
        \href{mailto:maletyon@proton.me}{maletyon@proton.me}}
\end{center}

\end{titlepage}

    \cleardoublepage

    % ============================================================================
% Fichier  : front/01_avant_propos.tex
% Version  : 1.0 -- novembre 2026
% ============================================================================

\chapter*{Avant-propos}
\addcontentsline{toc}{chapter}{Avant-propos}
\markboth{Avant-propos}{Avant-propos}

Presque toutes les organisations camerounaises de quelque taille ont un jour
commandé un schéma directeur. Beaucoup dorment dans un tiroir. Le document
existe, il est relié, il a été payé. Personne ne l'a ouvert depuis sa
signature. Les projets qui se font ne sont pas ceux qu'il annonçait.

Ce cours part de ce constat. Un schéma directeur échoue rarement par sa
technique. Il échoue parce qu'il ne dit pas à quelle stratégie il répond,
parce qu'il n'est pas financé ou parce que personne ne le pilote après sa
validation. Le lecteur apprendra donc à produire un document qui tient sur
trois preuves~: il est aligné sur la stratégie de l'organisation, il est
finançable dans le cadre budgétaire qui s'impose à elle, il est gouverné par
des instances qui se réunissent et qui décident.

Le cours clôt la spécialité Management des systèmes d'information. Il
suppose acquis la gouvernance stratégique des SI (UE MSN-2S3-MSI01),
l'architecture d'entreprise et l'urbanisation (UE MSN-2S3-MSI02), COBIT et
ITIL (UE MSN-2S3-MSI03), la gestion de projet (UE MSN-2S3-MGT01) et l'audit
des SI (UE MSN-2S3-MSI05). Il s'en sert sans les répéter. Il laisse au cours
de conduite du changement (UE MSN-2S4-MSI04) l'adhésion des personnes.

Trois partis pris le traversent. Le premier est de lire les théories dans
leurs textes d'origine, avec leurs dates et leurs limites, parce qu'un modèle
cité de seconde main devient vite un slogan. Le deuxième est de placer le
droit avant la méthode~: une administration n'achète pas, ne finance pas et
ne décide pas comme une entreprise, et un schéma directeur qui l'ignore est
inapplicable dès sa première ligne budgétaire. Le troisième est de coter ce
que l'on sait. Chaque constat du diagnostic et chaque exigence d'audit reçoit
une cote \OUI{}, \NON{} ou \NEK{}, appuyée sur une pièce. Un diagnostic qui
ne distingue pas ce qui est établi de ce qui est supposé fonde une cible sur
du sable.

Avant la méthode, deux parties situent le sujet. La première présente le
contexte international~: les normes et les référentiels, les doctrines
publiques, les stratégies numériques de l'Afrique. La seconde présente le
contexte national~: la stratégie numérique du Cameroun et le cadre juridique
qui s'impose à un schéma directeur. Les textes nationaux sont souvent
l'internalisation de règles venues d'ailleurs, d'où cet ordre.

Les trois organisations qui accompagnent le lecteur, l'ANPP, LiZbiZ et la
BANQUE SANAGA, sont fictives, comme leurs dirigeants, leurs projets, les
dates et les montants. Le cabinet AUDIX CONSEIL l'est aussi. Les théories,
les normes et les textes cités sont exacts. Les normes sous licence sont
paraphrasées.

Le livre est accompagné d'un livret des corrigés et de NDAP SUKULU, où les
fiches de questions et les dix laboratoires du cours permettent de pratiquer
chaque chapitre.

\bigskip
\hfill\begin{minipage}{0.5\textwidth}
\raggedleft
\IfFileExists{../../Logos/signature_avant_propos.png}{\includegraphics[height=1.6cm]{../../Logos/signature_avant_propos.png}}{\textit{MINKA MI NGUIDJOÏ Thierry Emmanuel}}\\
Yaoundé, novembre 2026
\end{minipage}

    \cleardoublepage

    % ============================================================================
% Fichier  : front/02_mode_emploi.tex (Mode d'emploi)
% Version  : 1.0 -- novembre 2026
% ============================================================================
% 30 h de cours magistral = 15 séances de 2 h ; TD (fiches, exercices) et TP
% (dix labs) faits par l'étudiant chez lui, hors des 30 h.
% ============================================================================

\chapter*{Mode d'emploi du cours}
\addcontentsline{toc}{chapter}{Mode d'emploi du cours}
\markboth{Mode d'emploi}{Mode d'emploi}

\section*{Structure générale}

Le cours compte un prologue et six parties, soit seize chapitres, pour
30~heures de cours magistral réparties en quinze séances de deux heures.
Les travaux dirigés et les travaux pratiques se font à la maison, hors de
ces 30~heures~: pour les travaux dirigés, la fiche de questions de chaque
chapitre et ses exercices de fin de chapitre~; pour les travaux pratiques,
les dix laboratoires de NDAP SUKULU, d'une heure chacun. Les corrigés des
exercices sont réunis dans un livret séparé.

\begin{tableaucours}{cL{6.4cm}cL{3.2cm}}{%
    \tableauHeader{Partie} &
    \tableauHeader{Thème} &
    \tableauHeader{Chap.} &
    \tableauHeader{Niveau}
}
0 & Prologue~: le schéma directeur dort dans un tiroir & P & introductif \\
I & Fondements~: stratégie, alignement, valeur & 1--3 & fondamental \\
II & Le contexte international & 4--5 & fondamental \\
III & Le contexte national & 6--7 & fondamental \\
IV & Méthode~: élaborer le schéma directeur & 8--12 & intermédiaire \\
V & Mise en œuvre~: piloter et maîtriser & 13--14 & intermédiaire \\
VI & Défis et recherche & 15--16 & expert à recherche \\
\end{tableaucours}

\section*{Les quinze séances}

\begin{tableaucours}{cL{6.8cm}L{4.2cm}}{%
    \tableauHeader{Séance} &
    \tableauHeader{Contenu} &
    \tableauHeader{Travail à la maison}
}
1 & Prologue et Chap.~1~: stratégie, SI, numérique & Fiches P et 1 \\
2 & Chap.~2~: les théories de la stratégie & Fiche 2, lab 01 \\
3 & Chap.~3~: l'alignement et la valeur & Fiche 3, lab 02 \\
4 & Chap.~4~: normes et référentiels & Fiche 4 \\
5 & Chap.~5~: doctrines et stratégies numériques & Fiche 5 \\
6 & Chap.~6 et~7~: stratégie numérique du Cameroun, cadre juridique &
Fiches 6 et 7, lab 03 \\
7 & Chap.~8~: cadrer et gouverner la démarche & Fiche 8 \\
8 & Chap.~9~: diagnostiquer l'existant & Fiche 9, lab 04 \\
9 & Chap.~10~: fixer la cible & Fiche 10, lab 05 \\
10 & Chap.~11~: portefeuille et trajectoire & Fiche 11, lab 06 \\
11 & Chap.~12~: chiffrer, financer, acheter, doter & Fiche 12, lab 07 \\
12 & Chap.~13~: piloter l'exécution & Fiche 13, lab 08 \\
13 & Chap.~14~: risques, sécurité, continuité, conformité & Fiche 14,
lab 09 \\
14 & Chap.~15~: données, IA, plateformes, sobriété & Fiche 15 \\
15 & Chap.~16~: auditer un schéma directeur & Fiche 16, lab 10 \\
\end{tableaucours}

Cette répartition est indicative~; l'enseignant l'adapte au calendrier et au
rythme du groupe.

\section*{Le gabarit des chapitres}

Chaque chapitre suit le même plan~: une épigraphe, le niveau et les
prérequis, les objectifs, une situation d'ouverture tirée de l'un des trois
cas, des sections numérotées, ce qu'il faut retenir, des points de
contrôle, trois à quatre exercices gradués, des lectures et une
transition.

\section*{Les encadrés}

\begin{tableaucours}{L{3.6cm}L{8.6cm}}{%
    \tableauHeader{Encadré} &
    \tableauHeader{Rôle}
}
Objectifs & Ce que l'on doit savoir faire à la fin du chapitre \\
Définition & Les notions, en gras dans le texte et indexées \\
Violation typique & Une erreur fréquente et la règle qui l'évite \\
Information & Un complément, un renvoi, une nuance \\
Mission (ANPP, LiZbiZ, BANQUE SANAGA) & Le cas~: les constatations, puis
l'analyse, avec un niveau de confiance \\
Démarche & Les étapes d'une procédure, dans l'ordre \\
Lecture & Un décompte de cotes \OUI{}~/ \NON{}~/ \NEK{} et sa lecture \\
Métier & Ce que la notion change pour chacun des métiers du cours \\
Texte juridique & Un texte cité avec sa référence et sa portée \\
Ce qu'il faut retenir & Les faits et les règles du chapitre \\
Exercice & Trois niveaux de difficulté~: \niveaubase, \niveaumoyen,
\niveauexpert \\
\end{tableaucours}

\section*{Les conventions}

\begin{itemize}
    \item Les cotes~: \OUI{} (fait établi par une pièce), \NON{} (écart
          établi par une pièce), \NEK{} (connaissance insuffisante pour
          trancher). Un \NEK{} est de preuve quand une pièce manque, de
          critère quand la référence manque.
    \item Constatation et analyse sont toujours séparées~; l'analyse porte
          un niveau de confiance (faible, moyenne, élevée).
    \item Les textes~: chaque loi, décret ou règlement est cité avec son
          numéro et sa date~; les normes sous licence sont paraphrasées~;
          la source de chaque texte figure à l'annexe~\ref{ann:textes}.
    \item Les montants sont en francs CFA. Un « milliard » vaut mille
          millions.
    \item Les cas sont fictifs~: noms en majuscules, villes réelles.
\end{itemize}

    \cleardoublepage

    % ==============================================================================
% Fichier : front/03_ndap_sukulu.tex
% Titre   : NDAP SUKULU, l'environnement compagnon du cours
% Auteur  : MINKA MI NGUIDJOÏ Thierry Emmanuel
% Version : 1.0 -- novembre 2026
% ==============================================================================
% Texte commun aux cours de l'auteur, mis à jour en novembre 2026 (version en
% ligne, banques de 100 questions, compte de démonstration) ; le tableau des
% labs est propre au cours (front/03b_labs.tex).
% ==============================================================================

\chapter*{NDAP SUKULU, la salle de classe compagnon}
\addcontentsline{toc}{chapter}{NDAP SUKULU, la salle de classe compagnon}
\markboth{NDAP SUKULU}{NDAP SUKULU}

\begin{center}
    \includegraphics[width=0.62\textwidth]{figures/logos/logo_ndap_sukulu.png}
\end{center}

\section*{Un nom venu du pays bassa}

\begin{quote}\itshape
Le soir tombe sur Éséka. Un enfant revient des champs avec son grand-père et
lui demande pourquoi il doit marcher si loin, chaque matin, pour aller à
l'école. Le vieil homme s'arrête au bord du chemin et répond, en bassa~:
« \textup{Ndap sukulu}, la salle de classe, ce ne sont pas les murs. Ce sont
les murs qui viennent après. La salle de classe, c'est l'endroit où quelqu'un
t'attend pour t'apprendre ce qu'il sait. » L'enfant insiste~: « Et qui
m'attend~? » Le grand-père sourit~: « \textup{Malet yon}. Ton enseignant. »
\end{quote}

En bassa, la langue paternelle de l'auteur, \textit{ndap sukulu} désigne la
salle de classe et \textit{malet yon} signifie « ton enseignant ».
L'environnement compagnon de ce cours porte ce nom parce qu'il prolonge la
salle de classe au-delà des heures de cours~: il attend l'étudiant le soir,
le week-end, à distance, avec des questions et des pièces à examiner. Il
est édité par MaletYon et mis à la disposition des étudiants et des
praticiens qui souhaitent se perfectionner
(\href{mailto:maletyon@proton.me}{maletyon@proton.me}).

\section*{Ce que l'on y trouve}

\begin{itemize}
    \item Une fiche de questions par chapitre. Chaque tentative tire vingt
          questions d'une banque de cent (QCM, vrai ou faux, textes à
          trous) et dure trente minutes.
    \item Dix laboratoires. Chaque lab place l'étudiant devant les pièces
          d'une mission d'audit du cours (registres, exports, journaux
          d'accès, contrats, tableaux de flux) et se termine par des points
          de contrôle.
    \item Un espace enseignant, qui corrige les remises, établit la
          synthèse de la classe et publie les résultats.
\end{itemize}

Les fiches et les labs sont des travaux à faire hors des heures de cours.
Ils ne remplacent ni le livre ni les exercices de fin de chapitre~: ils en
vérifient la maîtrise.

\section*{Comment y accéder}

\begin{enumerate}
    \item En ligne, à l'adresse
          \url{https://maletyon.github.io/cours-snsd/ndap-sukulu/}, dans un
          navigateur récent~; sur téléphone, ajouter \texttt{telephone/} à
          la fin de l'adresse.
    \item Hors ligne, en téléchargeant l'archive de l'édition ordinateur
          depuis le dépôt \texttt{MaletYon/cours-snsd}, dossier
          \texttt{ndap-sukulu}.
    \item Se connecter avec le matricule et le mot de passe remis par
          l'enseignant. Pour essayer avant d'être inscrit, un compte de
          démonstration est ouvert à tous~: matricule \texttt{ETU001}, mot
          de passe \texttt{eTU001}~; ses essais ne sont jamais transmis à
          l'enseignant.
    \item Pour un lab, vérifier d'abord l'empreinte SHA-256 des pièces
          (\texttt{sha256sum -c MANIFESTE\_SHA256.txt}), puis travailler sur
          une copie. Les données personnelles des pièces sont inventées.
\end{enumerate}

\begin{boxinfo}
Les fiches de questions sont aussi fournies au format texte dans le dossier
\texttt{fiches} du dépôt du cours, prêtes à être distribuées. Leurs
corrections figurent dans le livret de correction, à la suite des corrigés
des exercices de chaque chapitre.
\end{boxinfo}

\section*{Les laboratoires du cours}

% Tableau des labs du cours SNSD dans NDAP SUKULU (propre au cours).
Le cours compte dix laboratoires, disponibles sous le code \codeUE. Chacun
place l'étudiant dans l'une des trois missions du cabinet, devant les pièces
qu'un consultant réunit~: inventaires, questionnaires, fiches de projet,
budgets, tableaux de bord.

\begin{tableaucours}{lL{6.6cm}L{1.6cm}L{2.2cm}}{%
    \tableauHeader{Lab} &
    \tableauHeader{Sujet} &
    \tableauHeader{Chapitre} &
    \tableauHeader{Mission}
}
SNSD-LAB-01 & Cinq forces et chaîne de valeur & 2 & LiZbiZ \\
SNSD-LAB-02 & Mesurer la maturité de l'alignement & 3 & ANPP \\
SNSD-LAB-03 & Carte d'applicabilité des textes & 7 & Les trois \\
SNSD-LAB-04 & Inventaire applicatif et grille de McFarlan & 9 & ANPP \\
SNSD-LAB-05 & Comparer trois scénarios de cible & 10 & BANQUE SANAGA \\
SNSD-LAB-06 & Arbitrer un portefeuille sous contrainte & 11 & ANPP \\
SNSD-LAB-07 & Coût complet sur cinq ans & 12 & LiZbiZ \\
SNSD-LAB-08 & Tableau de bord d'exécution & 13 & ANPP \\
SNSD-LAB-09 & Registre des risques et continuité & 14 & BANQUE SANAGA \\
SNSD-LAB-10 & Audit intégrateur d'un schéma directeur & 16 & ANPP \\
\end{tableaucours}
\vspace{-0.8em}
\begin{center}{\small\itshape Les laboratoires du cours dans NDAP SUKULU}\end{center}


    \cleardoublepage

    \tableofcontents
    \cleardoublepage

    \listoftables
    \cleardoublepage

\mainmatter

    \part*{Prologue opérationnel}
    \addcontentsline{toc}{part}{Prologue Opérationnel}

        \setcounter{chapter}{-1}
        % ============================================================================
% Fichier  : partie0/P00_tiroir.tex
% Titre    : Le schéma directeur dort dans un tiroir
% Auteur   : MINKA MI NGUIDJOÏ Thierry Emmanuel
% Version  : 1.0 -- novembre 2026
% Corrigés : corriges/Corr_P00.tex
% ============================================================================

\chapter{Le schéma directeur dort dans un tiroir}
\label{chap:prologue}
\chaptermark{Prologue}

\epigraph{%
    « Un plan que personne ne suit n'est pas un mauvais plan.\\
    C'est un document. »%
}{Principe du cours}

\niveauChapitre{Introductif}{%
    Aucun prérequis
}

\begin{boxobjectifs}
\begin{itemize}
    \item Reconnaître un schéma directeur qui ne sert pas.
    \item Poser les trois questions du cours~: aligné, finançable, gouverné.
    \item Coter une proposition \OUI{}, \NON{} ou \NEK{} avec une pièce.
    \item Situer les trois organisations du cours.
\end{itemize}
\end{boxobjectifs}

\bigskip

% ============================================================================
\section*{Avant de commencer : « nous en avons déjà un »}
% ============================================================================

Yaoundé, mardi 14~janvier~2025. L'Agence nationale de paiement des prestations\footnote{L'ANPP, ses dirigeants,
ses projets et ses chiffres sont fictifs, comme LiZbiZ, la BANQUE SANAGA et le cabinet AUDIX CONSEIL. Ces
organisations figurent dans d'autres cours de l'auteur.}, l'ANPP, verse chaque mois des prestations à
186\,000~bénéficiaires par dix antennes régionales. Elle emploie 412~agents. Son conseil d'administration vient de
commander un schéma directeur pour 2025-2029 au cabinet AUDIX CONSEIL.

À la réunion de lancement, la cheffe de la division informatique, Mme~NGO~BIYONG, pose un classeur sur la
table. « Nous en avons déjà un. Il couvrait 2019 à 2023. Je suis arrivée en 2021 et personne ne m'en a parlé. »
Le directeur général, M.~OWONA, ajoute qu'en mars~2024 une panne a arrêté les paiements pendant quatre jours.
Le projet e-RECETTES, mis en service en 2023, n'avait pas de site de secours. « Le conseil veut un nouveau
document. Je voudrais surtout qu'il serve. »

M.~BEYALA, chef de mission, confie le classeur à Mme~EKANI et au consultant junior. Leur première tâche n'est
pas d'écrire. Elle est de comprendre pourquoi le document précédent n'a servi à rien.

\begin{boxalert}
\textbf{Violation typique.} Commander un nouveau schéma directeur sans avoir fait le bilan du précédent.

\smallskip
\textbf{Règle~:} les causes qui ont laissé dormir un document le laisseront dormir une seconde fois. Le bilan
du schéma directeur échu est la première pièce du nouveau.
\end{boxalert}

% ============================================================================
\section{Ce que contenait le schéma directeur 2019-2023}
\label{sec:p-contenu}
% ============================================================================

Le document compte 214~pages. Il a été rédigé en 2018 par un cabinet, validé par le conseil d'administration le
17~décembre~2018 et payé 86~millions de francs CFA. Mme~EKANI en dresse l'inventaire en une matinée.

\begin{tableaucours}{L{5.2cm}cL{5.6cm}}{%
    \tableauHeader{Rubrique} &
    \tableauHeader{Pages} &
    \tableauHeader{Ce que l'on y trouve}
}
État des lieux & 118 & Inventaire du matériel, poste par poste \\
Orientations & 9 & Sept axes, sans lien avec le plan stratégique de l'Agence \\
Projets & 61 & 23 fiches de projet, sans priorité ni responsable \\
Budget & 4 & Un total de 6,3~milliards, sans échéancier \\
Gouvernance & 2 & Un comité de pilotage, jamais réuni \\
Annexes & 20 & Glossaire, listes \\
\end{tableaucours}
\vspace{-0.8em}
\begin{center}{\small\itshape Tableau~P.1. Anatomie du schéma directeur 2019-2023 de l'ANPP}\end{center}

Plus de la moitié du document décrit du matériel qui a été remplacé depuis. Les orientations tiennent en neuf
pages. Le budget tient en quatre, pour un montant qui représente dix fois le budget informatique annuel de
l'Agence. Aucune ligne ne dit d'où viendrait l'argent.

% ============================================================================
\section{Ce qui s'est passé en cinq ans}
\label{sec:p-bilan}
% ============================================================================

\begin{boxscenario}{Mission ANPP. Le bilan du schéma directeur 2019-2023, 21 janvier 2025}
\textbf{Constatations.}
\begin{itemize}
    \item Sur les 23 projets inscrits, 6 ont été livrés, 4 ont été commencés puis arrêtés, 13 n'ont jamais
          démarré.
    \item Sur la même période, l'Agence a mené 11 projets qui ne figuraient pas au schéma directeur, dont
          e-RECETTES, le plus coûteux.
    \item Le comité de pilotage prévu ne s'est jamais réuni. Aucun procès-verbal n'existe.
    \item Le conseil d'administration n'a reçu aucun point d'avancement après la validation.
    \item Les dépenses informatiques réelles s'élèvent à 3,1~milliards sur cinq ans, soit la moitié du montant
          annoncé.
\end{itemize}

\textbf{Analyse} (confiance élevée sur les chiffres, moyenne sur les causes). Le document n'a pas guidé les
décisions. Les projets réalisés ont répondu à des urgences ou à des financements extérieurs. Personne n'était
chargé de comparer ce qui se faisait à ce qui était prévu.
\end{boxscenario}

M.~BEYALA écrit trois questions au tableau. Elles servent de fil à tout le cours.

\begin{enumerate}
    \item \textbf{Aligné~?} À quelle stratégie de l'Agence ce document répondait-il~?
    \item \textbf{Finançable~?} Avec quel argent, inscrit dans quel budget, engagé selon quelle procédure~?
    \item \textbf{Gouverné~?} Qui décidait, qui suivait, qui rendait compte~?
\end{enumerate}

% ============================================================================
\section{Huit premières propositions}
\label{sec:p-lecture}
% ============================================================================

Avant tout diagnostic, l'équipe cote huit propositions simples sur le document. Chacune reçoit une cote et une
pièce. C'est la convention de tout le cours.

\begin{boxlecture}{Le schéma directeur 2019-2023, 21 janvier 2025}
\begin{tableaucours}{L{7.6cm}cL{3.4cm}}{%
    \tableauHeader{Proposition} &
    \tableauHeader{Cote} &
    \tableauHeader{Preuve}
}
Le document cite le plan stratégique de l'Agence & \NON{} & Aucune mention \\
Chaque projet est rattaché à un objectif de l'Agence & \NON{} & 23 fiches lues \\
Les projets sont classés par priorité & \NON{} & Liste alphabétique \\
Le budget est réparti par année & \NON{} & Total unique \\
Une source de financement est indiquée & \NON{} & Aucune \\
Une instance de pilotage est désignée & \OUI{} & Page 198 \\
Cette instance s'est réunie & \NON{} & Aucun procès-verbal \\
Le conseil a suivi l'exécution & \NEK{} & Procès-verbaux de 2020 non retrouvés \\
\end{tableaucours}

\medskip
\textbf{Décompte.} Un \OUI{}, six \NON{}, un \NEK{} sur huit.
\hfill\lecture{0{,}125}{0{,}125}{1{,}06}

\medskip
\textbf{Lecture.} Le document n'est ni aligné ni finançable. Il désigne une instance, ce qui vaut un \OUI{}, mais
l'instance n'a pas fonctionné. Le \NEK{} est de preuve~: les procès-verbaux du conseil pour 2020 manquent aux
archives. Avec $U = 0{,}125$, l'incertitude ne gêne pas la conclusion. Six écarts sont établis.
\end{boxlecture}

\begin{boxdef}
\textbf{Schéma directeur.} Document qui fixe, pour une durée de trois à cinq ans, la cible du système
d'information d'une organisation, les projets qui y conduisent, leur ordre, leur coût et la manière dont
l'ensemble sera piloté. Il découle de la stratégie de l'organisation. Il n'est ni un inventaire ni un catalogue
de souhaits.
\end{boxdef}

% ============================================================================
\section{Le plan du cours}
\label{sec:p-plan}
% ============================================================================

Le cours suit l'ordre de la mission. Les fondements donnent le vocabulaire et les théories
(Chapitres~\ref{chap:vocabulaire} à~\ref{chap:alignement}). Le contexte international puis national dit ce qui
s'impose à l'Agence avant qu'elle ne décide (Chapitres~\ref{chap:normes} à~\ref{chap:droit}). La méthode
déroule les cinq temps d'un schéma directeur (Chapitres~\ref{chap:cadrer} à~\ref{chap:financer}). La mise en
œuvre traite le pilotage et la maîtrise des risques (Chapitres~\ref{chap:piloter} et~\ref{chap:maitriser}).
Les deux derniers chapitres ouvrent sur les données, l'intelligence artificielle et l'audit du schéma directeur
lui-même.

\begin{boxmetier}{Consultant}
Lire le document précédent avant toute réunion. Compter ce qui a été fait, ce qui ne l'a pas été et ce qui s'est fait en dehors.
\end{boxmetier}

\begin{boxmetier}{Directeur des systèmes d'information}
Garder la trace des décisions. Un projet hors schéma directeur n'est pas une faute. Un projet dont personne ne sait pourquoi il a été lancé en est une.
\end{boxmetier}

\begin{boxmetier}{Dirigeant}
Demander chaque année où en est le schéma directeur. Un document que le dirigeant ne réclame pas n'est suivi par personne.
\end{boxmetier}

\begin{boxresume}
\begin{itemize}
    \item Un schéma directeur fixe une cible, des projets ordonnés, un coût et un pilotage. Il découle de la stratégie.
    \item Trois questions le jugent~: est-il aligné, finançable, gouverné~?
    \item À l'ANPP, 6 projets sur 23 ont été livrés et 11 projets se sont faits hors du document.
    \item Chaque proposition se cote \OUI{}, \NON{} ou \NEK{}, avec une pièce.
    \item Le bilan du document échu ouvre le nouveau.
\end{itemize}
\end{boxresume}

\begin{boxinfo}
\textbf{Points de contrôle.} Vous avez compris ce chapitre si vous savez~:
\begin{itemize}
    \item définir un schéma directeur et dire ce qu'il n'est pas~;
    \item énoncer les trois questions du cours~;
    \item expliquer la différence entre un \NON{} et un \NEK{}.
\end{itemize}
\end{boxinfo}

\begin{boxexo}[title={\ding{45}\quad Exercice~P.1. Le bilan en chiffres \niveaubase}]
Calculez, pour le schéma directeur 2019-2023 de l'ANPP, le taux de projets livrés, le taux de projets hors document parmi les projets menés et le taux d'exécution du budget.
\end{boxexo}

\begin{boxexo}[title={\ding{45}\quad Exercice~P.2. Coter votre organisation \niveaumoyen}]
Prenez le schéma directeur d'une organisation que vous connaissez, ou à défaut son plan informatique. Cotez les huit propositions du prologue. Calculez $\mu$ et $U$. Rédigez la lecture en cinq lignes.
\end{boxexo}

\begin{boxexo}[title={\ding{45}\quad Exercice~P.3. La lettre au conseil \niveauexpert}]
Rédigez la note d'une page que M.~BEYALA adresse au président du conseil, M.~TCHUENTE, pour expliquer pourquoi la mission commence par un bilan. Séparez constatations et analyse.
\end{boxexo}

\bigskip

\begin{boxinfo}
\textbf{Pour aller plus loin.}
\begin{itemize}
    \item Le chapitre sur le schéma directeur dans Peppard et Ward~\cite{wardpeppard2016}.
    \item L'ISO/IEC~38500~\cite{iso38500}, pour le rôle de l'organe de gouvernance.
\end{itemize}
\end{boxinfo}

\begin{boxinfo}
\textbf{Transition.} Avant de reconstruire, il faut des mots exacts. Stratégie, système d'information, numérique, plan, schéma directeur, feuille de route~: le Chapitre~\ref{chap:vocabulaire} les sépare.
\end{boxinfo}


    \part{Fondements : stratégie, alignement, valeur}
    \label{part:fondements}

        % ============================================================================
% Fichier  : partie1/P01_C01_vocabulaire.tex
% Titre    : Stratégie, système d'information, numérique
% Auteur   : MINKA MI NGUIDJOÏ Thierry Emmanuel
% Version  : 1.0 -- novembre 2026
% Corrigés : corriges/Corr_C01.tex
% ============================================================================

\chapter{Stratégie, système d'information, numérique}
\label{chap:vocabulaire}
\chaptermark{Vocabulaire}

\epigraph{%
    « Informatiser, c'est faire plus vite.\\
    Passer au numérique, c'est faire autrement. »%
}{Principe du cours}

\niveauChapitre{Fondamental}{%
    Prologue~; gouvernance stratégique des SI (UE MSN-2S3-MSI01)
}

\begin{boxobjectifs}
\begin{itemize}
    \item Définir stratégie, système d'information et numérique.
    \item Distinguer quatre niveaux de stratégie.
    \item Situer le schéma directeur dans l'histoire des documents de planification.
    \item Énumérer les six pièces d'un schéma directeur.
    \item Dire ce qu'un schéma directeur n'est pas.
\end{itemize}
\end{boxobjectifs}

\bigskip

% ============================================================================
\section*{Avant de commencer : « notre stratégie, c'est la fibre »}
% ============================================================================

Douala, jeudi 14~janvier~2027. M.~ETEME, directeur général de LiZbiZ, reçoit M.~BEYALA. Son entreprise de
logistique veut doubler son chiffre d'affaires en trois ans. C'est le plan HORIZON~3. « Notre stratégie
numérique est prête. Nous installons la fibre dans les quatre sites et nous changeons tous les postes. »

M.~BEYALA demande ce que les clients verront de différent. M.~ETEME hésite. La fibre accélère les échanges entre
les entrepôts. Elle ne change rien pour le client qui attend un colis. Le concurrent qui vient de s'installer,
lui, envoie un message à chaque étape de la livraison.

\begin{boxalert}
\textbf{Violation typique.} Présenter un programme d'équipement comme une stratégie numérique.

\smallskip
\textbf{Règle~:} une stratégie numérique se formule du point de vue de celui que l'organisation sert. Si rien
ne change pour lui, il s'agit d'une stratégie informatique. Elle est nécessaire. Elle ne suffit pas.
\end{boxalert}

% ============================================================================
\section{Trois mots à séparer}
\label{sec:c01-mots}
% ============================================================================

\begin{boxdef}
\textbf{Stratégie.} Ensemble cohérent de choix par lesquels une organisation décide où elle veut aller, ce qu'elle
fera et ce qu'elle renonce à faire, compte tenu de ses moyens et de son environnement.
\end{boxdef}

Une \termecle{stratégie} se reconnaît à ses renoncements. Un texte qui promet tout à tous n'en est pas une. Dans
une entreprise, elle vise un avantage sur des concurrents. Dans une administration, elle vise l'accomplissement
d'une mission fixée par un texte, avec des moyens votés.

\begin{boxdef}
\textbf{Système d'information.} Ensemble organisé de personnes, de processus, de données et de moyens
techniques qui collecte, traite, conserve et diffuse l'information dont une organisation a besoin pour agir.
\end{boxdef}

Le \termecle{système d'information} dépasse l'informatique. Un registre papier tenu dans une antenne régionale
en fait partie. Un schéma directeur qui n'inventorie que des serveurs décrit un parc et non un système
d'information.

\begin{boxdef}
\textbf{Numérique.} Usage des technologies de l'information qui modifie le service rendu, la relation avec
l'usager ou le modèle d'activité lui-même et pas seulement l'outil de travail interne.
\end{boxdef}

La distinction entre informatique et \termecle{numérique} n'est pas de mode. Informatiser, c'est faire plus vite
ce que l'on faisait déjà. Passer au numérique, c'est faire autrement~: le bénéficiaire qui suit son dossier sur
son téléphone ne reçoit pas le même service que celui qui se déplace à un guichet équipé d'un ordinateur.

\begin{tableaucours}{L{3.0cm}L{4.4cm}L{4.6cm}}{%
    \tableauHeader{Notion} &
    \tableauHeader{Question à laquelle elle répond} &
    \tableauHeader{Exemple à l'ANPP}
}
Stratégie de l'organisation & Que voulons-nous devenir~? & Payer chaque bénéficiaire à date fixe, où qu'il vive \\
Stratégie numérique & Que change le numérique dans notre mission~? & Paiement sur compte mobile, suivi en ligne \\
Stratégie du SI & De quelles informations et applications avons-nous besoin~? & Un référentiel unique des bénéficiaires \\
Stratégie informatique & Avec quelles techniques et quels moyens~? & Hébergement, réseau des antennes, compétences \\
\end{tableaucours}
\vspace{-0.8em}
\begin{center}{\small\itshape Tableau~1.1. Quatre niveaux de stratégie}\end{center}

La séparation entre stratégie du système d'information et stratégie informatique vient d'Earl et a été reprise par
Peppard et Ward~\cite{wardpeppard2016}. La première dit le \textit{quoi}, tirée des besoins des métiers. La
seconde dit le \textit{comment}. Beaucoup de schémas directeurs ne traitent que la seconde.

% ============================================================================
\section{Du plan informatique à la stratégie numérique}
\label{sec:c01-histoire}
% ============================================================================

Le vocabulaire a changé avec la place de l'informatique dans l'organisation.

\begin{tableaucours}{L{2.2cm}L{3.2cm}L{6.6cm}}{%
    \tableauHeader{Période} &
    \tableauHeader{Document type} &
    \tableauHeader{Ce qu'il cherchait}
}
Années 1970 & Plan informatique & Programmer l'achat des machines et l'automatisation des tâches de masse \\
Années 1980 & Schéma directeur informatique & Relier les applications aux objectifs de l'organisation. En France, méthode RACINES \\
Années 1990 & Schéma directeur du système d'information & Aligner le système d'information sur la stratégie. Modèle de Henderson et Venkatraman \\
Années 2000 & Plan d'urbanisation, gouvernance & Maîtriser la cohérence d'ensemble et les droits de décision \\
Depuis 2010 & Stratégie numérique & Faire du numérique une composante de la stratégie elle-même \\
\end{tableaucours}
\vspace{-0.8em}
\begin{center}{\small\itshape Tableau~1.2. Cinq générations de documents}\end{center}

Ces générations ne se remplacent pas. Elles s'empilent. Une administration qui écrit sa stratégie numérique a
toujours besoin de programmer des achats de matériel. L'erreur est de croire qu'un nouveau mot dispense du
travail de l'ancien.

% ============================================================================
\section{Ce que contient un schéma directeur}
\label{sec:c01-contenu}
% ============================================================================

\begin{boxprocedure}
\textbf{Les six pièces d'un schéma directeur.}
\begin{enumerate}
    \item \textbf{Le rappel de la stratégie} de l'organisation et de ses contraintes.
    \item \textbf{Le diagnostic} de l'existant, avec ses forces et ses écarts.
    \item \textbf{La cible}~: ce que sera le système d'information à l'horizon retenu, avec ses principes.
    \item \textbf{Le portefeuille de projets}, hiérarchisé, avec les dépendances.
    \item \textbf{La trajectoire}~: le calendrier, le budget par année, les moyens humains.
    \item \textbf{La gouvernance}~: qui décide, qui suit, quels indicateurs, quelle révision.
\end{enumerate}
\end{boxprocedure}

Chaque pièce répond à l'une des trois questions du prologue. Les pièces 1 à 3 établissent l'alignement. La
pièce 5 établit le financement. La pièce 6 établit la gouvernance. La pièce 4 les relie.

\subsection{Ce qu'il n'est pas}

\begin{tableaucours}{L{4.2cm}L{7.8cm}}{%
    \tableauHeader{Confusion} &
    \tableauHeader{Différence}
}
Un inventaire & L'inventaire décrit ce qui existe. Le schéma directeur décide de ce qui existera \\
Un cahier des charges & Le cahier des charges décrit un projet. Le schéma directeur choisit les projets \\
Un plan d'urbanisation & Le plan d'urbanisation fixe des règles de construction. Le schéma directeur fixe un programme de travaux \\
Un budget & Le budget autorise une dépense annuelle. Le schéma directeur donne la raison de cette dépense sur plusieurs années \\
Une politique de sécurité & La politique fixe des règles. Le schéma directeur programme les moyens de les appliquer \\
\end{tableaucours}
\vspace{-0.8em}
\begin{center}{\small\itshape Tableau~1.3. Cinq confusions fréquentes}\end{center}

% ============================================================================
\section{L'horizon et la révision}
\label{sec:c01-horizon}
% ============================================================================

Un schéma directeur couvre le plus souvent trois à cinq ans. En dessous de trois ans, il ne laisse pas le temps de
conduire un projet structurant. Au-delà de cinq ans, les hypothèses techniques et budgétaires ne tiennent plus.

L'horizon ne dit rien de la fréquence de révision. Un document à cinq ans se révise chaque année. La trajectoire
est précise pour la première année, indicative pour les suivantes. On parle de \termecle{plan glissant}. À
l'ANPP, l'horizon retenu est 2027-2031, calé sur la fin de la Stratégie nationale de développement en 2030 et sur
le cadre budgétaire à moyen terme de l'Agence.

% ============================================================================
\section{Trois organisations, trois régimes}
\label{sec:c01-cas}
% ============================================================================

\begin{boxscenario}{Les trois missions d'AUDIX CONSEIL en 2027}
\begin{tableaucours}{L{2.6cm}L{3.4cm}L{5.8cm}}{%
    \tableauHeader{Organisation} &
    \tableauHeader{Nature} &
    \tableauHeader{Demande}
}
ANPP & Établissement public, Yaoundé, 412 agents & Schéma directeur 2027-2031, après l'échec du précédent \\
LiZbiZ & Société anonyme de logistique, Douala & Stratégie numérique à l'appui du plan de croissance HORIZON~3 \\
BANQUE SANAGA & Établissement de crédit, Douala & Schéma directeur exigé après la panne et l'Opération WOURI \\
\end{tableaucours}

\medskip
\textbf{Analyse} (confiance élevée). Les trois demandes portent le même nom et obéissent à trois régimes.
L'Agence achète selon le code des marchés publics et finance selon le régime financier de l'État. L'entreprise
décide seule sous le droit des sociétés. La banque rend compte à un superviseur communautaire. Le
Chapitre~\ref{chap:droit} en tire une carte d'applicabilité.
\end{boxscenario}

\begin{boxmetier}{Directeur des systèmes d'information}
Refuser d'écrire seul le schéma directeur. Les pièces 1 et 3 appartiennent aux métiers et à la direction générale.
\end{boxmetier}

\begin{boxmetier}{Dirigeant}
Dire la stratégie de l'organisation avant de demander celle du système d'information. Sans la première, la seconde est un inventaire.
\end{boxmetier}

\begin{boxmetier}{Consultant}
Demander à chaque interlocuteur ce qu'il appelle schéma directeur. Trois personnes donnent souvent trois définitions.
\end{boxmetier}

\begin{boxresume}
\begin{itemize}
    \item Une stratégie se reconnaît à ses renoncements.
    \item Le système d'information comprend des personnes, des processus, des données et des techniques.
    \item Quatre niveaux~: stratégie de l'organisation, stratégie numérique, stratégie du SI, stratégie informatique.
    \item Un schéma directeur réunit six pièces~: stratégie rappelée, diagnostic, cible, portefeuille, trajectoire, gouvernance.
    \item L'horizon est de trois à cinq ans, la révision est annuelle.
\end{itemize}
\end{boxresume}

\begin{boxinfo}
\textbf{Points de contrôle.} Vous avez compris ce chapitre si vous savez~:
\begin{itemize}
    \item distinguer stratégie numérique et stratégie informatique sur un exemple~;
    \item citer les six pièces d'un schéma directeur~;
    \item expliquer ce qu'est un plan glissant.
\end{itemize}
\end{boxinfo}

\begin{boxexo}[title={\ding{45}\quad Exercice~1.1. Classer dix phrases \niveaubase}]
Classez dans l'un des quatre niveaux du Tableau~1.1~: (a) migrer la messagerie~; (b) ouvrir un guichet en ligne~; (c) servir les zones rurales~; (d) créer un référentiel des agents~; (e) recruter deux administrateurs de bases de données~; (f) réduire de moitié le délai de paiement~; (g) permettre le suivi d'un dossier par SMS~; (h) remplacer les onduleurs~; (i) partager les données avec l'état civil~; (j) devenir le premier logisticien du pays.
\end{boxexo}

\begin{boxexo}[title={\ding{45}\quad Exercice~1.2. La stratégie de LiZbiZ \niveaumoyen}]
Réécrivez en dix lignes la stratégie numérique de LiZbiZ du point de vue de ses clients. Dites ensuite ce que devient le projet de fibre dans ce nouveau texte.
\end{boxexo}

\begin{boxexo}[title={\ding{45}\quad Exercice~1.3. Six pièces \niveaumoyen}]
Reprenez le Tableau~P.1 du prologue. Pour chacune des six pièces d'un schéma directeur, dites si le document de 2021 la contenait, en partie ou pas du tout.
\end{boxexo}

\begin{boxexo}[title={\ding{45}\quad Exercice~1.4. Horizon \niveauexpert}]
La BANQUE SANAGA veut un schéma directeur à dix ans « pour ne plus y revenir ». Rédigez la réponse de M.~BEYALA en séparant ce qui peut être fixé à dix ans de ce qui ne le peut pas.
\end{boxexo}

\bigskip

\begin{boxinfo}
\textbf{Pour aller plus loin.}
\begin{itemize}
    \item Peppard et Ward~\cite{wardpeppard2016}, chapitres 1 à 3.
    \item Bharadwaj et al.~\cite{bharadwaj2013}, sur la fusion des stratégies.
\end{itemize}
\end{boxinfo}

\begin{boxinfo}
\textbf{Transition.} Les mots sont posés. Reste à savoir comment une organisation choisit sa stratégie. Le Chapitre~\ref{chap:strategie} présente les théories que le consultant doit connaître dans leur texte d'origine.
\end{boxinfo}

        % ============================================================================
% Fichier  : partie1/P01_C02_theories_strategie.tex
% Titre    : Les Théories de la Stratégie
% Auteur   : MINKA MI NGUIDJOÏ Thierry Emmanuel
% Version  : 1.0 -- novembre 2026
% Corrigés : corriges/Corr_C02.tex
% ============================================================================

\chapter{Les Théories de la Stratégie}
\label{chap:strategie}
\chaptermark{Théories de la stratégie}

\epigraph{%
    « L'essence de la stratégie est de choisir\\
    ce qu'il ne faut pas faire. »%
}{D'après Michael Porter}

\niveauChapitre{Fondamental}{%
    Chapitre~\ref{chap:vocabulaire}
}

\begin{boxobjectifs}
\begin{itemize}
    \item Appliquer les cinq forces et la chaîne de valeur à une organisation.
    \item Utiliser la grille VRIO pour qualifier une ressource informatique.
    \item Distinguer stratégie délibérée et stratégie émergente.
    \item Reconnaître une rupture et un marché biface.
    \item Raisonner un choix de faire ou de faire faire.
    \item Transposer ces théories à une administration.
\end{itemize}
\end{boxobjectifs}

\bigskip

% ============================================================================
\section*{Avant de commencer : « le concurrent livre à moto »}
% ============================================================================

Douala, mardi 2~février~2027. Le directeur commercial de LiZbiZ ne s'inquiète pas. La jeune pousse qui livre à
moto dans la ville « ne fait pas le même métier ». Elle ne transporte que des petits colis, elle n'a pas
d'entrepôt et ses prix ne couvrent sans doute pas ses coûts.

Mme~EKANI demande les chiffres. En dix-huit mois, ce concurrent a pris 9~\% des livraisons urbaines. Ses clients
sont de petits commerçants en ligne que LiZbiZ ne démarche pas. Il leur donne le suivi du colis en temps réel.
Deux grands comptes de LiZbiZ viennent de demander le même service.

\begin{boxalert}
\textbf{Violation typique.} Juger un nouvel entrant sur la qualité de son offre du moment.

\smallskip
\textbf{Règle~:} une rupture commence par le bas du marché. La question utile n'est pas « fait-il aussi bien
que nous~? » mais « à qui suffit ce qu'il fait et à quelle vitesse progresse-t-il~? »
\end{boxalert}

% ============================================================================
\section{Lire l'Environnement : Porter}
\label{sec:c02-porter}
% ============================================================================

Michael Porter publie en 1980 une méthode d'analyse d'un secteur~\cite{porter1980}. La rentabilité d'une
entreprise dépend de cinq forces~: la rivalité entre concurrents, la menace des nouveaux entrants, la menace
des produits de substitution, le pouvoir des fournisseurs et celui des clients.

\begin{tableaucours}{L{3.4cm}L{4.0cm}L{4.6cm}}{%
    \tableauHeader{Force} &
    \tableauHeader{Question} &
    \tableauHeader{Chez LiZbiZ}
}
Rivalité & Combien de concurrents, de quelle taille~? & Six transporteurs nationaux, prix proches \\
Nouveaux entrants & Est-il facile d'entrer~? & Une jeune pousse livre à moto avec une application \\
Substituts & Le besoin peut-il être servi autrement~? & Points de retrait tenus par des commerçants \\
Fournisseurs & Peuvent-ils imposer leurs prix~? & Carburant et pièces, prix subis \\
Clients & Peuvent-ils partir ou négocier~? & Un client pèse 38~\% du chiffre d'affaires \\
\end{tableaucours}
\vspace{-0.8em}
\begin{center}{\small\itshape Tableau~2.1. Les cinq forces appliquées à LiZbiZ}\end{center}

En 1985, Porter ajoute la \termecle{chaîne de valeur}~\cite{porter1985}. L'entreprise est découpée en activités
principales (logistique d'approvisionnement, production, logistique de distribution, vente, service) et en
activités de soutien (infrastructure, ressources humaines, technologie, achats). L'avantage naît de la manière
dont ces activités sont exécutées et reliées.

La même année, Porter et Millar montrent que l'information traverse toute la chaîne~\cite{portermillar1985}.
Chaque activité produit et consomme de l'information. Les technologies peuvent donc abaisser un coût,
différencier une offre ou changer les règles d'un secteur. Ce texte est le premier à faire du système
d'information une arme stratégique et non un centre de coût.

\begin{boxinfo}
\textbf{Trois stratégies génériques.} Porter distingue la domination par les coûts, la différenciation et la
concentration sur un segment. Une entreprise qui poursuit les trois à la fois reste « enlisée au milieu ». Le
système d'information ne sert pas les trois de la même manière~: l'automatisation sert les coûts, la
connaissance du client sert la différenciation.
\end{boxinfo}

\textbf{Limite.} Les cinq forces supposent un secteur aux frontières nettes. Les plateformes numériques
brouillent ces frontières~: un opérateur de téléphonie devient un concurrent des banques.

% ============================================================================
\section{Regarder au-dedans : Ressources et Capacités}
\label{sec:c02-ressources}
% ============================================================================

La théorie des ressources renverse le regard. L'avantage ne vient pas d'abord de la position sur un marché. Il
vient de ce que l'entreprise possède et que les autres ne peuvent pas copier. Wernerfelt l'énonce en
1984~\cite{wernerfelt1984}. Barney en donne en 1991 les conditions~\cite{barney1991}~: une ressource fonde un
avantage durable si elle a de la valeur, si elle est rare, si elle est difficile à imiter et si rien ne la
remplace. La grille VRIO, qui en découle, ajoute que l'organisation doit savoir l'exploiter.

\begin{tableaucours}{L{4.2cm}ccccL{3.0cm}}{%
    \tableauHeader{Ressource de LiZbiZ} &
    \tableauHeader{Valeur} &
    \tableauHeader{Rare} &
    \tableauHeader{Inimitable} &
    \tableauHeader{Exploitée} &
    \tableauHeader{Conclusion}
}
Postes de travail neufs & Oui & Non & Non & Oui & Parité \\
Logiciel de gestion d'entrepôt du marché & Oui & Non & Non & En partie & Parité \\
Dix ans de données de tournées & Oui & Oui & Oui & Non & Avantage inexploité \\
Réseau de 40 correspondants en zone rurale & Oui & Oui & Oui & Oui & Avantage durable \\
\end{tableaucours}
\vspace{-0.8em}
\begin{center}{\small\itshape Tableau~2.2. Grille VRIO de quatre ressources}\end{center}

La grille dit quelque chose d'utile au schéma directeur. Le matériel que tout le monde peut acheter ne donne
aucun avantage. Les données que seule l'entreprise détient en donnent un, à condition de les exploiter.

En 1997, Teece, Pisano et Shuen ajoutent les \termecle{capacités dynamiques}~\cite{teece1997}~: l'aptitude à
recomposer ses ressources quand l'environnement change. Pour un schéma directeur, cela justifie des dépenses qui
ne rapportent rien dans l'immédiat, comme une architecture qui permet de changer de fournisseur ou une équipe
capable d'intégrer vite un nouveau service.

% ============================================================================
\section{La Stratégie se Fait aussi en Marchant}
\label{sec:c02-mintzberg}
% ============================================================================

Mintzberg et Waters observent en 1985 que la stratégie réalisée diffère presque toujours de la stratégie
voulue~\cite{mintzberg1985}. Une partie de la stratégie \textit{délibérée} n'est jamais exécutée. Une stratégie
\textit{émergente} apparaît, faite de décisions prises au fil des événements.

Le bilan de l'ANPP en est l'illustration. Six projets délibérés ont abouti. Onze projets émergents se sont
ajoutés. Mintzberg n'y voit pas un échec en soi. L'échec est de ne pas s'en apercevoir. Un schéma directeur doit
donc prévoir comment un projet non prévu y entre et comment un projet prévu en sort.

% ============================================================================
\section{Rupture, Plateformes et Coûts de Transaction}
\label{sec:c02-numerique}
% ============================================================================

Christensen décrit en 1997 l'\termecle{innovation de rupture}~\cite{christensen1997}. Un nouvel entrant propose
un produit moins bon et moins cher à des clients que les entreprises en place négligent. Il s'améliore, puis
remonte vers leur clientèle. Les entreprises établies ne le voient pas venir parce qu'elles écoutent leurs
meilleurs clients.

Rochet et Tirole formalisent en 2003 les \termecle{marchés bifaces}~\cite{rochettirole2003}. Une plateforme
sert deux groupes qui ont besoin l'un de l'autre. La valeur pour chaque groupe croît avec la taille de l'autre.
Le paiement par téléphone mobile en est l'exemple le plus proche~: plus il y a de commerçants qui l'acceptent,
plus les porteurs l'utilisent. Parker, Van Alstyne et Choudary en ont tiré une méthode~\cite{parker2016}.

La théorie des coûts de transaction éclaire le choix de faire ou de faire faire. Coase demande en 1937 pourquoi
les entreprises existent~\cite{coase1937}. Williamson précise en 1975 qu'il vaut mieux faire soi-même quand
l'actif est très spécifique, quand l'incertitude est forte et quand l'échange est fréquent~\cite{williamson1975}.

\begin{tableaucours}{L{4.4cm}L{3.6cm}L{4.0cm}}{%
    \tableauHeader{Activité informatique} &
    \tableauHeader{Spécificité} &
    \tableauHeader{Tendance}
}
Messagerie & Faible & Faire faire \\
Hébergement standard & Faible & Faire faire \\
Application de calcul des prestations & Forte & Garder la maîtrise \\
Référentiel des bénéficiaires & Forte & Garder la maîtrise \\
\end{tableaucours}
\vspace{-0.8em}
\begin{center}{\small\itshape Tableau~2.3. Faire ou faire faire à l'ANPP}\end{center}

% ============================================================================
\section{Relier Objectifs et Mesures}
\label{sec:c02-bsc}
% ============================================================================

Kaplan et Norton proposent en 1992 le \termecle{tableau de bord prospectif}~\cite{kaplannorton1992}. La
performance se lit sur quatre axes~: finances, clients, processus internes, apprentissage. Chaque objectif reçoit
une mesure, une cible et une initiative. Le modèle sera repris au Chapitre~\ref{chap:piloter} pour piloter
l'exécution du schéma directeur.

% ============================================================================
\section{Et dans une Administration ?}
\label{sec:c02-public}
% ============================================================================

Ces théories sont nées dans l'entreprise. Leur transposition demande trois précautions.

\begin{enumerate}
    \item \textbf{La mission remplace le marché.} L'ANPP n'a pas de concurrent. Les cinq forces se relisent en
          parties prenantes~: tutelle, bailleurs, bénéficiaires, opérateurs de paiement, fournisseurs.
    \item \textbf{La valeur est publique.} Elle se mesure en délai, en accès, en confiance et non en marge.
    \item \textbf{Les ressources sont contraintes.} Une agence ne recrute ni n'achète librement. La grille VRIO
          reste utile pour repérer ce qu'elle détient d'unique, par exemple son fichier des bénéficiaires.
\end{enumerate}

\begin{boxscenario}{Mission LiZbiZ. La note stratégique du 4 février 2027}
\textbf{Constatations.} Le plan HORIZON~3 vise le doublement du chiffre d'affaires. Un client pèse 38~\% des
ventes. Une jeune pousse a pris 9~\% des livraisons urbaines à Douala en dix-huit mois. LiZbiZ détient dix ans
de données de tournées qu'elle n'exploite pas.

\textbf{Analyse} (confiance moyenne). La menace vient d'un nouvel entrant de rupture sur le segment urbain.
L'avantage de LiZbiZ est ailleurs~: son réseau rural et ses données. La stratégie numérique doit défendre le
premier par le suivi des colis et exploiter les secondes pour optimiser les tournées. La fibre n'est qu'un
moyen.
\end{boxscenario}

\begin{boxmetier}{Consultant}
Citer une théorie avec son auteur et sa date. Dire ce qu'elle explique et ce qu'elle n'explique pas.
\end{boxmetier}

\begin{boxmetier}{Directeur des systèmes d'information}
Passer chaque grand poste de dépense à la grille VRIO. Ce qui n'est ni rare ni inimitable se standardise ou s'achète.
\end{boxmetier}

\begin{boxmetier}{Dirigeant}
Demander une fois par an quels projets sont nés hors du plan. Leur nombre dit ce que le plan n'avait pas vu.
\end{boxmetier}

\begin{boxresume}
\begin{itemize}
    \item Porter~: cinq forces, chaîne de valeur, trois stratégies génériques. L'information traverse toute la chaîne.
    \item Barney~: une ressource fonde un avantage si elle a de la valeur, si elle est rare, inimitable et exploitée.
    \item Teece~: les capacités dynamiques justifient d'investir dans l'aptitude à changer.
    \item Mintzberg~: la stratégie réalisée mêle le délibéré et l'émergent.
    \item Christensen, Rochet et Tirole, Williamson~: rupture, plateformes, faire ou faire faire.
    \item Dans une administration, la mission remplace le marché et la valeur est publique.
\end{itemize}
\end{boxresume}

\begin{boxinfo}
\textbf{Points de contrôle.} Vous avez compris ce chapitre si vous savez~:
\begin{itemize}
    \item nommer les cinq forces~;
    \item énoncer les quatre conditions de Barney~;
    \item donner un exemple local de marché biface~;
    \item dire quand il vaut mieux faire soi-même.
\end{itemize}
\end{boxinfo}

\begin{boxexo}[title={\ding{45}\quad Exercice~2.1. Cinq forces \niveaubase}]
Appliquez les cinq forces à une banque de détail camerounaise face au paiement par téléphone mobile. Une ligne par force.
\end{boxexo}

\begin{boxexo}[title={\ding{45}\quad Exercice~2.2. VRIO \niveaumoyen}]
Passez à la grille VRIO quatre ressources de l'ANPP~: le fichier des bénéficiaires, le centre de données, l'application e-RECETTES, le réseau des dix antennes.
\end{boxexo}

\begin{boxexo}[title={\ding{45}\quad Exercice~2.3. Délibéré, émergent \niveaumoyen}]
Classez les onze projets hors schéma directeur de l'ANPP en trois familles possibles~: réponse à une urgence, opportunité de financement, demande de la tutelle. Que doit prévoir le nouveau document pour chacune~?
\end{boxexo}

\begin{boxexo}[title={\ding{45}\quad Exercice~2.4. Faire ou faire faire \niveauexpert}]
La BANQUE SANAGA envisage de confier l'exploitation de son système central à un prestataire. Appliquez les trois critères de Williamson. Concluez, puis dites ce que le superviseur exigera.
\end{boxexo}

\bigskip

\begin{boxinfo}
\textbf{Pour aller plus loin.}
\begin{itemize}
    \item Porter~\cite{porter1980,porter1985} et Porter et Millar~\cite{portermillar1985}.
    \item Barney~\cite{barney1991} et Teece et al.~\cite{teece1997}.
    \item Mintzberg et Waters~\cite{mintzberg1985}.
    \item Christensen~\cite{christensen1997}.
\end{itemize}
\end{boxinfo}

\begin{boxinfo}
\textbf{Transition.} Ces théories disent comment une organisation choisit sa route. Elles ne disent pas comment le système d'information la suit. C'est la question de l'alignement et de la valeur, objet du Chapitre~\ref{chap:alignement}.
\end{boxinfo}

        % ============================================================================
% Fichier  : partie1/P01_C03_alignement_valeur.tex
% Titre    : L'alignement et la valeur
% Auteur   : MINKA MI NGUIDJOÏ Thierry Emmanuel
% Version  : 1.0 -- novembre 2026
% Corrigés : corriges/Corr_C03.tex
% ============================================================================

\chapter{L'alignement et la valeur}
\label{chap:alignement}
\chaptermark{Alignement et valeur}

\epigraph{%
    « Un bénéfice n'est pas livré par un projet.\\
    Il est obtenu par un métier qui change. »%
}{D'après Ward et Daniel}

\niveauChapitre{Fondamental}{%
    Chapitres~\ref{chap:vocabulaire} et~\ref{chap:strategie}
}

\begin{boxobjectifs}
\begin{itemize}
    \item Situer une organisation sur les stades de Nolan.
    \item Tirer des besoins d'information des facteurs clés de succès.
    \item Classer un portefeuille applicatif dans la grille de McFarlan.
    \item Décrire le modèle d'alignement stratégique et ses quatre perspectives.
    \item Coter la maturité de l'alignement.
    \item Discuter le paradoxe de la productivité et la thèse de Carr.
\end{itemize}
\end{boxobjectifs}

\bigskip

% ============================================================================
\section*{Avant de commencer : « l'informatique ne suit pas »}
% ============================================================================

Yaoundé, lundi 8~février~2027. Le directeur des prestations de l'ANPP résume son avis en une phrase~: « L'informatique
ne suit pas. » Il attend depuis deux ans un état quotidien des paiements. Mme~NGO~BIYONG répond qu'elle n'a
jamais reçu cette demande par écrit et qu'elle a appris par hasard la réforme des prestations de 2025.

Chacun a raison. La division informatique n'assiste pas aux réunions de direction. Les métiers ne savent pas quels
projets sont en cours. Personne n'arbitre. Mme~EKANI note que le problème n'est pas technique. Il porte un nom
depuis 1993.

\begin{boxalert}
\textbf{Violation typique.} Traiter un défaut d'alignement par un achat de matériel ou de logiciel.

\smallskip
\textbf{Règle~:} l'alignement est une affaire de décisions et de personnes. Il se construit par une instance
qui arbitre, par des responsables métiers désignés et par une mesure de la valeur obtenue.
\end{boxalert}

% ============================================================================
\section{Mesurer la maturité : Nolan}
\label{sec:c03-nolan}
% ============================================================================

Richard Nolan observe dans les années 1970 que la dépense informatique d'une organisation suit une courbe en S
répétée. Il en tire un modèle en six stades~\cite{nolan1979}~: initiation, contagion, contrôle, intégration,
administration des données, maturité. À chaque stade correspondent un type d'applications, une organisation et
un mode de contrôle des dépenses.

Le modèle a vieilli dans ses détails. Son idée reste juste~: une organisation ne saute pas de stade. Une agence
dont les applications ont proliféré sans contrôle ne peut pas passer d'un coup à la gestion des données. Il lui
faut d'abord reprendre la main sur ses dépenses et son inventaire.

% ============================================================================
\section{Partir des besoins des dirigeants : Rockart}
\label{sec:c03-rockart}
% ============================================================================

John Rockart propose en 1979 de partir des \termecle{facteurs clés de succès}~\cite{rockart1979}. Ce sont les
quelques domaines où les résultats doivent être bons pour que l'organisation réussisse. On les obtient par
entretien avec les dirigeants. De chaque facteur on tire les informations nécessaires, puis les applications.

\begin{tableaucours}{L{4.2cm}L{3.8cm}L{4.0cm}}{%
    \tableauHeader{Facteur clé de succès de l'ANPP} &
    \tableauHeader{Information nécessaire} &
    \tableauHeader{Conséquence pour le SI}
}
Payer à date fixe & État des paiements par antenne, chaque jour & Tableau de bord des paiements \\
Ne payer que les ayants droit & Fichier unique, à jour des décès & Lien avec l'état civil \\
Rendre compte à la tutelle & Dépense par programme & Comptabilité analytique \\
Rester joignable & Demandes et délais de réponse & Guichet unique des demandes \\
\end{tableaucours}
\vspace{-0.8em}
\begin{center}{\small\itshape Tableau~3.1. Des facteurs clés de succès aux applications}\end{center}

% ============================================================================
\section{Classer les applications : McFarlan}
\label{sec:c03-mcfarlan}
% ============================================================================

Warren McFarlan publie en 1984 une grille qui classe les applications selon leur importance pour l'activité
d'aujourd'hui et pour celle de demain~\cite{mcfarlan1984}. Peppard et Ward en ont donné la version la plus
employée~\cite{wardpeppard2016}.

\begin{tableaucours}{L{2.6cm}L{4.4cm}L{5.0cm}}{%
    \tableauHeader{Famille} &
    \tableauHeader{Rôle} &
    \tableauHeader{Conduite à tenir}
}
Stratégique & Décisive pour la stratégie future & Investir, innover, garder la maîtrise \\
Opérationnelle clé & L'activité s'arrête sans elle & Fiabiliser, sécuriser, ne pas rater \\
À fort potentiel & Pourrait devenir stratégique & Expérimenter à petit budget, trancher vite \\
De soutien & Utile, sans enjeu propre & Standardiser, acheter, réduire le coût \\
\end{tableaucours}
\vspace{-0.8em}
\begin{center}{\small\itshape Tableau~3.2. La grille des applications}\end{center}

La grille interdit de gérer toutes les applications de la même manière. Une application de soutien développée sur
mesure est un gaspillage. Une application opérationnelle clé sans secours est un risque. À l'ANPP, e-RECETTES
est opérationnelle clé et n'avait pas de site de secours.

% ============================================================================
\section{Le modèle d'alignement stratégique}
\label{sec:c03-sam}
% ============================================================================

Henderson et Venkatraman publient en 1993 le modèle qui a donné son nom à l'alignement~\cite{henderson1993}. Il
croise deux axes. Le premier oppose l'affaire et les technologies. Le second oppose l'externe (la stratégie) et
l'interne (l'organisation et les processus).

\begin{tableaucours}{L{3.0cm}L{4.4cm}L{4.6cm}}{%
    \tableauHeader{} &
    \tableauHeader{Affaire} &
    \tableauHeader{Technologies de l'information}
}
Externe & Stratégie d'affaire~: offre, compétences distinctives, alliances & Stratégie des TI~: périmètre technologique, compétences, partenariats \\
Interne & Organisation~: structure, processus, compétences & Infrastructure et processus des TI~: architecture, processus, compétences \\
\end{tableaucours}
\vspace{-0.8em}
\begin{center}{\small\itshape Tableau~3.3. Les quatre domaines du modèle d'alignement}\end{center}

L'alignement a deux dimensions. L'\textit{adéquation stratégique} relie l'externe et l'interne d'un même côté.
L'\textit{intégration fonctionnelle} relie l'affaire et les technologies. Les auteurs décrivent quatre parcours.

\begin{tableaucours}{L{3.6cm}L{4.2cm}L{4.2cm}}{%
    \tableauHeader{Perspective} &
    \tableauHeader{Point de départ} &
    \tableauHeader{Rôle du DSI}
}
Exécution de la stratégie & La stratégie d'affaire commande l'organisation, qui commande le SI & Exécutant \\
Transformation technologique & La stratégie d'affaire commande la stratégie des TI & Architecte \\
Potentiel concurrentiel & Une technologie ouvre une stratégie d'affaire nouvelle & Catalyseur \\
Niveau de service & La stratégie des TI commande le service rendu aux métiers & Prestataire \\
\end{tableaucours}
\vspace{-0.8em}
\begin{center}{\small\itshape Tableau~3.4. Quatre perspectives d'alignement}\end{center}

La plupart des schémas directeurs ne connaissent que la première perspective. La troisième est celle du
numérique~: c'est la technologie qui suggère la stratégie. Le paiement sur compte mobile n'a été pensable pour
l'ANPP que le jour où les opérateurs l'ont offert.

\subsection{Mesurer l'alignement : Luftman}

Jerry Luftman propose en 2000 une mesure de la maturité de l'alignement sur six critères et cinq
niveaux~\cite{luftman2000}~: communication, mesure de la valeur, gouvernance, partenariat, périmètre et
architecture, compétences. Le cours l'emploie sous forme d'exigences cotées, pour rester dans sa convention.

\begin{boxlecture}{L'alignement à l'ANPP, 9 février 2027}
\begin{tableaucours}{L{7.6cm}cL{3.4cm}}{%
    \tableauHeader{Exigence} &
    \tableauHeader{Cote} &
    \tableauHeader{Preuve}
}
Communication~: la division informatique participe aux réunions de direction & \NON{} & Feuilles de présence 2026 \\
Communication~: les métiers connaissent les projets en cours & \NON{} & Entretiens, 14 sur 18 \\
Valeur~: chaque projet a un bénéfice attendu écrit & \NON{} & Fiches de projet \\
Valeur~: un bilan est fait après la mise en service & \NON{} & Aucun bilan \\
Gouvernance~: une instance arbitre les projets & \NON{} & Aucune instance \\
Gouvernance~: le budget suit les priorités arbitrées & \NEK{} & Clé de répartition non documentée \\
Partenariat~: chaque projet a un responsable métier & \OUI{} & Notes de désignation, 9 sur 11 \\
Architecture~: une cartographie applicative est tenue & \NON{} & Aucune \\
Compétences~: un plan de formation existe & \OUI{} & Plan 2026 \\
Compétences~: les postes clés ont un suppléant & \NEK{} & Organigramme incomplet \\
\end{tableaucours}

\medskip
\textbf{Décompte.} Deux \OUI{}, six \NON{}, deux \NEK{} sur dix.
\hfill\lecture{0{,}20}{0{,}20}{1{,}37}

\medskip
\textbf{Lecture.} L'alignement n'est pas établi~: $\mu = 0{,}20$. Six écarts touchent la communication, la valeur, la
gouvernance et l'architecture. Les deux \NEK{} sont de preuve et se lèveront pendant le diagnostic. Le point
d'appui est le partenariat~: les métiers désignent déjà un responsable par projet.
\end{boxlecture}

% ============================================================================
\section{Graduer l'ambition : Venkatraman}
\label{sec:c03-venkatraman}
% ============================================================================

Venkatraman décrit en 1994 cinq niveaux de transformation par les technologies~\cite{venkatraman1994}. Les deux
premiers sont évolutifs~: l'exploitation localisée, puis l'intégration interne. Les trois suivants sont
révolutionnaires~: la refonte des processus, la refonte du réseau d'affaires, la redéfinition du périmètre
d'activité. Les bénéfices croissent avec le niveau, les changements d'organisation aussi.

Un schéma directeur doit dire à quel niveau il se place. Informatiser le guichet d'une antenne relève du niveau
un. Faire payer les prestations par un opérateur de téléphonie relève du niveau quatre et suppose des accords,
des textes et une autre organisation.

% ============================================================================
\section{La valeur en question}
\label{sec:c03-valeur}
% ============================================================================

Robert Solow écrit en 1987 que l'on voit des ordinateurs partout sauf dans les statistiques de
productivité~\cite{solow1987}. Brynjolfsson examine ce \termecle{paradoxe de la productivité} en
1993~\cite{brynjolfsson1993} et avance quatre explications~: des erreurs de mesure, des délais d'apprentissage,
une redistribution des gains entre entreprises et une mauvaise gestion des investissements. Ses travaux
ultérieurs montrent que la valeur apparaît quand l'investissement technique s'accompagne d'un changement
d'organisation.

Nicholas Carr soutient en 2003 la thèse adverse~\cite{carr2003}. Les technologies de l'information deviennent
une commodité, comme l'électricité. Ce que tous peuvent acheter ne distingue personne. Il recommande de dépenser
moins, de suivre plutôt que de mener, de gérer les risques plutôt que les opportunités.

\begin{boxinfo}
\textbf{Que retenir de la controverse.} Carr a raison pour l'infrastructure. La théorie des ressources
(Chapitre~\ref{chap:strategie}) dit la même chose. Il a tort pour l'usage~: deux organisations dotées des mêmes
outils n'en tirent pas la même valeur. La valeur vient des données, des processus et des compétences, non du
matériel.
\end{boxinfo}

Ward et Daniel proposent en 2006 la \termecle{gestion des bénéfices}~\cite{warddaniel2006}. Un bénéfice n'est pas
livré par un projet. Il est obtenu par un métier qui change sa manière de travailler. Chaque bénéfice a donc un
propriétaire, une mesure et une date, distincts de ceux du projet.

% ============================================================================
\section{Gouvernance, modèle opérationnel et fusion des stratégies}
\label{sec:c03-gouvernance}
% ============================================================================

Weill et Ross montrent en 2004 que les organisations les plus performantes ont clarifié qui décide de quoi en
matière de technologies~\cite{weillross2004}. Ils distinguent cinq décisions~: les principes, l'architecture,
l'infrastructure, les besoins applicatifs, l'investissement. Un schéma directeur dit pour chacune qui décide et
qui est consulté.

Ross, Weill et Robertson ajoutent en 2006 le \termecle{modèle opérationnel}~\cite{ross2006}. Une organisation
choisit son degré de standardisation des processus et son degré d'intégration des données. Quatre modèles en
résultent~: diversification, coordination, réplication, unification. L'ANPP, qui applique la même procédure
dans dix antennes sur un même fichier, relève de l'unification. LiZbiZ, dont les sites partagent les clients
mais pas toujours les procédures, relève de la coordination.

Bharadwaj et ses coauteurs proposent en 2013 de ne plus parler d'alignement~\cite{bharadwaj2013}. La stratégie
numérique n'est pas une stratégie fonctionnelle alignée sur la stratégie d'affaire. Elle en est une composante.
Westerman, Bonnet et McAfee croisent en 2014 deux capacités~\cite{westerman2014}~: la capacité numérique (ce que
l'on sait faire avec la technologie) et la capacité de conduite (vision, gouvernance, engagement). Seules les
organisations fortes sur les deux transforment durablement leur activité.

\begin{boxmetier}{Directeur des systèmes d'information}
Tenir la grille des applications à jour et la présenter chaque année. Demander pour chaque projet le nom du métier qui en obtiendra le bénéfice.
\end{boxmetier}

\begin{boxmetier}{Responsable métier}
Écrire le bénéfice attendu avant le projet et le mesurer après. Sans cela, l'investissement ne se justifie ni ne se défend.
\end{boxmetier}

\begin{boxmetier}{Dirigeant}
Faire entrer le responsable du système d'information dans les réunions où la stratégie se discute.
\end{boxmetier}

\begin{boxresume}
\begin{itemize}
    \item Nolan~: une organisation ne saute pas de stade de maturité.
    \item Rockart~: les facteurs clés de succès donnent les besoins d'information.
    \item McFarlan~: quatre familles d'applications, quatre conduites.
    \item Henderson et Venkatraman~: quatre domaines, deux dimensions, quatre perspectives.
    \item Luftman~: l'alignement se mesure sur six critères.
    \item La valeur vient de l'usage et du changement d'organisation, non de l'outil.
    \item Weill et Ross~: dire qui décide de quoi. Bharadwaj~: la stratégie numérique est une composante de la stratégie.
\end{itemize}
\end{boxresume}

\begin{boxinfo}
\textbf{Points de contrôle.} Vous avez compris ce chapitre si vous savez~:
\begin{itemize}
    \item dessiner les quatre domaines du modèle d'alignement~;
    \item classer une application dans la grille de McFarlan~;
    \item citer deux explications du paradoxe de la productivité~;
    \item dire ce qu'est un propriétaire de bénéfice.
\end{itemize}
\end{boxinfo}

\begin{boxexo}[title={\ding{45}\quad Exercice~3.1. Grille de McFarlan \niveaubase}]
Classez~: la paie de l'ANPP, e-RECETTES, un pilote de paiement sur compte mobile dans une antenne, la messagerie, le fichier des bénéficiaires, un site d'information du public.
\end{boxexo}

\begin{boxexo}[title={\ding{45}\quad Exercice~3.2. Quelle perspective \niveaumoyen}]
Pour chacune des trois organisations du cours, dites quelle perspective du Tableau~3.4 décrit le mieux la situation actuelle et laquelle vous recommandez. Justifiez en trois lignes.
\end{boxexo}

\begin{boxexo}[title={\ding{45}\quad Exercice~3.3. Lever deux NEK \niveaumoyen}]
Quelles pièces demandez-vous pour lever les deux \NEK{} de la lecture de l'ANPP~? Recalculez $\mu$ et $U$ dans le cas le plus favorable et dans le moins favorable.
\end{boxexo}

\begin{boxexo}[title={\ding{45}\quad Exercice~3.4. Carr à la banque \niveauexpert}]
Le directeur financier de la BANQUE SANAGA cite Carr pour réduire de 30~\% le budget informatique. Rédigez une réponse d'une page qui lui donne raison sur un point et tort sur un autre, avec des exemples tirés de la banque.
\end{boxexo}

\bigskip

\begin{boxinfo}
\textbf{Pour aller plus loin.}
\begin{itemize}
    \item Henderson et Venkatraman~\cite{henderson1993}.
    \item Luftman~\cite{luftman2000}.
    \item Carr~\cite{carr2003} et Brynjolfsson~\cite{brynjolfsson1993}.
    \item Weill et Ross~\cite{weillross2004}~; Ross et al.~\cite{ross2006}.
\end{itemize}
\end{boxinfo}

\begin{boxinfo}
\textbf{Transition.} Les théories donnent des idées. Les normes et les référentiels donnent des critères. Avant de regarder le Cameroun, le Chapitre~\ref{chap:normes} présente ceux que le monde emploie.
\end{boxinfo}


    \part{Le contexte international}
    \label{part:international}

        % ============================================================================
% Fichier  : partie2/P02_C04_normes_referentiels.tex
% Titre    : Normes et Référentiels
% Auteur   : MINKA MI NGUIDJOÏ Thierry Emmanuel
% Version  : 1.0 -- novembre 2026
% Corrigés : corriges/Corr_C04.tex
% ============================================================================

\chapter{Normes et Référentiels}
\label{chap:normes}
\chaptermark{Normes et référentiels}

\epigraph{%
    « Le conseil ne rédige pas le schéma directeur.\\
    Il le demande, l'approuve et le surveille. »%
}{D'après l'ISO/IEC 38500}

\niveauChapitre{Fondamental}{%
    Chapitre~\ref{chap:alignement}~; COBIT et ITIL (UE MSN-2S3-MSI03)~; architecture d'entreprise (UE MSN-2S3-MSI02)
}

\begin{boxobjectifs}
\begin{itemize}
    \item Distinguer texte obligatoire, norme certifiable et cadre volontaire.
    \item Décrire les trois tâches de l'ISO/IEC~38500.
    \item Citer les objectifs de COBIT~2019 qui portent le schéma directeur.
    \item Situer TOGAF, ITIL~4 et l'ISO~21504.
    \item Déclarer les référentiels d'un schéma directeur.
\end{itemize}
\end{boxobjectifs}

\bigskip

% ============================================================================
\section*{Avant de commencer : « conforme à COBIT »}
% ============================================================================

Douala, mercredi 20~janvier~2027. Le schéma directeur que la BANQUE SANAGA\footnote{La BANQUE SANAGA est fictive.
Sa panne de 2026 et l'Opération WOURI qui a suivi figurent dans d'autres cours de l'auteur.} avait remis à son
superviseur en 2024 portait en couverture la mention « conforme à COBIT, ITIL et TOGAF ». Mme~NKENG, responsable
du contrôle de second niveau, a cherché sur quoi reposait cette phrase. Elle n'a trouvé ni liste des objectifs
retenus, ni édition des référentiels, ni évaluation.

Le superviseur a posé la même question après la panne. La banque n'a rien pu produire. Mme~NKENG résume~: « Nous
avons écrit trois noms sur une couverture. »

\begin{boxalert}
\textbf{Violation typique.} Écrire « conforme à » devant le nom d'un cadre volontaire.

\smallskip
\textbf{Règle~:} on est conforme à un texte ou certifié selon une norme, sur un périmètre et à une date. Un
cadre volontaire s'emploie et se déclare~: lesquels de ses éléments, dans quelle édition, pour quel usage.
\end{boxalert}

% ============================================================================
\section{Trois Natures de Référentiel}
\label{sec:c04-natures}
% ============================================================================

\begin{tableaucours}{L{2.8cm}L{3.0cm}L{2.8cm}L{3.4cm}}{%
    \tableauHeader{Nature} &
    \tableauHeader{Force} &
    \tableauHeader{Exemple} &
    \tableauHeader{Preuve attendue}
}
Texte obligatoire & S'impose, sous sanction & Code des marchés publics & Conformité à l'article \\
Norme certifiable & Volontaire, vérifiée par un tiers & ISO/IEC~27001 & Certificat, périmètre, date \\
Cadre volontaire & Volontaire, déclaré & COBIT 2019, TOGAF & Usage déclaré et documenté \\
Norme de recommandation & Volontaire, non certifiable & ISO/IEC~38500, ISO~31000 & Usage déclaré \\
\end{tableaucours}
\vspace{-0.8em}
\begin{center}{\small\itshape Tableau~4.1. Quatre natures de référentiel}\end{center}

Une norme ou un cadre ne devient contraignant que de trois manières~: un texte y renvoie, un contrat l'exige ou
l'organisation s'y engage par une décision de son organe délibérant. Hors de ces cas, c'est une référence que
l'on emprunte et que l'on déclare.

% ============================================================================
\section{Gouverner : ISO/IEC 38500 et ISO 37000}
\label{sec:c04-38500}
% ============================================================================

L'ISO/IEC~38500 traite de la gouvernance des technologies de l'information~\cite{iso38500}. Sa troisième édition
date de février~2024 et s'aligne sur l'ISO~37000 de 2021, qui traite de la gouvernance des
organismes~\cite{iso37000}. Elle s'adresse à l'organe de gouvernance, c'est-à-dire au conseil et non à la
direction informatique.

Son modèle tient en trois tâches que le conseil exerce sur l'usage des technologies~: \textbf{évaluer} l'usage
actuel et futur, \textbf{diriger} la préparation des stratégies et des politiques, \textbf{surveiller} la
conformité et la performance. Elle énonce des principes de bonne gouvernance, parmi lesquels la responsabilité,
la stratégie, l'acquisition, la performance, la conformité et le comportement humain.

\begin{boxinfo}
\textbf{Portée pour le schéma directeur.} La norme place le schéma directeur du bon côté de la table. Le conseil
ne rédige pas le document. Il le demande, l'approuve et en surveille l'exécution. Le bilan de l'ANPP au prologue
montre un conseil qui a approuvé sans surveiller~: deux tâches sur trois manquaient.
\end{boxinfo}

% ============================================================================
\section{Gouverner et Gérer : COBIT 2019}
\label{sec:c04-cobit}
% ============================================================================

COBIT~2019 est le cadre de l'ISACA~\cite{cobit2019}. Il compte quarante objectifs, répartis en un domaine de
gouvernance et quatre domaines de gestion. Le cours de gouvernance (UE MSN-2S3-MSI03) l'a présenté. Six
objectifs concernent directement le schéma directeur.

\begin{tableaucours}{L{1.6cm}L{4.2cm}L{6.2cm}}{%
    \tableauHeader{Objectif} &
    \tableauHeader{Intitulé résumé} &
    \tableauHeader{Usage pour le schéma directeur}
}
EDM01 & Cadre de gouvernance & Qui décide du schéma directeur \\
EDM02 & Livraison des bénéfices & Suivi de la valeur des investissements \\
APO02 & Stratégie gérée & Diagnostic, cible, feuille de route \\
APO03 & Architecture d'entreprise gérée & Cible d'architecture et plan de migration \\
APO05 & Portefeuille géré & Sélection et arbitrage des projets \\
APO06 & Budget et coûts gérés & Chiffrage et suivi financier \\
\end{tableaucours}
\vspace{-0.8em}
\begin{center}{\small\itshape Tableau~4.2. Six objectifs de COBIT 2019 pour le schéma directeur}\end{center}

COBIT propose aussi une \termecle{cascade des objectifs}~: des besoins des parties prenantes aux objectifs de
l'entreprise, puis aux objectifs d'alignement, puis aux objectifs de gouvernance et de gestion. C'est un outil
commode pour démontrer l'alignement pièce par pièce.

% ============================================================================
\section{Dessiner la Cible : TOGAF, ISO 42010, ArchiMate}
\label{sec:c04-togaf}
% ============================================================================

TOGAF est la norme d'architecture d'entreprise de l'Open Group, dont la dixième édition date de
2022~\cite{togaf10}. Sa méthode de développement de l'architecture parcourt des phases~: vision, architecture
métier, architecture des systèmes d'information (données et applications), architecture technique, opportunités
et solutions, planification de la migration, gouvernance de la mise en œuvre, gestion du changement. Le cours
d'architecture (UE MSN-2S3-MSI02) la détaille.

Pour le schéma directeur, trois phases comptent. La vision donne les principes. Les phases d'architecture donnent
l'existant et la cible. La planification de la migration donne les \textit{architectures de transition}, qui
sont les étapes de la trajectoire.

L'ISO/IEC/IEEE~42010 fixe la manière de décrire une architecture~\cite{iso42010}~: des parties prenantes ont
des préoccupations, auxquelles répondent des points de vue. ArchiMate est le langage de modélisation de l'Open
Group, organisé en couches métier, application et technologie.

% ============================================================================
\section{Servir : ITIL 4 et ISO/IEC 20000}
\label{sec:c04-itil}
% ============================================================================

ITIL~4 décrit un \termecle{système de valeur des services}~\cite{itil4}. La demande y entre, la valeur en sort.
Le système réunit des principes directeurs, une gouvernance, une chaîne de valeur des services, des pratiques et
une amélioration continue. ITIL~4 comporte un module consacré à la stratégie numérique et informatique.

Pour le schéma directeur, ITIL apporte une idée que les documents oublient~: un projet livré devient un service à
exploiter. Le coût de fonctionnement, les compétences d'exploitation et les niveaux de service doivent figurer à
la trajectoire. L'ISO/IEC~20000-1 est la norme certifiable de gestion des services.

% ============================================================================
\section{Arbitrer : les Normes de Portefeuille}
\label{sec:c04-portefeuille}
% ============================================================================

L'ISO~21502 traite du management de projet~\cite{iso21502}. L'ISO~21504 traite du management de
portefeuille~\cite{iso21504}. La distinction est celle du cours~: le projet fait bien les choses, le portefeuille
choisit les bonnes choses à faire.

\begin{tableaucours}{L{2.6cm}L{4.4cm}L{5.0cm}}{%
    \tableauHeader{Niveau} &
    \tableauHeader{Question} &
    \tableauHeader{Référentiels}
}
Portefeuille & Quels projets, dans quel ordre, avec quels moyens~? & ISO~21504, MoP, COBIT APO05 \\
Programme & Comment obtenir un bénéfice par plusieurs projets liés~? & MSP, ISO~21503 \\
Projet & Comment livrer dans les délais, le coût et la qualité~? & ISO~21502, PMBOK, PRINCE2 \\
\end{tableaucours}
\vspace{-0.8em}
\begin{center}{\small\itshape Tableau~4.3. Portefeuille, programme, projet}\end{center}

% ============================================================================
\section{Maîtriser : Risque, Sécurité, Continuité, Intelligence Artificielle}
\label{sec:c04-maitrise}
% ============================================================================

\begin{tableaucours}{L{3.8cm}L{4.0cm}L{4.2cm}}{%
    \tableauHeader{Norme} &
    \tableauHeader{Objet} &
    \tableauHeader{Place dans le schéma directeur}
}
ISO~31000:2018~\cite{iso31000} & Gestion du risque, lignes directrices & Risques du schéma directeur lui-même \\
ISO/IEC~27001:2022~\cite{iso27001} & Sécurité de l'information, certifiable & Volet sécurité de la cible \\
ISO~22301:2019~\cite{iso22301} & Continuité d'activité, certifiable & Secours et reprise \\
ISO/IEC~42001:2023~\cite{iso42001} & Système de management de l'IA, certifiable & Projets d'intelligence artificielle \\
ISO~56002:2019~\cite{iso56002} & Management de l'innovation & Expérimentations \\
EN~16234-1:2019~\cite{ecf} & Référentiel européen des compétences numériques & Volet compétences \\
\end{tableaucours}
\vspace{-0.8em}
\begin{center}{\small\itshape Tableau~4.4. Normes de maîtrise utiles au schéma directeur}\end{center}

% ============================================================================
\section{Choisir ses Référentiels}
\label{sec:c04-choisir}
% ============================================================================

\begin{boxprocedure}
\textbf{Déclarer les référentiels d'un schéma directeur.}
\begin{enumerate}
    \item Recenser les textes obligatoires qui s'appliquent à l'organisation (Chapitre~\ref{chap:droit}).
    \item Recenser les référentiels exigés par un contrat, un bailleur ou un superviseur.
    \item Choisir, pour chaque besoin restant, un seul référentiel volontaire.
    \item Écrire dans le document la liste retenue, l'édition de chaque référentiel et son usage.
    \item Ne jamais écrire « conforme à » pour un cadre volontaire. Écrire « établi selon ».
\end{enumerate}
\end{boxprocedure}

\begin{boxscenario}{Mission ANPP. Les référentiels retenus, 26 janvier 2027}
\textbf{Constatations.} Le cahier des charges du conseil cite « les bonnes pratiques » sans les nommer. Le
bailleur qui finance une partie des projets demande un suivi de portefeuille. La division informatique compte
deux agents certifiés ITIL et aucun architecte.

\textbf{Analyse} (confiance élevée). L'équipe retient l'ISO/IEC~38500 pour le rôle du conseil, COBIT~2019 pour
les six objectifs du Tableau~4.2, l'ISO~21504 pour le portefeuille et l'ISO~31000 pour les risques. TOGAF n'est
pas retenu comme méthode complète, faute d'architecte~: la cible sera décrite par des cartographies simples.
Ce choix est écrit dans la note de cadrage.
\end{boxscenario}

\begin{boxlecture}{Les référentiels à l'ANPP avant la mission}
\begin{tableaucours}{L{7.6cm}cL{3.4cm}}{%
    \tableauHeader{Exigence} &
    \tableauHeader{Cote} &
    \tableauHeader{Preuve}
}
Les référentiels employés sont nommés avec leur édition & \NON{} & Schéma directeur 2021 \\
Le conseil exerce les trois tâches de l'ISO/IEC~38500 & \NON{} & Approuve sans surveiller \\
Un processus de portefeuille existe & \NON{} & Aucune procédure \\
Un registre des risques informatiques est tenu & \OUI{} & Registre de 2026 \\
La gestion des services suit un cadre déclaré & \NEK{} & Pratiques ITIL non documentées \\
Une certification est détenue sur un périmètre & \NON{} & Aucune \\
\end{tableaucours}

\medskip
\textbf{Décompte.} Un \OUI{}, quatre \NON{}, un \NEK{} sur six.
\hfill\lecture{0{,}167}{0{,}167}{1{,}25}

\medskip
\textbf{Lecture.} Un seul \OUI{}. Le registre des risques existe depuis la panne de mars~2026. Le \NEK{} est de preuve~:
deux agents appliquent ITIL sans procédure écrite. L'Agence n'a pas besoin de certificat à ce stade. Elle a
besoin de nommer ce qu'elle emploie.
\end{boxlecture}

\begin{boxmetier}{Administrateur}
Demander chaque année trois choses sur les technologies~: ce qui a été évalué, ce qui a été décidé, ce qui a été constaté.
\end{boxmetier}

\begin{boxmetier}{Directeur des systèmes d'information}
Tenir la liste des référentiels employés, avec leur édition. Retirer de tout document les mentions qui ne reposent sur aucune pièce.
\end{boxmetier}

\begin{boxmetier}{Consultant}
Choisir peu de référentiels et les employer vraiment. Quatre référentiels déclarés et suivis valent mieux que dix cités.
\end{boxmetier}

\begin{boxresume}
\begin{itemize}
    \item Un cadre volontaire se déclare, une norme se certifie, un texte oblige.
    \item ISO/IEC~38500:2024~: le conseil évalue, dirige et surveille.
    \item COBIT~2019~: EDM01, EDM02, APO02, APO03, APO05 et APO06 portent le schéma directeur.
    \item TOGAF donne la cible et les architectures de transition.
    \item ITIL~4 rappelle qu'un projet livré devient un service à exploiter.
    \item L'ISO~21504 traite le portefeuille, l'ISO~21502 le projet.
\end{itemize}
\end{boxresume}

\begin{boxinfo}
\textbf{Points de contrôle.} Vous avez compris ce chapitre si vous savez~:
\begin{itemize}
    \item nommer les trois tâches de l'ISO/IEC~38500~;
    \item dire quand un cadre volontaire devient contraignant~;
    \item distinguer portefeuille, programme et projet.
\end{itemize}
\end{boxinfo}

\begin{boxexo}[title={\ding{45}\quad Exercice~4.1. Quelle nature \niveaubase}]
Classez dans le Tableau~4.1~: COBIT~2019, ISO/IEC~27001, le code des marchés publics, ISO~31000, TOGAF, ISO~22301, un règlement du superviseur bancaire, ITIL~4.
\end{boxexo}

\begin{boxexo}[title={\ding{45}\quad Exercice~4.2. Trois tâches \niveaumoyen}]
Reprenez le bilan du prologue. Pour chacune des trois tâches de l'ISO/IEC~38500, dites ce que le conseil de l'ANPP a fait et ce qu'il aurait dû faire.
\end{boxexo}

\begin{boxexo}[title={\ding{45}\quad Exercice~4.3. La couverture \niveaumoyen}]
Réécrivez la mention de couverture du schéma directeur de la BANQUE SANAGA pour qu'elle soit exacte. Dites quelles pièces doivent l'accompagner.
\end{boxexo}

\begin{boxexo}[title={\ding{45}\quad Exercice~4.4. Cascade \niveauexpert}]
Appliquez la cascade des objectifs de COBIT à l'objectif « payer chaque bénéficiaire à date fixe ». Donnez un objectif d'alignement et deux objectifs de gestion. Indiquez la pièce qui prouverait chacun.
\end{boxexo}

\bigskip

\begin{boxinfo}
\textbf{Pour aller plus loin.}
\begin{itemize}
    \item L'ISO/IEC~38500~\cite{iso38500} et l'ISO~37000~\cite{iso37000}.
    \item COBIT~2019~\cite{cobit2019}, objectifs APO02 et APO05.
    \item La norme TOGAF~\cite{togaf10}.
    \item L'ISO~21504~\cite{iso21504}.
\end{itemize}
\end{boxinfo}

\begin{boxinfo}
\textbf{Transition.} Les normes viennent d'organismes techniques. Les États et les organisations internationales ont aussi leurs doctrines et leurs stratégies numériques. Le Chapitre~\ref{chap:doctrines} les présente, de l'OCDE à l'Union africaine.
\end{boxinfo}

        % ============================================================================
% Fichier  : partie2/P02_C05_doctrines_strategies.tex
% Titre    : Doctrines et Stratégies Numériques dans le Monde et en Afrique
% Auteur   : MINKA MI NGUIDJOÏ Thierry Emmanuel
% Version  : 1.0 -- novembre 2026
% Corrigés : corriges/Corr_C05.tex
% ============================================================================

\chapter{Doctrines et Stratégies Numériques dans le Monde et en Afrique}
\label{chap:doctrines}
\chaptermark{Doctrines et stratégies}

\epigraph{%
    « Un État ne numérise pas service par service.\\
    Il construit des socles. »%
}{Doctrine de l'infrastructure publique numérique}

\niveauChapitre{Fondamental}{%
    Chapitre~\ref{chap:normes}
}

\begin{boxobjectifs}
\begin{itemize}
    \item Citer la recommandation de l'OCDE de 2014 et le Pacte numérique mondial de 2024.
    \item Employer trois indices de comparaison dans un diagnostic.
    \item Présenter les six textes de l'Union africaine utiles à une stratégie numérique.
    \item Décrire la méthode RACINES, l'urbanisation et le pilotage par la valeur.
    \item Définir infrastructure publique numérique, souveraineté et sobriété.
\end{itemize}
\end{boxobjectifs}

\bigskip

% ============================================================================
\section*{Avant de commencer : « le bailleur parle une autre langue »}
% ============================================================================

Yaoundé, lundi 1\textsuperscript{er}~février~2027. M.~OWONA transmet à l'équipe la convention de financement
signée avec un bailleur. Trois expressions y reviennent~: « infrastructure publique numérique », « Pacte numérique
mondial », « interopérabilité ». Le directeur financier de l'Agence, M.~ABDOULAYE, avoue ne pas savoir ce que le
bailleur attend au juste. « Nous avons toujours présenté nos projets un par un. »

Mme~EKANI explique que le vocabulaire a changé. Les bailleurs ne financent plus une application de plus. Ils
financent des projets qui se branchent sur des socles communs. Un schéma directeur qui ignore ce vocabulaire ne
sera pas financé.

\begin{boxalert}
\textbf{Violation typique.} Écrire un schéma directeur d'administration comme si l'organisation était seule.

\smallskip
\textbf{Règle~:} une administration s'inscrit dans une stratégie nationale, elle-même inscrite dans des
engagements continentaux et mondiaux. Le document cite ces niveaux et dit comment ses projets s'y rattachent.
\end{boxalert}

% ============================================================================
\section{Les Organisations Internationales}
\label{sec:c05-monde}
% ============================================================================

\begin{boxjuris}
\textbf{Recommandation du Conseil de l'OCDE sur les stratégies numériques gouvernementales}, adoptée le
15~juillet~2014 (OECD/LEGAL/0406)~\cite{ocde2014}.

\smallskip
\textbf{Portée.} C'est le premier instrument international consacré à l'administration numérique. Il s'organise
en trois piliers et douze principes. Il demande de passer de l'administration \textit{électronique}, qui
met en ligne des procédures existantes, à l'administration \textit{numérique}, qui repense le service. Le
Cameroun n'est pas membre de l'OCDE~: le texte ne l'oblige pas. Il sert de référence.
\end{boxjuris}

Trois instruments de mesure permettent à un État de se situer.

\begin{tableaucours}{L{3.4cm}L{2.6cm}L{6.0cm}}{%
    \tableauHeader{Instrument} &
    \tableauHeader{Auteur} &
    \tableauHeader{Ce qu'il mesure}
}
Indice de développement de l'administration en ligne (EGDI) & Nations Unies & Services en ligne, infrastructures de télécommunication, capital humain \\
Indice de maturité GovTech & Banque mondiale & Systèmes centraux de l'État, services en ligne, participation, cadres habilitants \\
Indice de l'administration numérique & OCDE & Mise en œuvre de la recommandation de 2014 dans ses pays membres et partenaires \\
\end{tableaucours}
\vspace{-0.8em}
\begin{center}{\small\itshape Tableau~5.1. Trois instruments de comparaison}\end{center}

Ces indices ne disent pas quoi faire. Ils servent au diagnostic~: un schéma directeur d'administration cite le
rang du pays, la composante la plus faible et ce que l'organisation peut y changer à son échelle.

\begin{boxjuris}
\textbf{Pacte numérique mondial}, annexe~I de la résolution 79/1 de l'Assemblée générale des Nations Unies,
adoptée le 22~septembre~2024~\cite{onu2024pacte}.

\smallskip
\textbf{Portée.} Le Pacte fixe des objectifs communs~: réduire les fractures numériques, élargir les bénéfices
de l'économie numérique, garantir un espace numérique sûr et respectueux des droits, gouverner les données de
façon équitable et encadrer l'intelligence artificielle. C'est un engagement politique des États, sans
sanction. Il oriente les financements des bailleurs.
\end{boxjuris}

\subsection{L'infrastructure publique numérique}

Une doctrine s'est imposée depuis 2023 sous le nom d'\termecle{infrastructure publique numérique}. Elle désigne
des socles partagés par toutes les administrations et ouverts au secteur privé~: l'identité numérique, le
paiement, l'échange de données. L'idée est qu'un État ne numérise pas service par service. Il construit des
socles, puis chaque service s'y branche.

Pour le schéma directeur d'une agence, la conséquence est directe. Avant de bâtir son propre système
d'identification des bénéficiaires, l'ANPP doit savoir ce que l'État prévoit en matière d'identité. Bâtir seul
ce qui existera en commun est le gaspillage le plus fréquent des administrations.

% ============================================================================
\section{L'Afrique}
\label{sec:c05-afrique}
% ============================================================================

\begin{tableaucours}{L{4.2cm}L{2.6cm}L{5.2cm}}{%
    \tableauHeader{Texte} &
    \tableauHeader{Date} &
    \tableauHeader{Apport}
}
Convention de Malabo~\cite{malabo2014} & 27 juin 2014 & Transactions électroniques, données personnelles, cybersécurité \\
Stratégie de transformation numérique pour l'Afrique~\cite{ua2020dts} & 2020, pour 2020-2030 & Marché unique numérique d'ici 2030, piliers et secteurs prioritaires \\
Cadre stratégique en matière de données~\cite{ua2022donnees} & Février 2022 & Gouvernance et circulation des données \\
Protocole de la ZLECAf sur le commerce numérique~\cite{zlecaf2024} & Février 2024, annexes en février 2025 & Règles communes du commerce numérique \\
Stratégie continentale sur l'intelligence artificielle~\cite{ua2024ia} & 18 et 19 juillet 2024 & Cadres nationaux, compétences, usage responsable \\
Pacte numérique africain~\cite{ua2024pacte} & Juillet 2024 & Position commune de l'Afrique, adossée au Pacte mondial \\
\end{tableaucours}
\vspace{-0.8em}
\begin{center}{\small\itshape Tableau~5.2. Six textes de l'Union africaine}\end{center}

La Stratégie de transformation numérique pour l'Afrique est le texte pivot. Elle vise un marché unique numérique
d'ici 2030 et demande aux États de mettre leurs politiques en cohérence. Elle s'appuie sur des fondations
(infrastructures, compétences, cadre réglementaire, innovation) et désigne des secteurs prioritaires, parmi
lesquels l'administration, la finance, la santé, l'éducation, l'agriculture et le commerce.

Le protocole de la ZLECAf sur le commerce numérique a été adopté par la Conférence de l'Union africaine en
février~2024. Ses huit annexes, adoptées le 15~février~2025, couvrent notamment les identités numériques, les
paiements transfrontaliers et les transferts de données. Il n'entre en vigueur qu'après ratification par
vingt-deux États. Une entreprise comme LiZbiZ, qui livre au Tchad et en République centrafricaine, doit le
suivre.

\begin{boxinfo}
\textbf{Internalisation.} Ces textes ne s'appliquent pas directement à une organisation. Ils passent par le
droit national. La loi camerounaise de 2024 sur les données personnelles reprend des principes de Malabo. La
stratégie nationale d'intelligence artificielle de 2025 cite la stratégie continentale. Le consultant lit donc
les textes continentaux pour comprendre où va le droit national.
\end{boxinfo}

% ============================================================================
\section{La Doctrine Française du Schéma Directeur}
\label{sec:c05-france}
% ============================================================================

Le mot « schéma directeur » vient de l'administration française, qui l'a emprunté à l'urbanisme. Dans les
années 1980, le ministère français de l'Industrie diffuse la méthode RACINES, pour « rationalisation des choix
informatiques ». Elle fixe la démarche que la plupart des cabinets d'Afrique francophone appliquent encore.

\begin{boxprocedure}
\textbf{Les cinq étapes de la méthode RACINES.}
\begin{enumerate}
    \item Évaluation de l'existant et orientations.
    \item Scénarios.
    \item Plan d'action.
    \item Mise en œuvre.
    \item Suivi.
\end{enumerate}
\end{boxprocedure}

La méthode installe trois instances~: un comité directeur qui décide, un groupe de projet qui travaille et des
groupes d'utilisateurs qui expriment les besoins. Le cours en garde la structure et y ajoute ce qui lui
manque~: le financement et la mesure de la valeur.

\subsection{L'urbanisation}

Dans les années 1990, l'\termecle{urbanisation} du système d'information transpose le vocabulaire de la ville.
Le système est découpé en zones, quartiers et îlots. Un plan d'occupation des sols fixe ce qui peut se
construire et où. Les règles sont simples~: un îlot rend un service et un seul, les échanges passent par des
interfaces, une donnée a un seul maître. Longépé en a donné l'exposé de référence~\cite{longepe2001}.

\subsection{Le pilotage par la valeur}

Dans les années 2000, l'administration française se dote d'une méthode d'analyse de la valeur des projets, MAREVA.
Elle note un projet sur plusieurs axes~: la rentabilité pour l'État, la maîtrise du risque, la valeur pour les
usagers, la valeur pour les agents et la nécessité. L'intérêt est de ne pas réduire un projet public à son
rendement financier. Le Chapitre~\ref{chap:portefeuille} en reprend le principe.

Le Cigref, association de grandes entreprises et administrations françaises, publie des travaux sur le pilotage
économique du système d'information et une nomenclature des métiers. Ce sont des documents de praticiens, utiles
pour bâtir un budget et un plan de compétences.

% ============================================================================
\section{Deux Doctrines Récentes}
\label{sec:c05-recent}
% ============================================================================

La \termecle{souveraineté numérique} désigne la capacité d'un État ou d'une organisation à garder la maîtrise de
ses données, de ses infrastructures et de ses choix techniques. Elle pèse sur le choix d'un hébergement. Elle ne
se confond pas avec l'autarcie~: héberger chez soi un système que l'on ne sait pas exploiter n'est pas une
souveraineté.

La \termecle{sobriété numérique} demande de mesurer l'empreinte matérielle et énergétique du système
d'information et d'en tenir compte dans les choix. Dans un pays où l'électricité est irrégulière et le matériel
importé, c'est aussi une affaire de coût~: allonger la durée de vie des postes et réduire le nombre de serveurs
relèvent de la bonne gestion avant de relever de l'écologie.

% ============================================================================
\section{Ce que le Consultant en Retient}
\label{sec:c05-usage}
% ============================================================================

\begin{boxscenario}{Mission ANPP. La note de contexte international, 2 février 2027}
\textbf{Constatations.} Le bailleur qui cofinance les projets de l'Agence cite le Pacte numérique mondial et
l'infrastructure publique numérique dans sa convention. Le schéma directeur de 2021 ne mentionnait aucun texte
international. L'Agence a développé son propre identifiant des bénéficiaires.

\textbf{Analyse} (confiance moyenne). Le nouveau document doit montrer au bailleur que les projets de l'Agence
s'inscrivent dans des socles communs. L'identifiant propre à l'Agence est un point de fragilité~: il faudra le
relier à l'identité nationale quand elle sera disponible. Cette dépendance est inscrite au registre des risques.
\end{boxscenario}

\begin{boxmetier}{Consultant}
Lire la convention du bailleur avant d'écrire la première ligne. Son vocabulaire dit ce qui sera financé.
\end{boxmetier}

\begin{boxmetier}{Dirigeant d'administration}
Demander quels socles communs l'État prépare avant d'approuver un projet d'identification, de paiement ou d'échange de données.
\end{boxmetier}

\begin{boxmetier}{Directeur des systèmes d'information}
Concevoir chaque application pour qu'elle puisse se brancher plus tard sur un socle commun~: interfaces documentées, identifiants convertibles.
\end{boxmetier}

\begin{boxresume}
\begin{itemize}
    \item OCDE, 15~juillet~2014~: passer de l'administration électronique à l'administration numérique.
    \item Pacte numérique mondial, 22~septembre~2024~: engagement politique des États.
    \item Union africaine~: stratégie 2020-2030, cadre des données de 2022, stratégie d'intelligence artificielle et Pacte numérique africain de 2024, protocole de la ZLECAf.
    \item RACINES~: existant et orientations, scénarios, plan d'action, mise en œuvre, suivi.
    \item Urbanisation~: zones, quartiers, îlots, une donnée a un seul maître.
    \item Infrastructure publique numérique~: identité, paiement, échange de données comme socles communs.
\end{itemize}
\end{boxresume}

\begin{boxinfo}
\textbf{Points de contrôle.} Vous avez compris ce chapitre si vous savez~:
\begin{itemize}
    \item dire pourquoi la recommandation de l'OCDE n'oblige pas le Cameroun~;
    \item citer trois socles de l'infrastructure publique numérique~;
    \item énumérer les cinq étapes de RACINES.
\end{itemize}
\end{boxinfo}

\begin{boxexo}[title={\ding{45}\quad Exercice~5.1. Dates \niveaubase}]
Associez chaque texte à sa date~: recommandation de l'OCDE, Pacte numérique mondial, Convention de Malabo, stratégie continentale sur l'intelligence artificielle, protocole de la ZLECAf sur le commerce numérique.
\end{boxexo}

\begin{boxexo}[title={\ding{45}\quad Exercice~5.2. Socles \niveaumoyen}]
Pour l'ANPP, dites quels projets du futur schéma directeur dépendent d'un socle d'identité, d'un socle de paiement et d'un socle d'échange de données. Que faire tant que le socle n'existe pas~?
\end{boxexo}

\begin{boxexo}[title={\ding{45}\quad Exercice~5.3. RACINES aujourd'hui \niveaumoyen}]
Comparez les cinq étapes de RACINES aux six pièces du schéma directeur (Chapitre~\ref{chap:vocabulaire}). Qu'est-ce que la méthode des années 1980 ne traite pas~?
\end{boxexo}

\begin{boxexo}[title={\ding{45}\quad Exercice~5.4. Souveraineté \niveauexpert}]
La BANQUE SANAGA hésite entre un hébergement à Douala qu'elle exploiterait elle-même et un hébergement chez un prestataire étranger certifié. Rédigez une page qui distingue souveraineté, sécurité et coût, sans conclure à la place de la banque.
\end{boxexo}

\bigskip

\begin{boxinfo}
\textbf{Pour aller plus loin.}
\begin{itemize}
    \item La recommandation de l'OCDE~\cite{ocde2014}.
    \item La Stratégie de transformation numérique pour l'Afrique~\cite{ua2020dts}.
    \item Le Pacte numérique africain~\cite{ua2024pacte}.
    \item Longépé~\cite{longepe2001}, pour l'urbanisation.
\end{itemize}
\end{boxinfo}

\begin{boxinfo}
\textbf{Transition.} Le monde et le continent ont leurs textes. Le Cameroun a les siens, avec ses plans, ses acteurs et ses hésitations. Le Chapitre~\ref{chap:cameroun} les présente.
\end{boxinfo}


    \part{Le contexte national}
    \label{part:national}

        % ============================================================================
% Fichier  : partie3/P03_C06_strategie_cameroun.tex
% Titre    : La stratégie numérique du Cameroun
% Auteur   : MINKA MI NGUIDJOÏ Thierry Emmanuel
% Version  : 1.0 -- novembre 2026
% Corrigés : corriges/Corr_C06.tex
% ============================================================================

\chapter{La stratégie numérique du Cameroun}
\label{chap:cameroun}
\chaptermark{Stratégie du Cameroun}

\epigraph{%
    « Un plan attendu ne se cite pas comme un plan en vigueur.\\
    Il ne s'ignore pas non plus. »%
}{Principe du cours}

\niveauChapitre{Fondamental}{%
    Chapitre~\ref{chap:doctrines}
}

\begin{boxobjectifs}
\begin{itemize}
    \item Situer la Vision 2035, la SND30 et le plan Cameroun numérique 2020.
    \item Dire l'état exact du plan Cameroun numérique 2026-2030.
    \item Présenter la stratégie nationale d'intelligence artificielle et une stratégie sectorielle.
    \item Nommer les acteurs publics du numérique et leur rôle.
    \item Établir le rattachement stratégique d'un schéma directeur.
\end{itemize}
\end{boxobjectifs}

\bigskip

% ============================================================================
\section*{Avant de commencer : « conformément au plan Cameroun numérique »}
% ============================================================================

Yaoundé, jeudi 4~février~2027. Le projet de note que le consultant junior soumet à Mme~EKANI commence ainsi~:
« Conformément au plan Cameroun numérique 2026-2030, l'ANPP engage sa transformation. » Mme~EKANI demande où
il a lu ce plan. Il en a entendu parler dans un séminaire.

Elle lui fait chercher le texte d'approbation. Il n'y en a pas. Le plan est finalisé et déposé. Le Premier
ministre ne l'a pas encore approuvé. La phrase du consultant est donc fausse~: on ne se conforme pas à un texte
qui n'est pas en vigueur.

\begin{boxalert}
\textbf{Violation typique.} Rattacher un schéma directeur à un document d'orientation que l'on n'a pas lu ou
qui n'est pas approuvé.

\smallskip
\textbf{Règle~:} chaque document cité l'est avec sa date, son auteur et son état~: en vigueur, échu ou
attendu. Le rattachement se fait sur ce qui est en vigueur.
\end{boxalert}

% ============================================================================
\section{La chaîne des documents d'orientation}
\label{sec:c06-chaine}
% ============================================================================

\begin{tableaucours}{L{4.0cm}L{2.4cm}L{5.6cm}}{%
    \tableauHeader{Document} &
    \tableauHeader{Date} &
    \tableauHeader{Place du numérique}
}
Cameroun Vision 2035~\cite{vision2035} & 2009 & Cadre de long terme~: devenir un pays émergent, démocratique et uni dans sa diversité \\
Document de stratégie pour la croissance et l'emploi & 2009, pour 2010-2019 & Première phase de la Vision \\
Plan stratégique Cameroun numérique 2020~\cite{camnum2020} & Mai 2016 & Premier document national consacré à l'économie numérique \\
Stratégie nationale de développement, SND30~\cite{snd30} & 2020, pour 2020-2030 & Deuxième phase de la Vision. Le numérique parmi les sous-secteurs retenus \\
Stratégie nationale d'intelligence artificielle~\cite{snia2025} & Juillet 2025 & Sept piliers, dont la gouvernance et la souveraineté \\
Plan stratégique national de santé numérique~\cite{psnsn2026} & Mars 2026, pour 2026-2030 & Stratégie d'un secteur, huit axes, 29~milliards de FCFA \\
\end{tableaucours}
\vspace{-0.8em}
\begin{center}{\small\itshape Tableau~6.1. Les documents d'orientation du Cameroun}\end{center}

La Vision 2035 a été élaborée par le ministère chargé de l'économie et de la planification. Elle se met en œuvre
par phases de dix ans. La Stratégie nationale de développement 2020-2030, dite SND30, en est la deuxième phase.
Elle vise une transformation structurelle de l'économie et compte le numérique parmi les sous-secteurs sur
lesquels elle s'appuie. Elle demande la modernisation de l'administration et le développement de la
gouvernance électronique.

Tout schéma directeur d'une administration camerounaise se rattache donc à la SND30. C'est le premier maillon de
la démonstration d'alignement~: la mission de l'organisation découle d'un programme, qui découle d'une stratégie
sectorielle, qui découle de la SND30.

% ============================================================================
\section{Cameroun numérique 2020, son bilan et sa suite}
\label{sec:c06-plan}
% ============================================================================

Le Plan stratégique Cameroun numérique 2020 a été publié par le ministère des Postes et Télécommunications en
mai~2016. Sa vision était de faire du Cameroun un pays numérique à l'horizon 2020. Il s'organisait autour du
développement des infrastructures à large bande, de la production de contenus numériques, de la transformation
numérique de l'administration et des entreprises, de la confiance numérique et des compétences.

Son exécution a été partielle. La SND30 relève elle-même que le niveau d'utilisation des technologies reste
faible malgré les investissements. Le plan a pourtant installé un vocabulaire et des acteurs que l'on retrouve
dans tous les documents suivants.

\begin{boxinfo}
\textbf{Le plan 2026-2030.} Le Plan de développement de l'économie numérique au Cameroun 2026-2030, édition
2025~\cite{pden2026}, est finalisé et déposé. À la date de ce cours, il est en cours de validation dans les
services du Premier ministre et n'a pas encore reçu son approbation. Il n'est donc pas en cours de mise en
œuvre. Un schéma directeur rédigé aujourd'hui le mentionne dans ces termes exacts, se rattache à la SND30, qui
est en vigueur et prévoit de se réviser à l'approbation du plan. La section~\ref{sec:c06-projet} en décrit le
contenu.
\end{boxinfo}

Cette situation est fréquente et il faut savoir la traiter. Un document d'orientation attendu ne se cite pas
comme une source en vigueur. Il ne s'ignore pas non plus. Le schéma directeur écrit ce qui est établi, ce qui est
attendu et ce qui changera.

\subsection{Le projet d'accélération de la transformation numérique}

Le Projet d'accélération de la transformation numérique du Cameroun, le PATNUC, est conduit par le ministère des
Postes et Télécommunications avec l'appui d'un bailleur. Il comprend un volet de connectivité des zones rurales
et un volet d'appui au cadre stratégique, dont est issue la stratégie d'intelligence artificielle. Une
administration qui prépare son schéma directeur vérifie si l'un de ses projets peut s'y inscrire.

% ============================================================================
\section{La stratégie d'intelligence artificielle}
\label{sec:c06-ia}
% ============================================================================

La Stratégie nationale d'intelligence artificielle a été présentée lors de la deuxième édition des Concertations
nationales sur l'intelligence artificielle, à Yaoundé, les 7 et 8~juillet~2025. Elle repose sur sept piliers,
parmi lesquels la gouvernance de l'intelligence artificielle et la souveraineté numérique, les données et
l'infrastructure et une intelligence artificielle multilingue et inclusive. Elle annonce une loi sur
l'intelligence artificielle et une autorité compétente.

Pour un schéma directeur, deux conséquences. Tout projet d'intelligence artificielle se rattache à cette
stratégie. Tout projet doit prévoir qu'un cadre légal viendra~: on documente dès maintenant les données
employées, les décisions prises et les contrôles humains (Chapitre~\ref{chap:donnees}).

% ============================================================================
\section{Les acteurs}
\label{sec:c06-acteurs}
% ============================================================================

\begin{tableaucours}{L{3.6cm}L{8.4cm}}{%
    \tableauHeader{Acteur} &
    \tableauHeader{Rôle pour un schéma directeur}
}
Ministère des Postes et Télécommunications & Politique du secteur, plans nationaux, tutelle des agences du numérique \\
ANTIC & Agence nationale des technologies de l'information et de la communication. Sécurité des systèmes d'information, audits de sécurité, expertise auprès des administrations, interopérabilité \\
ART & Agence de régulation des télécommunications. Opérateurs, licences, qualité de service \\
CAMTEL et opérateurs & Infrastructures, connectivité, services de paiement par téléphone \\
Ministère chargé de l'économie et de la planification & SND30, programmation des investissements publics \\
Ministère des Finances & Budget, cadre de dépenses à moyen terme, systèmes financiers de l'État \\
Ministère des Marchés publics et ARMP & Passation et régulation des marchés \\
Autorité de protection des données & Prévue par la loi n°~2024/017. Autorisations et contrôle des traitements \\
Contrôle supérieur de l'État & Contrôle des projets et de la gestion publics \\
\end{tableaucours}
\vspace{-0.8em}
\begin{center}{\small\itshape Tableau~6.2. Les acteurs publics du numérique}\end{center}

L'ANTIC est régie par le décret n°~2019/150 du 22~mars~2019~\cite{decret2019150}, qui a remplacé le décret
n°~2012/180 du 10~avril~2012~\cite{decret2012180}. Le texte de 2012 la chargeait déjà de favoriser
l'interopérabilité entre les systèmes d'information et de fournir son expertise aux administrations pour la
conception de leurs outils. Elle conduit les audits de sécurité obligatoires prévus par le décret
n°~2012/1643/PM~\cite{decret20121643}.

\begin{boxinfo}
\textbf{Qui impose un schéma directeur~?} Aucun texte général n'impose un schéma directeur à toute
administration. L'obligation figure dans les textes organiques. Celui de la Présidence de la République charge
sa direction de l'informatique d'élaborer et de mettre en œuvre le schéma directeur d'informatisation des
services. Beaucoup de ministères ont une disposition voisine. Le schéma directeur 2020-2024 d'un ministère,
publié en ligne, se réfère aux prescriptions de l'ANTIC pour ses achats~\cite{minepded2020sdi}. Le consultant
cherche donc le texte organique de l'organisation avant d'écrire que le schéma directeur est obligatoire.
\end{boxinfo}

% ============================================================================
\section{Les stratégies sectorielles}
\label{sec:c06-secteurs}
% ============================================================================

La santé donne l'exemple le plus abouti. Un premier plan stratégique de santé numérique couvrait 2020-2024. Le
second, pour 2026-2030, a été lancé le 4~mars~2026. Il compte huit axes, dont la gouvernance, le cadre
juridique, les ressources humaines, les infrastructures et l'interopérabilité. Son budget est d'environ
29~milliards de francs CFA. Le ministère y reconnaît un niveau de maturité limité.

Un hôpital public qui rédige son schéma directeur se rattache à ce plan avant de se rattacher à la SND30. La
règle est générale~: le document le plus proche d'abord, puis la chaîne jusqu'à la stratégie nationale.

% ============================================================================
\section{Le projet de plan 2026-2030}
\label{sec:c06-projet}
% ============================================================================

Le projet de plan s'ouvre par un état des lieux. Il y constate que le plan Cameroun numérique 2020 a connu un
taux d'exécution très faible et il en donne cinq causes~: un cadre institutionnel et réglementaire inadapté,
des ressources financières insuffisantes, des ressources humaines qualifiées insuffisantes, une faible
coordination des activités transversales et une prise en compte insuffisante des questions de genre. Le
prologue de ce cours a fait le même constat à l'échelle d'une agence. Un plan national et un schéma directeur
d'organisme échouent pour les mêmes raisons.

Le document se compose de quatre parties~: un état des lieux, un diagnostic, un plan d'actions assorti d'un
plan d'actions prioritaires, un dispositif de mise en œuvre et de suivi-évaluation. Il se termine par une
déclaration de politique.

\begin{tableaucours}{L{3.6cm}L{8.4cm}}{%
    \tableauHeader{Élément} &
    \tableauHeader{Contenu du projet de plan}
}
Objectif global & Développer l'économie numérique au Cameroun \\
Cinq piliers & Infrastructures des communications électroniques, plateformes numériques, services financiers numériques, entrepreneuriat numérique, compétences numériques \\
Quatre axes & Amélioration de la gouvernance~; renforcement du cadre juridique et institutionnel~; exploitation de l'économie numérique~; utilisation du numérique \\
Mise en œuvre & À court terme, un comité interministériel de l'économie numérique. À moyen terme, la création d'une agence nationale de l'économie numérique \\
Suivi-évaluation & Suivi des actions, évaluation de l'impact, tableau de bord \\
Financement & Un fonds d'investissement pour l'économie numérique, des partenariats public-privé, des alliances internationales et des fonds de garantie, un programme de subventions pour l'inclusion numérique \\
\end{tableaucours}
\vspace{-0.8em}
\begin{center}{\small\itshape Tableau~6.3. Le projet de Plan de développement de l'économie numérique 2026-2030}\end{center}

Trois enseignements pour un schéma directeur. Le plan se rattache lui-même au premier pilier de la SND30, la
transformation structurelle de l'économie~: la chaîne de rattachement de ce chapitre est donc la bonne. Il
prévoit de réviser des textes, dont la loi de 2010 sur la cybersécurité~: la carte d'applicabilité du
Chapitre~\ref{chap:droit} devra être relue à son approbation. Il annonce des institutions nouvelles~: la liste
des acteurs du Tableau~6.2 changera.

Tant que l'approbation n'est pas donnée, aucun de ces éléments n'oblige. Le schéma directeur de l'ANPP les
note comme des évolutions attendues et n'en tire aucun projet.

% ============================================================================
\section{Rattacher un schéma directeur}
\label{sec:c06-rattacher}
% ============================================================================

\begin{boxprocedure}
\textbf{Établir le rattachement stratégique.}
\begin{enumerate}
    \item Trouver le texte qui crée l'organisation et fixe ses missions.
    \item Trouver le programme budgétaire auquel elle contribue.
    \item Trouver la stratégie de son secteur, si elle existe, avec sa période.
    \item Citer la SND30 et l'objectif concerné.
    \item Noter les documents attendus et non approuvés, avec leur état exact.
    \item Écrire le tout en une page, en tête du schéma directeur.
\end{enumerate}
\end{boxprocedure}

\begin{boxscenario}{Mission ANPP. Le rattachement stratégique, 5 février 2027}
\textbf{Constatations.} Le décret de création de l'Agence fixe trois missions~: liquider les prestations, les
payer, tenir le fichier des bénéficiaires. L'Agence contribue à un programme budgétaire de protection sociale.
Son secteur n'a pas de stratégie numérique propre. Le plan stratégique 2026-2030 de l'Agence, adopté en
décembre~2025, compte quatre objectifs.

\textbf{Analyse} (confiance élevée). Le rattachement se fait à quatre niveaux~: missions, plan stratégique de
l'Agence, programme budgétaire, SND30. Le plan Cameroun numérique 2026-2030 est noté comme attendu. L'absence de
stratégie sectorielle est écrite~: elle explique que l'Agence doive fixer seule ses principes d'interopérabilité.
\end{boxscenario}

\begin{boxlecture}{Le rattachement stratégique de l'ANPP}
\begin{tableaucours}{L{7.6cm}cL{3.4cm}}{%
    \tableauHeader{Exigence} &
    \tableauHeader{Cote} &
    \tableauHeader{Preuve}
}
Les missions sont tirées d'un texte cité & \OUI{} & Décret de création \\
Un plan stratégique de l'organisation est en vigueur & \OUI{} & Plan 2026-2030 \\
Le programme budgétaire est identifié & \OUI{} & Loi de finances \\
La stratégie du secteur est citée & \NEK{} & Aucune stratégie sectorielle \\
Le lien avec la SND30 est écrit & \NON{} & Absent du plan de l'Agence \\
Les documents attendus sont notés avec leur état & \NON{} & Aucune mention \\
\end{tableaucours}

\medskip
\textbf{Décompte.} Trois \OUI{}, deux \NON{}, un \NEK{} sur six.
\hfill\lecture{0{,}50}{0{,}167}{1{,}46}

\medskip
\textbf{Lecture.} Le socle existe~: trois \OUI{} portent sur des textes en vigueur. Le \NEK{} est de critère, puisqu'il
n'y a pas de stratégie sectorielle à citer. Les deux \NON{} se corrigent en une page. C'est le premier livrable
de la mission.
\end{boxlecture}

\begin{boxmetier}{Consultant}
Vérifier l'état de chaque document d'orientation à la date de la mission. Garder la preuve de la vérification.
\end{boxmetier}

\begin{boxmetier}{Dirigeant d'administration}
Faire inscrire dans le schéma directeur une clause de révision à l'approbation des plans nationaux attendus.
\end{boxmetier}

\begin{boxmetier}{Directeur des systèmes d'information}
Prendre l'attache de l'ANTIC au début de la démarche, pour la sécurité et l'interopérabilité.
\end{boxmetier}

\begin{boxresume}
\begin{itemize}
    \item La Vision 2035 date de 2009. La SND30 en est la deuxième phase, pour 2020-2030.
    \item Cameroun numérique 2020 date de mai~2016. Son exécution a été partielle.
    \item Le plan 2026-2030 est finalisé et déposé, sans approbation du Premier ministre à ce jour.
    \item La stratégie d'intelligence artificielle a été présentée les 7 et 8~juillet~2025.
    \item L'ANTIC est régie par le décret n°~2019/150 du 22~mars~2019.
    \item Aucun texte général n'impose un schéma directeur. L'obligation vient des textes organiques.
\end{itemize}
\end{boxresume}

\begin{boxinfo}
\textbf{Points de contrôle.} Vous avez compris ce chapitre si vous savez~:
\begin{itemize}
    \item dire sur quel document en vigueur une administration se rattache~;
    \item énoncer l'état du plan 2026-2030~;
    \item citer quatre acteurs publics et leur rôle.
\end{itemize}
\end{boxinfo}

\begin{boxexo}[title={\ding{45}\quad Exercice~6.1. En vigueur, échu, attendu \niveaubase}]
Classez~: Vision 2035, Cameroun numérique 2020, SND30, Cameroun numérique 2026-2030, plan de santé numérique 2020-2024, plan de santé numérique 2026-2030, stratégie d'intelligence artificielle de 2025.
\end{boxexo}

\begin{boxexo}[title={\ding{45}\quad Exercice~6.2. Une page \niveaumoyen}]
Rédigez la page de rattachement stratégique de l'ANPP, selon la démarche en six points.
\end{boxexo}

\begin{boxexo}[title={\ding{45}\quad Exercice~6.3. Qui saisir \niveaumoyen}]
L'ANPP veut ouvrir un paiement sur compte mobile. Quels acteurs du Tableau~6.2 sont concernés et pour quoi~?
\end{boxexo}

\begin{boxexo}[title={\ding{45}\quad Exercice~6.4. Texte organique \niveauexpert}]
Cherchez le texte organique d'un ministère ou d'un organisme public de votre choix. Dit-il qui élabore le schéma directeur~? Qui l'approuve~? Rédigez votre constat en séparant ce que le texte dit de ce que la pratique ajoute.
\end{boxexo}

\bigskip

\begin{boxinfo}
\textbf{Pour aller plus loin.}
\begin{itemize}
    \item La SND30~\cite{snd30}.
    \item Le plan Cameroun numérique 2020~\cite{camnum2020}.
    \item Le plan de santé numérique 2026-2030~\cite{psnsn2026}.
    \item Le décret n°~2019/150 sur l'ANTIC~\cite{decret2019150}.
\end{itemize}
\end{boxinfo}

\begin{boxinfo}
\textbf{Transition.} Les orientations disent où aller. Le droit dit comment une organisation a le droit d'y aller~: acheter, financer, décider, protéger. Le Chapitre~\ref{chap:droit} dresse le cadre juridique d'un schéma directeur.
\end{boxinfo}

        % ============================================================================
% Fichier  : partie3/P03_C07_cadre_juridique.tex
% Titre    : Le cadre juridique d'un schéma directeur
% Auteur   : MINKA MI NGUIDJOÏ Thierry Emmanuel
% Version  : 1.0 -- novembre 2026
% Corrigés : corriges/Corr_C07.tex
% ============================================================================

\chapter{Le cadre juridique d'un schéma directeur}
\label{chap:droit}
\chaptermark{Cadre juridique}

\epigraph{%
    « Une administration n'achète pas, ne finance pas\\
    et ne décide pas comme une entreprise. »%
}{Principe du cours}

\niveauChapitre{Fondamental}{%
    Chapitre~\ref{chap:cameroun}
}

\begin{boxobjectifs}
\begin{itemize}
    \item Citer les lois du numérique et dire ce que chacune impose au schéma directeur.
    \item Distinguer autorisations d'engagement et crédits de paiement.
    \item Dire quel organe approuve un schéma directeur.
    \item Énoncer quatre effets du code des marchés publics sur une trajectoire.
    \item Citer les textes communautaires applicables au système d'information d'une banque.
    \item Dresser une carte d'applicabilité.
\end{itemize}
\end{boxobjectifs}

\bigskip

% ============================================================================
\section*{Avant de commencer : « nous lancerons le projet en mars »}
% ============================================================================

Yaoundé, mercredi 12~février~2025. Le directeur des prestations veut un tableau de bord des paiements « pour
mars ». Mme~EKANI demande à M.~ABDOULAYE, directeur financier, si la dépense est inscrite au budget 2025. Elle ne
l'est pas. Elle demande si le marché figure à la programmation de l'année. Il n'y figure pas.

M.~ABDOULAYE fait le compte. Il faut une inscription au budget, donc une délibération du conseil. Il faut ensuite
un dossier d'appel d'offres, un délai de publication, une commission, une attribution, une notification. « Votre
mars est un mars 2026. » Le directeur des prestations découvre que le calendrier d'un projet public commence
bien avant le premier jour de travail.

\begin{boxalert}
\textbf{Violation typique.} Bâtir la trajectoire d'une administration avec des délais d'entreprise.

\smallskip
\textbf{Règle~:} dans une organisation soumise au droit public, la trajectoire compte trois délais avant le
démarrage d'un projet~: l'inscription budgétaire, la programmation du marché et la passation.
\end{boxalert}

% ============================================================================
\section{Cinq familles de règles}
\label{sec:c07-familles}
% ============================================================================

Un schéma directeur rencontre le droit à cinq moments~: quand il traite des données, quand il sécurise, quand il
finance, quand il achète et quand il fait décider.

\begin{tableaucours}{L{2.8cm}L{5.0cm}L{4.2cm}}{%
    \tableauHeader{Famille} &
    \tableauHeader{Textes principaux} &
    \tableauHeader{Question pour le schéma directeur}
}
Numérique & Lois n°~2010/012, 2010/013, 2010/021, 2023/009 et 2024/017 & Que peut-on traiter, échanger, mettre en ligne~? \\
Finances & Lois n°~2018/011 et 2018/012 & Comment inscrire une dépense sur plusieurs années~? \\
Organisation & Lois n°~2017/010 et 2017/011 & Qui approuve le document~? \\
Achat & Décrets n°~2018/366 et 2018/355 & Comment acquérir~? \\
Secteur & Règlements de la CEMAC et de la COBAC, droit OHADA & Que demande le superviseur ou le droit des affaires~? \\
\end{tableaucours}
\vspace{-0.8em}
\begin{center}{\small\itshape Tableau~7.1. Cinq familles de règles}\end{center}

% ============================================================================
\section{Les lois du numérique}
\label{sec:c07-numerique}
% ============================================================================

Trois lois ont été promulguées le même jour, le 21~décembre~2010.

La loi n°~2010/012 relative à la cybersécurité et à la cybercriminalité~\cite{loi2010012} impose aux exploitants
de systèmes d'information de les protéger, prévoit un audit de sécurité obligatoire, que le décret
n°~2012/1643/PM organise~\cite{decret20121643} et protège la vie privée dans les communications. Pour le schéma
directeur, elle fonde le volet sécurité et le budget de l'audit.

La loi n°~2010/013 régissant les communications électroniques, modifiée par la loi n°~2015/006 du
20~avril~2015~\cite{loi2010013}, encadre les réseaux et les services. Une organisation qui veut exploiter son
propre réseau entre ses sites, ou qui achète de la connectivité, s'y réfère.

La loi n°~2010/021 régissant le commerce électronique~\cite{loi2010021}, précisée par le décret n°~2011/1521/PM
du 15~juin~2011~\cite{decret20111521}, encadre les activités commerciales en ligne. Elle intéresse LiZbiZ dès
qu'elle vend ou contracte par voie électronique.

S'y ajoutent la loi-cadre n°~2011/012 du 6~mai~2011 sur la protection du consommateur~\cite{loi2011012} et la
loi n°~2023/009 du 25~juillet~2023 portant charte de protection des enfants en ligne~\cite{loi2023009}.

\begin{boxjuris}
\textbf{\loicam{n°~2024/017} du 23~décembre~2024 relative à la protection des données à caractère
personnel}~\cite{loi2024017}.

\smallskip
\textbf{Portée pour le schéma directeur.} La loi soumet les traitements à une autorisation préalable, encadre
les transferts hors du Cameroun, impose des mesures de sécurité et prévoit une évaluation d'impact pour les
traitements à risque élevé. Tout projet qui crée ou modifie un traitement de données personnelles porte donc
une tâche de conformité, un délai et un coût. Un schéma directeur qui ne les prévoit pas sous-estime chaque
projet. Le cours d'audit de protection des données (UE MSN-2S4-SEC02) en donne le détail.
\end{boxjuris}

% ============================================================================
\section{Les finances publiques}
\label{sec:c07-finances}
% ============================================================================

La loi n°~2018/012 du 11~juillet~2018 fixe le régime financier de l'État et des autres entités
publiques~\cite{loi2018012}. Elle organise le budget en programmes dotés d'objectifs et d'indicateurs et
distingue les \termecle{autorisations d'engagement}, qui permettent de s'engager sur plusieurs années, des
\termecle{crédits de paiement}, qui limitent ce qui peut être payé dans l'année. La loi n°~2018/011 du même jour
porte code de transparence et de bonne gouvernance~\cite{loi2018011}.

Deux conséquences. Un projet pluriannuel se chiffre en engagement total et en paiements annuels. Les deux
tableaux figurent au schéma directeur. Ensuite, un projet se rattache à un programme et à ses indicateurs. Un
projet qui ne contribue à aucun indicateur de programme a peu de chances d'être financé.

\begin{boxinfo}
\textbf{Le cadre de dépenses à moyen terme.} L'administration programme ses dépenses sur trois ans. La première
année du schéma directeur se cale sur le budget voté. Les deux suivantes se calent sur ce cadre. Au-delà, la
trajectoire est indicative, ce que le document doit écrire.
\end{boxinfo}

% ============================================================================
\section{Qui approuve}
\label{sec:c07-organes}
% ============================================================================

La loi n°~2017/010 du 12~juillet~2017 porte statut général des établissements publics~\cite{loi2017010}. La loi
n°~2017/011 du même jour porte statut général des entreprises publiques~\cite{loi2017011}. Elles ont remplacé la
loi de 1999. Dans les deux cas, un conseil d'administration définit les orientations et approuve les budgets et
les plans et un directeur général exécute.

Le schéma directeur engage l'organisation sur plusieurs exercices. Il relève donc du conseil. Un document
approuvé par le seul directeur général, ou par la seule direction informatique, n'engage pas les budgets à
venir. C'est l'une des causes du sort du document de 2019 à l'ANPP~: le conseil l'avait approuvé, mais aucune
délibération n'en tirait de conséquence budgétaire.

% ============================================================================
\section{Acheter}
\label{sec:c07-marches}
% ============================================================================

Le décret n°~2018/366 du 20~juin~2018 porte code des marchés publics~\cite{decret2018366}. Il s'applique à
l'État, aux collectivités et aux établissements publics. Les entreprises publiques relèvent du décret
n°~2018/355 du 12~juin~2018, qui fixe des règles communes à leurs marchés~\cite{decret2018355}.

Le code pèse sur le schéma directeur de quatre manières.

\begin{enumerate}
    \item \textbf{Le délai.} Entre la décision de lancer un projet et la notification d'un marché, il s'écoule
          plusieurs mois. La trajectoire doit les compter.
    \item \textbf{L'allotissement.} Découper un projet en lots change la concurrence et le calendrier. Le
          fractionnement destiné à éviter une procédure est interdit.
    \item \textbf{La programmation.} Les marchés de l'année sont programmés à l'avance. Un projet absent de la
          programmation ne peut pas être lancé.
    \item \textbf{La définition du besoin.} Le besoin est décrit par ses fonctions et ses performances, non par
          une marque. Le schéma directeur ne nomme donc pas de produit.
\end{enumerate}

\begin{boxalert}
\textbf{Violation typique.} Inscrire au schéma directeur le nom d'un éditeur ou d'un matériel.

\smallskip
\textbf{Règle~:} le schéma directeur décrit des besoins et des exigences. Le choix du produit résulte de la
mise en concurrence. Un document qui nomme la solution prépare un contentieux.
\end{boxalert}

% ============================================================================
\section{Le secteur bancaire}
\label{sec:c07-banque}
% ============================================================================

Un établissement de crédit de la CEMAC est supervisé par la Commission bancaire de l'Afrique centrale, la COBAC.
Cinq textes concernent son système d'information.

\begin{tableaucours}{L{5.2cm}L{6.8cm}}{%
    \tableauHeader{Texte} &
    \tableauHeader{Exigence pour le système d'information}
}
Règlement n°~04/08/CEMAC/UMAC/COBAC du 6~octobre~2008~\cite{cemac200804} & Gouvernement d'entreprise~: rôle du conseil d'administration \\
Règlement COBAC R-2008/01~\cite{cobac200801} & Plan de continuité des activités \\
Règlement COBAC R-2016/04~\cite{cobac201604} & Contrôle interne, y compris les activités externalisées \\
Règlement n°~04/18/CEMAC/UMAC/COBAC du 21~décembre~2018~\cite{cemac2018paiement} & Services de paiement, en vigueur depuis le 1\textsuperscript{er}~janvier~2019 \\
Lettre circulaire LC-COB/04 du 21~janvier~2022~\cite{cobac2022lc} & Audit indépendant de la sécurité, politique de sécurité selon les normes internationales \\
\end{tableaucours}
\vspace{-0.8em}
\begin{center}{\small\itshape Tableau~7.2. Textes communautaires pour le système d'information d'une banque}\end{center}

La lettre circulaire de 2022 demandait aux établissements de faire auditer leur système d'information par des
experts indépendants avant le 30~juin~2022 et de se doter d'une politique de sécurité fondée sur les normes
internationales. Le rapport annuel 2024 de la COBAC signale l'adoption d'un règlement relatif à la gestion du
risque informatique~\cite{cobac2024rapport}. Pour la BANQUE SANAGA, le schéma directeur n'est donc pas un choix.
Il répond à un superviseur qui lit le document et vérifie son exécution.

Le règlement n°~03/16/CEMAC/UMAC/CM du 21~décembre~2016 sur les systèmes, moyens et incidents de
paiement~\cite{cemac2016paiement} et le règlement de 2018 sur les services de paiement s'appliquent aussi aux
opérateurs de paiement par téléphone. L'ANPP les rencontre si elle paie des prestations sur compte mobile.

% ============================================================================
\section{Le droit des affaires}
\label{sec:c07-ohada}
% ============================================================================

LiZbiZ est une société anonyme régie par les Actes uniformes de l'OHADA. L'Acte uniforme relatif au droit
comptable et à l'information financière, adopté à Brazzaville le 26~janvier~2017~\cite{ohada2017audcif},
intéresse le schéma directeur à deux titres. Il fixe les conditions dans lesquelles une dépense informatique
s'inscrit à l'actif ou en charge, ce qui commande la présentation du budget. Il impose de conserver dix ans les
livres comptables et les pièces justificatives, ce qui commande l'archivage et la réversibilité des systèmes
comptables.

% ============================================================================
\section{La carte d'applicabilité}
\label{sec:c07-carte}
% ============================================================================

\begin{tableaucours}{L{4.4cm}ccc}{%
    \tableauHeader{Texte} &
    \tableauHeader{ANPP} &
    \tableauHeader{LiZbiZ} &
    \tableauHeader{BANQUE SANAGA}
}
Loi n°~2010/012, sécurité & Oui & Oui & Oui \\
Loi n°~2024/017, données personnelles & Oui & Oui & Oui \\
Loi n°~2010/021, commerce électronique & Non & Oui & Oui \\
Loi n°~2018/012, régime financier & Oui & Non & Non \\
Loi n°~2017/010, établissements publics & Oui & Non & Non \\
Décret n°~2018/366, marchés publics & Oui & Non & Non \\
Actes uniformes de l'OHADA & Non & Oui & Oui \\
Règlements de la COBAC & Non & Non & Oui \\
Règlements de paiement de la CEMAC & Si paiement mobile & Non & Oui \\
\end{tableaucours}
\vspace{-0.8em}
\begin{center}{\small\itshape Tableau~7.3. Carte d'applicabilité des textes aux trois organisations}\end{center}

\begin{boxprocedure}
\textbf{Dresser la carte d'applicabilité d'une organisation.}
\begin{enumerate}
    \item Qualifier l'organisation~: administration, établissement public, entreprise publique, société privée,
          établissement supervisé.
    \item Lister les textes de chaque famille du Tableau~7.1.
    \item Pour chacun, répondre oui, non ou sous condition et écrire la condition.
    \item Pour chaque oui, noter l'exigence qui touche le schéma directeur et la pièce qui la prouvera.
    \item Dater la carte. Les textes changent.
\end{enumerate}
\end{boxprocedure}

\begin{boxlecture}{Le cadre juridique dans le schéma directeur 2019 de l'ANPP}
\begin{tableaucours}{L{7.6cm}cL{3.4cm}}{%
    \tableauHeader{Exigence} &
    \tableauHeader{Cote} &
    \tableauHeader{Preuve}
}
Le document est approuvé par le conseil & \OUI{} & Délibération du 17 décembre 2018 \\
Une délibération en tire les conséquences budgétaires & \NON{} & Aucune \\
Les projets sont chiffrés en engagements et en paiements & \NON{} & Total unique \\
Chaque projet est rattaché à un programme budgétaire & \NON{} & Aucun rattachement \\
Les délais de passation des marchés sont comptés & \NON{} & Calendrier sans procédure \\
Aucun produit n'est nommé & \NON{} & Onze marques citées \\
L'audit de sécurité obligatoire est prévu & \NON{} & Absent \\
Les traitements de données personnelles sont recensés & \NEK{} & Loi postérieure au document \\
\end{tableaucours}

\medskip
\textbf{Décompte.} Un \OUI{}, six \NON{}, un \NEK{} sur huit.
\hfill\lecture{0{,}125}{0{,}125}{1{,}06}

\medskip
\textbf{Lecture.} Un \OUI{} sur huit. Le document de 2019 était approuvé et inapplicable. Le \NEK{} est de critère~: la
loi sur les données personnelles date de 2024 et ne pouvait pas être citée en 2018. Les six \NON{} tiennent à
l'ignorance du droit public par ses rédacteurs. Aucun n'est technique.
\end{boxlecture}

\begin{boxmetier}{Directeur financier}
Fournir au consultant le budget voté, le cadre de dépenses à moyen terme et la programmation des marchés. Sans ces trois pièces, la trajectoire est une fiction.
\end{boxmetier}

\begin{boxmetier}{Consultant}
Dresser la carte d'applicabilité avant le diagnostic. Elle dit quelles pièces demander et à qui.
\end{boxmetier}

\begin{boxmetier}{Directeur des systèmes d'information}
Écrire les besoins par leurs fonctions. Tenir un calendrier des marchés à côté du calendrier des projets.
\end{boxmetier}

\begin{boxresume}
\begin{itemize}
    \item Trois lois du 21~décembre~2010~: cybersécurité, communications électroniques, commerce électronique.
    \item La loi n°~2024/017 ajoute à chaque projet une tâche de conformité.
    \item La loi n°~2018/012 distingue autorisations d'engagement et crédits de paiement.
    \item Le conseil d'administration approuve le schéma directeur (lois n°~2017/010 et 2017/011).
    \item Le code des marchés publics (décret n°~2018/366) impose des délais et interdit de nommer un produit.
    \item Une banque répond à la COBAC~: continuité, contrôle interne, audit de sécurité.
    \item L'Acte uniforme comptable de l'OHADA impose dix ans de conservation.
\end{itemize}
\end{boxresume}

\begin{boxinfo}
\textbf{Points de contrôle.} Vous avez compris ce chapitre si vous savez~:
\begin{itemize}
    \item dire pourquoi un schéma directeur ne nomme pas de produit~;
    \item expliquer la différence entre engagement et paiement~;
    \item citer trois textes de la COBAC ou de la CEMAC~;
    \item dresser la carte d'applicabilité d'une organisation.
\end{itemize}
\end{boxinfo}

\begin{boxexo}[title={\ding{45}\quad Exercice~7.1. Quelle famille \niveaubase}]
Rangez dans les cinq familles du Tableau~7.1~: le décret n°~2018/366, la loi n°~2018/012, la loi n°~2024/017, le règlement COBAC R-2008/01, la loi n°~2017/010, la loi n°~2010/021, l'Acte uniforme comptable.
\end{boxexo}

\begin{boxexo}[title={\ding{45}\quad Exercice~7.2. Engagement et paiement \niveaumoyen}]
Un projet de l'ANPP coûte 900~millions sur trois ans~: 200, 450 puis 250. Dressez le tableau des autorisations d'engagement et des crédits de paiement. Que se passe-t-il si les crédits de la deuxième année sont réduits à 300~?
\end{boxexo}

\begin{boxexo}[title={\ding{45}\quad Exercice~7.3. Carte d'une université \niveaumoyen}]
Dressez la carte d'applicabilité d'une université d'État qui veut un schéma directeur numérique. Notez pour trois textes l'exigence et la pièce.
\end{boxexo}

\begin{boxexo}[title={\ding{45}\quad Exercice~7.4. La banque sous-traite \niveauexpert}]
La BANQUE SANAGA veut confier son hébergement à un prestataire étranger. Listez les textes concernés dans chacune des cinq familles et les questions que le superviseur posera. Séparez ce que les textes exigent de ce que la prudence recommande.
\end{boxexo}

\bigskip

\begin{boxinfo}
\textbf{Pour aller plus loin.}
\begin{itemize}
    \item La loi n°~2018/012~\cite{loi2018012}.
    \item Le code des marchés publics~\cite{decret2018366}.
    \item Les lois n°~2017/010 et 2017/011~\cite{loi2017010,loi2017011}.
    \item La lettre circulaire LC-COB/04~\cite{cobac2022lc}.
\end{itemize}
\end{boxinfo}

\begin{boxinfo}
\textbf{Transition.} Le contexte est posé~: théories, normes, orientations, droit. La méthode peut commencer. Son premier temps est le cadrage, objet du Chapitre~\ref{chap:cadrer}.
\end{boxinfo}


    \part{Méthode : élaborer le schéma directeur}
    \label{part:methode}

        % ============================================================================
% Fichier  : partie4/P04_C08_cadrer.tex
% Titre    : Cadrer et Gouverner la Démarche
% Auteur   : MINKA MI NGUIDJOÏ Thierry Emmanuel
% Version  : 1.0 -- novembre 2026
% Corrigés : corriges/Corr_C08.tex
% ============================================================================

\chapter{Cadrer et Gouverner la Démarche}
\label{chap:cadrer}
\chaptermark{Cadrer}

\epigraph{%
    « Ce qui n'est pas écrit au départ\\
    se discute à l'arrivée. »%
}{Principe du cours}

\niveauChapitre{Fondamental}{%
    Chapitres~\ref{chap:alignement} et \ref{chap:droit}
}

\begin{boxobjectifs}
\begin{itemize}
    \item Répondre par écrit aux six questions du cadrage.
    \item Définir un périmètre par trois listes et une liste d'exclusions.
    \item Composer les instances d'une démarche de schéma directeur.
    \item Placer les parties prenantes sur une grille pouvoir-intérêt.
    \item Dresser une matrice des responsabilités.
    \item Rédiger une note de cadrage.
\end{itemize}
\end{boxobjectifs}

\bigskip

% ============================================================================
\section*{Avant de commencer : « qui préside ? »}
% ============================================================================

Yaoundé, jeudi 14~janvier~2027. Réunion de lancement dans la salle du conseil de l'ANPP. M.~BEYALA projette la
composition du comité de pilotage. La première ligne porte le nom de Mme~NGO BIYONG comme présidente. C'est ce
que prévoyait le cahier des charges.

Mme~NGO BIYONG lève la main la première. « Si je préside, ce sera encore le plan de l'informatique. En 2021
c'était le cas. Vous savez ce qu'il est devenu. » M.~OWONA, le directeur général, hésite. Son agenda est plein.
M.~BEYALA lui pose une seule question~: « Qui, dans cette salle, peut retirer un milliard à une direction pour le
donner à une autre~? » Le directeur général préside.

\begin{boxalert}
\textbf{Violation typique.} Commencer le diagnostic avant d'avoir une note de cadrage signée.

\smallskip
\textbf{Règle~:} le cadrage se termine par une signature. Sans elle, chaque désaccord ultérieur rouvre le
périmètre, l'horizon ou la présidence.
\end{boxalert}

% ============================================================================
\section{Ce que Cadrer Veut Dire}
\label{sec:c08-cadrer}
% ============================================================================

Cadrer, c'est répondre par écrit à six questions avant de commencer. Tant qu'une réponse manque, le diagnostic
ne démarre pas.

\begin{tableaucours}{L{3.0cm}L{4.4cm}L{4.6cm}}{%
    \tableauHeader{Question} &
    \tableauHeader{Ce qu'il faut écrire} &
    \tableauHeader{Réponse de l'ANPP}
}
Pourquoi & Le fait qui déclenche la démarche & Document de 2021 échu, panne de mars 2026 \\
Pour qui & Le commanditaire et l'organe qui approuve & Directeur général, conseil d'administration \\
Sur quoi & Le périmètre~: entités, métiers, systèmes & Siège et dix antennes, 47 applications \\
Jusqu'où & L'horizon et le niveau d'ambition & 2027-2031, enveloppe à fixer \\
Avec qui & Les instances et leurs membres & Trois instances, Tableau~8.2 \\
Pour quand & Le calendrier et les livrables & 12 janvier au 24 juin 2027 \\
\end{tableaucours}
\vspace{-0.8em}
\begin{center}{\small\itshape Tableau~8.1. Les six questions du cadrage}\end{center}

Le produit du cadrage est une \termecle{note de cadrage}. Elle tient en dix pages, elle est signée par le
commanditaire et elle sert de référence pendant toute la mission. L'Annexe~B en donne un modèle.

% ============================================================================
\section{Le Périmètre et l'Horizon}
\label{sec:c08-perimetre}
% ============================================================================

Le périmètre se définit par trois listes. La première nomme les entités~: le siège, les antennes, les filiales.
La deuxième nomme les métiers~: à l'ANPP, l'immatriculation des bénéficiaires, la liquidation des droits, le
paiement, le recouvrement, les fonctions de soutien. La troisième nomme les objets~: les applications, les
infrastructures, les données, l'organisation informatique, les contrats.

Ce qui n'est pas dans le périmètre s'écrit aussi. L'ANPP a exclu la téléphonie des antennes, gérée par un
contrat distinct qui court jusqu'en 2029. L'exclusion est datée et motivée.

L'horizon dépend du rythme de l'organisation. Trois ans conviennent à une entreprise dont le marché bouge.
Cinq ans conviennent à une administration dont le plan stratégique court sur cinq ans. Au-delà, le document ne
prévoit plus, il rêve. L'ANPP retient 2027-2031 pour couvrir son plan stratégique 2026-2030 et l'exercice qui
suit.

\begin{boxinfo}
\textbf{Le niveau d'ambition.} Le commanditaire choisit entre trois ambitions avant le diagnostic. Remettre à
niveau, c'est rattraper un retard. Aligner, c'est servir le plan stratégique de l'organisation. Transformer,
c'est changer la manière de rendre le service. Le directeur général de l'ANPP a choisi d'aligner. Le choix
borne la mission~: aucun scénario de transformation ne sera chiffré.
\end{boxinfo}

% ============================================================================
\section{Les Instances}
\label{sec:c08-instances}
% ============================================================================

La norme ISO/IEC~38500 demande à l'organe de gouvernance d'évaluer, de diriger et de surveiller~\cite{iso38500}.
Appliquée à la démarche, elle interdit de confier le schéma directeur à la seule direction informatique.
Trois instances suffisent.

\begin{tableaucours}{L{2.8cm}L{4.0cm}L{2.4cm}L{2.6cm}}{%
    \tableauHeader{Instance} &
    \tableauHeader{Composition à l'ANPP} &
    \tableauHeader{Rythme} &
    \tableauHeader{Rôle}
}
Comité de pilotage & Directeur général, les quatre directeurs, cheffe de la division informatique & Fin de chaque phase & Décide \\
Équipe projet & Mme~NGO BIYONG, deux agents de la division, le cabinet & Chaque semaine & Produit \\
Groupes de travail & Un par métier, cinq à sept personnes dont un chef d'antenne & Selon la phase & Informe et valide le détail \\
\end{tableaucours}
\vspace{-0.8em}
\begin{center}{\small\itshape Tableau~8.2. Les trois instances de la mission}\end{center}

Le conseil d'administration n'est pas une instance de la démarche. Il intervient deux fois~: il est informé du
lancement et il approuve le document. M.~TCHUENTE, son président, a demandé un point d'étape en avril.

\begin{boxalert}
\textbf{Violation typique.} Un comité de pilotage présidé par le responsable informatique.

\smallskip
\textbf{Règle~:} le comité de pilotage est présidé par celui qui arbitre les budgets. Si le directeur général
délègue, il délègue à un directeur de métier, jamais à l'informatique, qui serait juge et partie.
\end{boxalert}

% ============================================================================
\section{Les Parties Prenantes}
\label{sec:c08-parties}
% ============================================================================

Une partie prenante est une personne ou un groupe qui peut peser sur la démarche ou que la démarche touche. On
les recense, puis on les place sur une grille à deux axes~: leur pouvoir sur la décision, leur intérêt pour le
sujet.

\begin{tableaucours}{L{3.2cm}L{4.2cm}L{4.6cm}}{%
    \tableauHeader{Position} &
    \tableauHeader{Exemple à l'ANPP} &
    \tableauHeader{Conduite}
}
Pouvoir fort, intérêt fort & Directeur général, directeur financier & Associer à chaque décision \\
Pouvoir fort, intérêt faible & Tutelle, conseil d'administration & Informer brièvement, à dates fixes \\
Pouvoir faible, intérêt fort & Chefs d'antenne, agents de guichet & Consulter, restituer \\
Pouvoir faible, intérêt faible & Fournisseurs sortants & Surveiller \\
\end{tableaucours}
\vspace{-0.8em}
\begin{center}{\small\itshape Tableau~8.3. La grille des parties prenantes}\end{center}

Deux parties prenantes sont souvent oubliées. Les bénéficiaires ne siègent nulle part, pourtant le service
leur est destiné. L'ANPP a prévu deux réunions avec leurs associations. Les agents de la division informatique
craignent pour leur poste. Mme~NGO BIYONG les a réunis avant le lancement.

% ============================================================================
\section{Qui Fait Quoi}
\label{sec:c08-raci}
% ============================================================================

Une matrice des responsabilités attribue à chaque livrable quatre rôles. Le réalisateur produit. L'approbateur
décide et il est seul. Le consulté donne un avis avant. L'informé reçoit le résultat.

\begin{tableaucours}{L{3.6cm}cccc}{%
    \tableauHeader{Livrable} &
    \tableauHeader{Direction générale} &
    \tableauHeader{Division informatique} &
    \tableauHeader{Directions de métier} &
    \tableauHeader{Cabinet}
}
Note de cadrage & A & C & C & R \\
Rapport de diagnostic & I & C & C & R \\
Cible & A & C & C & R \\
Portefeuille de projets & A & R & C & R \\
Plan de financement & A & C & I & R \\
Schéma directeur & A & C & C & R \\
\end{tableaucours}
\vspace{-0.8em}
\begin{center}{\small\itshape Tableau~8.4. Matrice des responsabilités (R réalise, A approuve, C consulté, I informé)}\end{center}

Le directeur financier, M.~ABDOULAYE, figure parmi les directions de métier. Pour le plan de financement, la
matrice de l'ANPP le fait passer de consulté à réalisateur associé. Le cabinet ne connaît pas le budget de
l'Agence mieux que lui.

% ============================================================================
\section{Le Calendrier et les Livrables}
\label{sec:c08-calendrier}
% ============================================================================

\begin{tableaucours}{L{0.6cm}L{3.2cm}L{3.6cm}L{4.4cm}}{%
    \tableauHeader{N°} &
    \tableauHeader{Phase} &
    \tableauHeader{Dates en 2027} &
    \tableauHeader{Livrable}
}
1 & Cadrage & 12 au 29 janvier & Note de cadrage \\
2 & Diagnostic & 1\textsuperscript{er} février au 19 mars & Rapport de diagnostic \\
3 & Cible & 22 mars au 23 avril & Dossier de cible, scénarios \\
4 & Portefeuille et trajectoire & 26 avril au 21 mai & Portefeuille, feuille de route \\
5 & Chiffrage et financement & 24 mai au 11 juin & Plan de financement \\
6 & Validation & 14 au 24 juin & Schéma directeur 2027-2031 \\
\end{tableaucours}
\vspace{-0.8em}
\begin{center}{\small\itshape Tableau~8.5. Les six phases de la mission ANPP}\end{center}

Chaque phase se clôt par un comité de pilotage qui approuve le livrable. Une phase non approuvée n'est pas
terminée. La règle a un coût~: le calendrier dépend de l'agenda du directeur général. Elle a un bénéfice~: le
document final ne contient que des décisions déjà prises.

Le schéma directeur de 2021 comptait 214 pages. La note de cadrage de 2027 fixe un plafond de 60 pages hors
annexes, plus une synthèse de quatre pages pour le conseil.

\begin{boxprocedure}
\textbf{Rédiger une note de cadrage.}
\begin{enumerate}
    \item Écrire le fait déclencheur en trois lignes.
    \item Nommer le commanditaire et l'organe qui approuvera.
    \item Dresser les trois listes du périmètre et la liste des exclusions.
    \item Fixer l'horizon et le niveau d'ambition.
    \item Composer les instances, avec les noms.
    \item Dresser la grille des parties prenantes et la matrice des responsabilités.
    \item Poser le calendrier, les livrables et le plafond de pages.
    \item Faire signer.
\end{enumerate}
\end{boxprocedure}

% ============================================================================
\section{Communiquer}
\label{sec:c08-communiquer}
% ============================================================================

Un schéma directeur que personne ne connaît ne s'exécute pas. Le plan de communication tient en un tableau~: à
qui, quoi, quand, par quel moyen. À l'ANPP, une note du directeur général annonce la mission à tous les agents
le 15~janvier. Une page d'information paraît à la fin de chaque phase. Les chefs d'antenne reçoivent le
diagnostic avant le comité de pilotage, non après.

\begin{boxlecture}{La note de cadrage de la mission ANPP, 29 janvier 2027}
\begin{tableaucours}{L{7.6cm}cL{3.4cm}}{%
    \tableauHeader{Exigence} &
    \tableauHeader{Cote} &
    \tableauHeader{Preuve}
}
Le fait déclencheur est écrit & \OUI{} & Section 1 de la note \\
Le commanditaire et l'organe d'approbation sont nommés & \OUI{} & Section 2 \\
Le périmètre et les exclusions sont listés & \OUI{} & Section 3, téléphonie exclue \\
L'horizon et l'ambition sont fixés & \OUI{} & 2027-2031, aligner \\
Le comité de pilotage est présidé par un arbitre des budgets & \OUI{} & Directeur général \\
Les bénéficiaires sont représentés ou consultés & \OUI{} & Deux réunions prévues \\
Chaque livrable a un approbateur unique & \OUI{} & Tableau 8.4 \\
L'enveloppe financière est fixée & \NON{} & Renvoyée à la phase 5 \\
La note est signée par le commanditaire & \OUI{} & Signature du 29 janvier \\
La tutelle a été informée du lancement & \NEK{} & Courrier annoncé, copie non remise \\
\end{tableaucours}

\medskip
\textbf{Décompte.} Huit \OUI{}, un \NON{}, un \NEK{} sur dix.
\hfill\lecture{0{,}80}{0{,}10}{0{,}92}

\medskip
\textbf{Lecture.} Huit \OUI{} sur dix. Le \NON{} est un choix assumé~: le directeur général ne veut pas fixer d'enveloppe
avant le diagnostic. Le risque est connu. Une cible dessinée sans plafond coûte toujours plus que ce que
l'organisation peut payer. Le \NEK{} est de preuve. Il se lève en obtenant la copie du courrier.
\end{boxlecture}

\begin{boxmetier}{Directeur général}
Présider le comité de pilotage. Signer la note de cadrage. Dire avant le diagnostic quelle ambition il retient.
\end{boxmetier}

\begin{boxmetier}{Consultant}
Refuser de commencer sans note signée. Tenir le plafond de pages. Préparer chaque comité par un document de décision d'une page.
\end{boxmetier}

\begin{boxmetier}{Directeur des systèmes d'information}
Animer l'équipe projet sans présider. Rassurer ses agents. Ouvrir ses contrats et ses comptes au cabinet.
\end{boxmetier}

\begin{boxresume}
\begin{itemize}
    \item Cadrer, c'est répondre par écrit à six questions~: pourquoi, pour qui, sur quoi, jusqu'où, avec qui, pour quand.
    \item Le périmètre tient en trois listes et une liste d'exclusions datées.
    \item Trois instances~: comité de pilotage, équipe projet, groupes de travail.
    \item Le comité de pilotage est présidé par celui qui arbitre les budgets.
    \item Chaque livrable a un seul approbateur.
    \item Une phase non approuvée n'est pas terminée.
    \item La note de cadrage est signée avant le diagnostic.
\end{itemize}
\end{boxresume}

\begin{boxinfo}
\textbf{Points de contrôle.} Vous avez compris ce chapitre si vous savez~:
\begin{itemize}
    \item rédiger une note de cadrage de dix pages~;
    \item composer un comité de pilotage et justifier sa présidence~;
    \item placer dix parties prenantes sur la grille~;
    \item dresser la matrice des responsabilités d'une mission.
\end{itemize}
\end{boxinfo}

\begin{boxexo}[title={\ding{45}\quad Exercice~8.1. Les six questions \niveaubase}]
Répondez aux six questions du cadrage pour LiZbiZ, à partir de ce que les chapitres précédents en disent. Signalez les réponses que vous ne pouvez pas donner.
\end{boxexo}

\begin{boxexo}[title={\ding{45}\quad Exercice~8.2. La présidence \niveaumoyen}]
À la BANQUE SANAGA, la directrice générale propose que le directeur des risques préside le comité de pilotage. Donnez deux arguments pour et deux arguments contre. Concluez.
\end{boxexo}

\begin{boxexo}[title={\ding{45}\quad Exercice~8.3. La grille \niveaumoyen}]
Placez sur la grille du Tableau~8.3 les parties prenantes suivantes de l'ANPP~: le PCA, le syndicat des agents, l'éditeur d'e-RECETTES, l'ANTIC, les associations de bénéficiaires, le contrôleur financier. Justifiez chaque position en une ligne.
\end{boxexo}

\begin{boxexo}[title={\ding{45}\quad Exercice~8.4. Le cadrage sans enveloppe \niveauexpert}]
Le directeur général de l'ANPP refuse de fixer une enveloppe avant le diagnostic. Rédigez la note d'une page que le chef de mission lui adresse pour proposer une fourchette, avec la méthode pour l'établir à partir des cinq exercices passés.
\end{boxexo}

\bigskip

\begin{boxinfo}
\textbf{Pour aller plus loin.}
\begin{itemize}
    \item ISO/IEC~38500~\cite{iso38500}.
    \item ISO~21502 sur le management de projet~\cite{iso21502}.
    \item Peppard et Ward, chapitre sur le processus de planification~\cite{wardpeppard2016}.
\end{itemize}
\end{boxinfo}

\begin{boxinfo}
\textbf{Transition.} La note est signée. Le cabinet peut regarder ce qui existe. C'est l'objet du Chapitre~\ref{chap:diagnostic}.
\end{boxinfo}

        % ============================================================================
% Fichier  : partie4/P04_C09_diagnostiquer.tex
% Titre    : Diagnostiquer l'existant
% Auteur   : MINKA MI NGUIDJOÏ Thierry Emmanuel
% Version  : 1.0 -- novembre 2026
% Corrigés : corriges/Corr_C09.tex
% ============================================================================

\chapter{Diagnostiquer l'existant}
\label{chap:diagnostic}
\chaptermark{Diagnostiquer}

\epigraph{%
    « On ne gère bien que ce que l'on a d'abord compté. »%
}{Principe du cours}

\niveauChapitre{Intermédiaire}{%
    Chapitre~\ref{chap:cadrer}
}

\begin{boxobjectifs}
\begin{itemize}
    \item Porter cinq regards sur l'existant.
    \item Conduire un entretien de diagnostic.
    \item Tenir un inventaire et une fiche d'application.
    \item Noter une application sur sa valeur et sa santé.
    \item Reconstituer la dépense informatique et la partager entre fonctionnement et projets.
    \item Rédiger un constat en trois temps.
\end{itemize}
\end{boxobjectifs}

\bigskip

% ============================================================================
\section*{Avant de commencer : « nous avons trente et une applications »}
% ============================================================================

Yaoundé, lundi 10~février~2025. Mme~NGO BIYONG remet au cabinet la liste des applications de l'Agence. Elle en
compte 31. Le même jour, Mme~EKANI revient de la direction des prestations avec une clé contenant un classeur
de calcul. Un chef de service l'a construit en 2017 pour corriger les pensions de réversion que l'application
officielle calcule mal. Il traite 4~200 dossiers par an. Personne à la division informatique ne le connaît.

En trois semaines, le cabinet trouve seize outils de cette sorte. Le consultant junior propose de les appeler
« applications pirates ». M.~BEYALA corrige~: « Ce sont des besoins que personne n'a écoutés. Vous les
inventoriez comme les autres. »

\begin{boxalert}
\textbf{Violation typique.} Établir l'inventaire à partir de la seule liste de la fonction informatique.

\smallskip
\textbf{Règle~:} l'inventaire se construit depuis les métiers. On demande à chaque service avec quoi il
travaille, puis on rapproche de la liste officielle.
\end{boxalert}

% ============================================================================
\section{Cinq regards sur l'existant}
\label{sec:c09-regards}
% ============================================================================

Le diagnostic décrit l'existant sans le juger d'avance. Il porte cinq regards. Chacun a ses sources et son
livrable.

\begin{tableaucours}{L{2.6cm}L{4.6cm}L{4.8cm}}{%
    \tableauHeader{Regard} &
    \tableauHeader{Question} &
    \tableauHeader{Sources}
}
Métier & Que fait l'organisation, que veut-elle faire~? & Plan stratégique, entretiens de direction \\
Applications & Quels logiciels, pour quels métiers, dans quel état~? & Inventaire, contrats, entretiens d'utilisateurs \\
Infrastructures & Sur quoi tout cela tourne-t-il~? & Inventaire technique, visites de salles \\
Organisation & Qui s'en occupe, avec quelles compétences~? & Organigramme, fiches de poste \\
Finances & Combien cela coûte-t-il, pour quoi~? & Comptes des cinq derniers exercices, marchés \\
\end{tableaucours}
\vspace{-0.8em}
\begin{center}{\small\itshape Tableau~9.1. Les cinq regards du diagnostic}\end{center}

Le diagnostic ne propose rien. Les idées de solution qui viennent pendant les entretiens se notent dans un
carnet à part. Elles serviront à la phase de cible.

% ============================================================================
\section{Collecter}
\label{sec:c09-collecter}
% ============================================================================

Trois moyens se complètent. Les documents disent ce qui a été décidé. Les entretiens disent ce qui se fait.
L'observation dit ce qui se passe.

À l'ANPP, le cabinet a demandé 34~documents le premier jour. Il en a reçu 21 en deux semaines. Il a conduit
38~entretiens d'une heure~: les directeurs, les chefs de service, six chefs d'antenne, les agents de la division
informatique. Il a passé une journée au guichet de l'antenne de Bafoussam et une journée à celle de Garoua.

\begin{boxprocedure}
\textbf{Conduire un entretien de diagnostic.}
\begin{enumerate}
    \item Envoyer trois questions la veille. Pas davantage.
    \item Commencer par le travail de la personne, non par l'informatique.
    \item Demander ce qui se passe quand le système est arrêté.
    \item Demander à voir~: un écran, un fichier, un registre.
    \item Noter les faits, les chiffres et leur source. Séparer les opinions.
    \item Envoyer le compte rendu sous trois jours et le faire valider.
\end{enumerate}
\end{boxprocedure}

Un document manquant est une information. Treize des 34~documents demandés n'existent pas à l'ANPP, dont
l'inventaire des applications, le plan de secours et la liste des habilitations. Le rapport de diagnostic les
énumère.

% ============================================================================
\section{L'inventaire des applications}
\label{sec:c09-inventaire}
% ============================================================================

L'inventaire est une liste où chaque application occupe une ligne. La \termecle{fiche d'application} en donne
le détail. Elle répond à dix questions~: nom, métier servi, responsable côté métier, nombre d'utilisateurs,
éditeur ou auteur, année de mise en service, technologie, contrat de maintenance, coût annuel, données
personnelles traitées.

La division informatique de l'ANPP croyait gérer 31~applications. Le cabinet en a trouvé 47. Les seize autres
sont des classeurs, des bases bureautiques ou des services en ligne souscrits par une direction avec sa
carte. Trois d'entre elles servent à calculer des droits.

\begin{boxdef}
\textbf{Informatique hors contrôle.} Ensemble des outils numériques employés dans une organisation sans que la
fonction informatique les ait fournis ni recensés. Elle signale presque toujours un besoin que l'offre
officielle ne couvre pas.
\end{boxdef}

% ============================================================================
\section{Deux notes par application}
\label{sec:c09-notes}
% ============================================================================

Chaque application reçoit deux notes de 1 à 5. La \termecle{valeur pour le métier} mesure le service rendu. La
\termecle{santé technique} mesure la capacité à durer. Chaque note se calcule sur quatre critères notés de 1 à 5
dont on prend la moyenne.

\begin{tableaucours}{L{5.6cm}L{6.4cm}}{%
    \tableauHeader{Valeur pour le métier} &
    \tableauHeader{Santé technique}
}
Contribution aux objectifs du plan stratégique & Technologie encore maintenue par son éditeur \\
Part de l'activité qui en dépend & Compétences disponibles pour la faire évoluer \\
Satisfaction des utilisateurs & Nombre d'incidents sur douze mois \\
Absence d'alternative manuelle & Documentation et sauvegardes vérifiées \\
\end{tableaucours}
\vspace{-0.8em}
\begin{center}{\small\itshape Tableau~9.2. Les huit critères de notation d'une application}\end{center}

Le seuil est fixé à 3. Le croisement des deux notes donne quatre décisions.

\begin{tableaucours}{L{3.0cm}L{3.0cm}L{2.4cm}L{3.6cm}}{%
    \tableauHeader{Valeur} &
    \tableauHeader{Santé} &
    \tableauHeader{Décision} &
    \tableauHeader{Applications de l'ANPP}
}
3 ou plus & 3 ou plus & Investir & 6 \\
3 ou plus & Moins de 3 & Refondre & 14 \\
Moins de 3 & 3 ou plus & Conserver & 18 \\
Moins de 3 & Moins de 3 & Retirer & 9 \\
\end{tableaucours}
\vspace{-0.8em}
\begin{center}{\small\itshape Tableau~9.3. Les quatre décisions et le résultat pour les 47 applications}\end{center}

Quatorze applications à refondre~: c'est le cœur du prochain schéma directeur. Elles comprennent la liquidation
des droits, écrite en 2004 dans un langage que deux agents seulement savent lire. Neuf applications à retirer~:
c'est une économie immédiate si l'on résilie leurs contrats.

\begin{boxalert}
\textbf{Violation typique.} Noter les applications en réunion, à main levée.

\smallskip
\textbf{Règle~:} chaque note renvoie à une pièce. Le nombre d'incidents vient du registre, la satisfaction
d'une enquête, la date de fin de maintenance du contrat. Une note sans pièce est une opinion.
\end{boxalert}

% ============================================================================
\section{Les infrastructures et la sécurité}
\label{sec:c09-infra}
% ============================================================================

Le regard technique recense les salles, les serveurs, le réseau, les postes, les sauvegardes. Il vérifie quatre
points qui commandent la continuité~: l'alimentation électrique, la climatisation, le lien entre les sites,
l'existence d'un site de secours.

À l'ANPP, la salle du siège abrite tous les serveurs. Les antennes y accèdent par un lien unique. La panne de
mars~2024 a duré quatre jours parce qu'un onduleur a cédé et que la dernière sauvegarde restaurable datait de
onze jours. Le diagnostic ne refait pas l'audit de sécurité. Il en reprend les constats et note s'il existe.
À l'ANPP, l'audit obligatoire n'a pas été réalisé depuis 2020.

% ============================================================================
\section{L'organisation et les compétences}
\label{sec:c09-organisation}
% ============================================================================

La division informatique de l'ANPP compte 19~agents pour 412~agents de l'Agence. Le référentiel européen des
compétences numériques, la norme EN~16234-1~\cite{ecf}, sert de grille. On y repère les compétences tenues, les
compétences tenues par une seule personne et les compétences absentes.

\begin{tableaucours}{L{4.6cm}cL{5.6cm}}{%
    \tableauHeader{Domaine} &
    \tableauHeader{Agents} &
    \tableauHeader{Constat}
}
Exploitation et assistance & 11 & Effectif suffisant, astreinte non organisée \\
Développement et maintenance & 5 & Deux agents seuls à connaître la liquidation \\
Réseau et sécurité & 2 & Aucun responsable de la sécurité désigné \\
Pilotage, projets, achats & 1 & La cheffe de division, seule \\
Données & 0 & Compétence absente \\
\end{tableaucours}
\vspace{-0.8em}
\begin{center}{\small\itshape Tableau~9.4. La division informatique de l'ANPP par domaine}\end{center}

% ============================================================================
\section{Les finances}
\label{sec:c09-finances}
% ============================================================================

Le regard financier reconstitue la dépense informatique réelle des cinq derniers exercices. Elle ne se lit pas
dans une seule ligne du budget. Il faut y ajouter les achats faits par les autres directions, les contrats
portés par les moyens généraux et la masse salariale de la division.

La dépense se partage en deux. Le \termecle{fonctionnement} maintient l'existant~: salaires, maintenance,
licences, liaisons, énergie. Le \termecle{projet} crée du nouveau. Le rapport entre les deux dit la marge de
manœuvre.

\begin{tableaucours}{L{4.2cm}rr}{%
    \tableauHeader{Nature} &
    \tableauHeader{Montant 2019-2023 (millions de FCFA)} &
    \tableauHeader{Part}
}
Fonctionnement & 2 418 & 78~\% \\
Projets & 682 & 22~\% \\
Total & 3 100 & 100~\% \\
\end{tableaucours}
\vspace{-0.8em}
\begin{center}{\small\itshape Tableau~9.5. La dépense informatique de l'ANPP sur cinq ans}\end{center}

Avec 78~\% de fonctionnement, l'ANPP consacre aux projets 136~millions par an en moyenne. Le schéma directeur de
2019 annonçait 6,3~milliards. L'écart était lisible dès le premier jour pour qui avait fait ce tableau.

% ============================================================================
\section{Restituer}
\label{sec:c09-restituer}
% ============================================================================

Le rapport de diagnostic tient en trente pages. Il s'ouvre par dix constats au plus, rangés par ordre de
gravité. Chaque constat suit la même forme~: le fait, la pièce, la conséquence pour le métier. Une synthèse en
quatre cases, forces, faiblesses, occasions, menaces, peut le clore. Elle résume, elle ne remplace pas les
constats.

\begin{boxscenario}{Mission ANPP. Les cinq premiers constats, 21 mars 2025}
\begin{enumerate}
    \item La liquidation des droits repose sur une application de 2004 que deux agents savent maintenir.
    \item Aucun site de secours. La dernière panne a suspendu les paiements quatre jours.
    \item Seize applications sur 47 échappent à la division informatique. Trois calculent des droits.
    \item Les projets reçoivent 22~\% de la dépense, soit 136~millions par an.
    \item Treize documents de référence n'existent pas, dont l'inventaire et le plan de secours.
\end{enumerate}
\end{boxscenario}

Le comité de pilotage valide le diagnostic. Les directeurs y lisent parfois des faits qui les gênent. Le
consultant ne retire pas un constat prouvé. Il peut en revoir la formulation.

\begin{boxlecture}{Le diagnostic de la mission ANPP, 21 mars 2025}
\begin{tableaucours}{L{7.6cm}cL{3.4cm}}{%
    \tableauHeader{Exigence} &
    \tableauHeader{Cote} &
    \tableauHeader{Preuve}
}
L'inventaire couvre toutes les applications trouvées & \OUI{} & 47 fiches \\
Chaque application a un responsable côté métier & \NON{} & 19 fiches sans responsable \\
Chaque note renvoie à une pièce & \OUI{} & Registre des pièces du diagnostic \\
La dépense de cinq exercices est reconstituée & \OUI{} & Tableau 9.5 \\
La dépense des autres directions est incluse & \NEK{} & Trois directions n'ont pas répondu \\
Les compétences critiques sont repérées & \OUI{} & Tableau 9.4 \\
Les comptes rendus d'entretien sont validés & \OUI{} & 38 validations \\
Le rapport est approuvé par le comité de pilotage & \OUI{} & Procès-verbal du 21 mars \\
\end{tableaucours}

\medskip
\textbf{Décompte.} Six \OUI{}, un \NON{}, un \NEK{} sur huit.
\hfill\lecture{0{,}75}{0{,}125}{1{,}06}

\medskip
\textbf{Lecture.} Six \OUI{} sur huit. Le \NON{} n'est pas une faute du cabinet. Il est un constat~: dix-neuf applications
n'ont pas de propriétaire. Le \NEK{} est de preuve. Tant que trois directions n'ont pas transmis leurs achats,
le total de 3,1~milliards est un plancher.
\end{boxlecture}

\begin{boxmetier}{Consultant}
Demander les documents le premier jour. Noter la source de chaque chiffre. Tenir un registre des pièces.
\end{boxmetier}

\begin{boxmetier}{Directeur des systèmes d'information}
Fournir l'inventaire qu'il a, sans le maquiller. Accepter que le diagnostic en trouve davantage.
\end{boxmetier}

\begin{boxmetier}{Contrôleur de gestion}
Reconstituer la dépense par nature et par direction. Signaler les achats hors budget informatique.
\end{boxmetier}

\begin{boxmetier}{Auditeur}
Reprendre les notes d'un échantillon d'applications et remonter aux pièces.
\end{boxmetier}

\begin{boxresume}
\begin{itemize}
    \item Le diagnostic porte cinq regards~: métier, applications, infrastructures, organisation, finances.
    \item Un document manquant est un constat.
    \item L'inventaire se construit depuis les métiers.
    \item Chaque application reçoit une note de valeur et une note de santé, d'où quatre décisions~: investir, refondre, conserver, retirer.
    \item La dépense se partage entre fonctionnement et projets.
    \item Un constat s'écrit en trois temps~: le fait, la pièce, la conséquence.
    \item Le diagnostic ne propose pas de solution.
\end{itemize}
\end{boxresume}

\begin{boxinfo}
\textbf{Points de contrôle.} Vous avez compris ce chapitre si vous savez~:
\begin{itemize}
    \item conduire un entretien et en rédiger le compte rendu~;
    \item remplir une fiche d'application~;
    \item calculer les deux notes d'une application et en tirer la décision~;
    \item reconstituer la part des projets dans la dépense.
\end{itemize}
\end{boxinfo}

\begin{boxexo}[title={\ding{45}\quad Exercice~9.1. Deux notes \niveaubase}]
Une application de l'ANPP obtient 4, 5, 2, 4 sur les critères de valeur et 1, 2, 3, 2 sur les critères de santé. Calculez les deux notes et donnez la décision.
\end{boxexo}

\begin{boxexo}[title={\ding{45}\quad Exercice~9.2. La dépense \niveaumoyen}]
Sur un exercice, une organisation dépense 310~millions de salaires informatiques, 240 de maintenance et de licences, 90 de liaisons et 160 de projets. Calculez la part des projets. Elle veut la porter à 35~\% sans augmenter le total. De combien doit baisser le fonctionnement~?
\end{boxexo}

\begin{boxexo}[title={\ding{45}\quad Exercice~9.3. Trois constats \niveaumoyen}]
Rédigez en trois temps les constats correspondant aux faits suivants~: l'audit de sécurité n'a pas été fait depuis 2020, la division n'a aucun agent chargé des données, dix-neuf applications n'ont pas de responsable côté métier.
\end{boxexo}

\begin{boxexo}[title={\ding{45}\quad Exercice~9.4. Les seize \niveauexpert}]
Proposez une conduite pour les seize applications hors contrôle de l'ANPP. Distinguez celles qui calculent des droits des autres. Dites ce que le schéma directeur doit prévoir pour que le phénomène ne se reproduise pas, sans l'interdire par principe.
\end{boxexo}

\bigskip

\begin{boxinfo}
\textbf{Pour aller plus loin.}
\begin{itemize}
    \item Peppard et Ward sur l'évaluation du portefeuille existant~\cite{wardpeppard2016}.
    \item COBIT~2019 pour les niveaux de capacité~\cite{cobit2019}.
    \item La norme EN~16234-1~\cite{ecf}.
\end{itemize}
\end{boxinfo}

\begin{boxinfo}
\textbf{Transition.} L'existant est connu. Reste à dire où aller. C'est l'objet du Chapitre~\ref{chap:cible}.
\end{boxinfo}

        % ============================================================================
% Fichier  : partie4/P04_C10_cible.tex
% Titre    : Fixer la cible
% Auteur   : MINKA MI NGUIDJOÏ Thierry Emmanuel
% Version  : 1.0 -- novembre 2026
% Corrigés : corriges/Corr_C10.tex
% ============================================================================

\chapter{Fixer la cible}
\label{chap:cible}
\chaptermark{Fixer la cible}

\epigraph{%
    « Il n'est pas de vent favorable\\
    pour celui qui ne sait où il va. »%
}{Sénèque, \textit{Lettres à Lucilius}, lettre 71}

\niveauChapitre{Intermédiaire}{%
    Chapitres~\ref{chap:normes} et \ref{chap:diagnostic}
}

\begin{boxobjectifs}
\begin{itemize}
    \item Déduire des capacités à partir d'objectifs stratégiques.
    \item Rédiger un principe directeur avec ses conséquences.
    \item Décrire une architecture sur quatre plans.
    \item Rapprocher un plan fonctionnel d'un inventaire.
    \item Conduire une analyse des écarts.
    \item Comparer des scénarios sur une grille pondérée.
\end{itemize}
\end{boxobjectifs}

\bigskip

% ============================================================================
\section*{Avant de commencer : « il nous faut un progiciel intégré »}
% ============================================================================

Yaoundé, lundi 22~mars~2027. Premier jour de la phase de cible. M.~ABDOULAYE arrive avec une plaquette. Un
éditeur lui a présenté un progiciel qui couvre « toute la gestion d'un organisme de protection sociale ». Il
propose d'en faire la cible. « Le diagnostic a montré que tout est à refaire. Refaisons tout d'un coup. »

M.~BEYALA ne dit pas non. Il demande quelle capacité du plan stratégique ce produit apporte, ce qu'il coûte sur
cinq ans et qui, à l'Agence, saurait conduire un projet de cette taille. Le produit devient le scénario~C. Il
sera noté comme les autres.

\begin{boxalert}
\textbf{Violation typique.} Partir d'une solution et chercher ensuite les besoins qu'elle couvre.

\smallskip
\textbf{Règle~:} la cible part des objectifs de l'organisation. Une solution proposée en cours de route
devient un scénario parmi d'autres.
\end{boxalert}

% ============================================================================
\section{De la stratégie à la cible}
\label{sec:c10-chaine}
% ============================================================================

La cible décrit le système d'information voulu à la fin de la période. Elle se déduit, elle ne s'invente pas.
La chaîne compte quatre maillons~: les objectifs de l'organisation, les capacités qu'ils demandent, les
principes qui guideront les choix, l'architecture qui en résulte.

\begin{tableaucours}{L{4.6cm}L{7.4cm}}{%
    \tableauHeader{Objectif du plan stratégique 2026-2030} &
    \tableauHeader{Capacité demandée au système d'information}
}
Payer chaque prestation à date fixe & Liquider et payer sans interruption, y compris en cas de sinistre \\
Ne payer que les ayants droit & Tenir un fichier unique des bénéficiaires, rapproché de l'état civil \\
Rapprocher l'Agence des bénéficiaires & Recevoir et suivre une demande depuis une antenne ou un téléphone \\
Rendre compte de chaque franc & Produire la dépense par programme et par antenne \\
\end{tableaucours}
\vspace{-0.8em}
\begin{center}{\small\itshape Tableau~10.1. Des quatre objectifs de l'ANPP aux capacités}\end{center}

Une capacité se formule par un verbe et un résultat. Elle ne nomme ni un logiciel ni une technique. « Disposer
d'un progiciel intégré » n'est pas une capacité. « Produire la dépense par programme » en est une.

% ============================================================================
\section{Les principes directeurs}
\label{sec:c10-principes}
% ============================================================================

Un \termecle{principe directeur} est une règle durable qui tranche d'avance une famille de décisions. Le cadre
TOGAF demande de l'écrire en quatre parties~: un nom, un énoncé, une justification, des
conséquences~\cite{togaf10}. Les conséquences font la valeur du principe. Un principe qui ne gêne personne ne
décide rien.

\begin{tableaucours}{L{3.2cm}L{4.4cm}L{4.4cm}}{%
    \tableauHeader{Nom} &
    \tableauHeader{Énoncé} &
    \tableauHeader{Conséquence acceptée}
}
Donnée unique & Un bénéficiaire est enregistré une seule fois & Les antennes renoncent à leurs fichiers locaux \\
Acheter avant de construire & Un besoin commun se couvre par un produit du marché & Le métier adapte ses usages au produit \\
Secours d'abord & Aucune application clé sans reprise éprouvée & Un projet sans budget de secours n'est pas lancé \\
Propriétaire désigné & Chaque application a un responsable côté métier & Un directeur signe pour chaque application \\
Ouverture & Les échanges emploient des formats publics & Certains produits sont écartés \\
Sobriété & On retire avant d'ajouter & Tout projet nomme ce qu'il remplace \\
\end{tableaucours}
\vspace{-0.8em}
\begin{center}{\small\itshape Tableau~10.2. Les six principes directeurs de l'ANPP}\end{center}

Six à dix principes suffisent. Le comité de pilotage les approuve un par un, avec leurs conséquences. À l'ANPP,
le principe de la donnée unique a demandé une heure de débat. Les chefs d'antenne tenaient à leurs fichiers.

% ============================================================================
\section{Les quatre plans de l'architecture}
\label{sec:c10-plans}
% ============================================================================

La norme ISO/IEC/IEEE~42010 décrit une architecture par des vues adaptées aux préoccupations de chaque
partie~\cite{iso42010}. La pratique française de l'urbanisation en retient quatre~\cite{longepe2001}. TOGAF en
retient quatre aussi, en regroupant données et applications~\cite{togaf10}.

\begin{tableaucours}{L{2.6cm}L{4.4cm}L{5.0cm}}{%
    \tableauHeader{Plan} &
    \tableauHeader{Contenu} &
    \tableauHeader{Qui le lit}
}
Métier & Processus, acteurs, événements & Les directions de métier \\
Fonctionnel & Blocs de fonctions, indépendants des logiciels & L'équipe projet \\
Applicatif & Logiciels, données, échanges & La fonction informatique \\
Technique & Serveurs, réseaux, sites, hébergement & L'exploitation \\
\end{tableaucours}
\vspace{-0.8em}
\begin{center}{\small\itshape Tableau~10.3. Les quatre plans}\end{center}

Le plan fonctionnel est le plus utile au schéma directeur. Il découpe le système en \termecle{zones}, en
quartiers et en blocs, à la manière d'un plan de ville. Chaque fonction appartient à un seul bloc. Chaque
donnée a un seul bloc maître.

\begin{boxscenario}{Mission ANPP. Le plan fonctionnel cible, 9 avril 2027}
Le plan cible de l'ANPP compte cinq zones.
\begin{enumerate}
    \item \textbf{Référentiels.} Bénéficiaires, employeurs, barèmes. Zone maîtresse des données communes.
    \item \textbf{Production.} Immatriculation, liquidation, paiement, recouvrement.
    \item \textbf{Relation.} Guichet, téléphone, demandes en ligne, réclamations.
    \item \textbf{Soutien.} Comptabilité, budget, ressources humaines, achats.
    \item \textbf{Pilotage.} Tableaux de bord, statistiques, contrôle.
\end{enumerate}
Les 47 applications du diagnostic se placent sur ce plan. Sept blocs sont couverts par deux applications ou
plus. Quatre blocs ne sont couverts par aucune, dont le référentiel des bénéficiaires.
\end{boxscenario}

Un bloc couvert deux fois signale une redondance à réduire. Un bloc vide signale un besoin. Le rapprochement
du plan fonctionnel et de l'inventaire donne la liste brute des chantiers.

% ============================================================================
\section{L'analyse des écarts}
\label{sec:c10-ecarts}
% ============================================================================

Pour chaque capacité, l'analyse des écarts compare l'existant à la cible et qualifie la distance.

\begin{tableaucours}{L{4.0cm}L{3.6cm}L{2.2cm}L{2.2cm}}{%
    \tableauHeader{Capacité} &
    \tableauHeader{Existant} &
    \tableauHeader{Écart} &
    \tableauHeader{Chantier}
}
Liquider et payer sans interruption & Site unique, application de 2004 & Fort & Refonte, secours \\
Fichier unique des bénéficiaires & Onze fichiers & Fort & Référentiel \\
Recevoir une demande à distance & Guichet seul & Moyen & Guichet unique \\
Dépense par programme & Comptabilité générale seule & Moyen & Comptabilité analytique \\
\end{tableaucours}
\vspace{-0.8em}
\begin{center}{\small\itshape Tableau~10.4. L'analyse des écarts de l'ANPP}\end{center}

% ============================================================================
\section{Comparer des scénarios}
\label{sec:c10-scenarios}
% ============================================================================

Une seule cible ne se discute pas, elle s'accepte ou se refuse. Le cabinet en présente donc deux ou trois, qui
atteignent les mêmes capacités par des voies différentes. Chaque scénario est décrit de la même manière~: ce
qu'il fait, ce qu'il coûte en ordre de grandeur, ce qu'il demande à l'organisation, ce qu'il risque.

\begin{tableaucours}{L{2.8cm}L{3.0cm}L{3.0cm}L{3.2cm}}{%
    \tableauHeader{} &
    \tableauHeader{A. Consolider} &
    \tableauHeader{B. Refondre par étapes} &
    \tableauHeader{C. Remplacer d'un bloc}
}
Principe & Sécuriser l'existant, différer la refonte & Secours puis refonte de la production en trois lots & Progiciel intégré unique \\
Ordre de grandeur & 1,9 milliard & 4,2 milliards & 7,5 milliards \\
Capacités atteintes en 2031 & 1 sur 4 & 4 sur 4 & 4 sur 4 \\
Effort interne & Faible & Soutenu & Très fort \\
Risque principal & L'application de 2004 cède & Coexistence de l'ancien et du nouveau & Échec d'un projet unique \\
\end{tableaucours}
\vspace{-0.8em}
\begin{center}{\small\itshape Tableau~10.5. Les trois scénarios présentés au comité de pilotage}\end{center}

Le choix se fait sur une grille de critères pondérés, fixée avant de noter. Le comité de pilotage de l'ANPP a
retenu quatre critères.

\begin{tableaucours}{L{4.4cm}cccc}{%
    \tableauHeader{Critère} &
    \tableauHeader{Poids} &
    \tableauHeader{A} &
    \tableauHeader{B} &
    \tableauHeader{C}
}
Contribution aux quatre objectifs & 40~\% & 2 & 4 & 5 \\
Soutenabilité financière & 25~\% & 5 & 3 & 1 \\
Capacité de l'Agence à conduire & 20~\% & 5 & 3 & 1 \\
Maîtrise du risque & 15~\% & 2 & 4 & 2 \\
Note pondérée &  & 3,35 & 3,55 & 2,75 \\
\end{tableaucours}
\vspace{-0.8em}
\begin{center}{\small\itshape Tableau~10.6. La notation des scénarios (notes de 1 à 5)}\end{center}

Le calcul pour le scénario~B~: $0{,}40 \times 4 + 0{,}25 \times 3 + 0{,}20 \times 3 + 0{,}15 \times 4 = 3{,}55$.
L'écart entre A et B est mince. Le comité a retenu B le 23~avril après avoir constaté que A laissait trois
objectifs du plan stratégique sans réponse. La grille éclaire. Elle ne décide pas à la place du comité.

\begin{boxalert}
\textbf{Violation typique.} Présenter trois scénarios dont deux sont faits pour être écartés.

\smallskip
\textbf{Règle~:} chaque scénario doit pouvoir être choisi par un dirigeant raisonnable. Sinon le cabinet a
décidé seul et le comité n'a fait qu'entériner.
\end{boxalert}

% ============================================================================
\section{Décrire la cible retenue}
\label{sec:c10-decrire}
% ============================================================================

Le dossier de cible tient en quinze pages. Il contient les capacités, les principes, le plan fonctionnel, le
scénario retenu avec ses motifs et la liste des chantiers. Il ne contient pas de calendrier ni de budget
détaillé. Ils viennent aux deux phases suivantes.

La cible porte aussi sur l'organisation. Celle de l'ANPP prévoit trois changements~: un responsable de la
sécurité rattaché au directeur général, une cellule des données de trois agents, un propriétaire désigné pour
chaque application.

\begin{boxlecture}{Le dossier de cible de l'ANPP, 23 avril 2027}
\begin{tableaucours}{L{7.6cm}cL{3.4cm}}{%
    \tableauHeader{Exigence} &
    \tableauHeader{Cote} &
    \tableauHeader{Preuve}
}
Chaque capacité se rattache à un objectif du plan stratégique & \OUI{} & Tableau 10.1 \\
Les principes sont approuvés avec leurs conséquences & \OUI{} & Procès-verbal du 9 avril \\
Le plan fonctionnel situe toutes les applications & \OUI{} & 47 applications placées \\
Chaque donnée commune a un bloc maître & \OUI{} & Zone des référentiels \\
Deux scénarios au moins étaient éligibles & \OUI{} & A et B, écart de 0,20 \\
Les critères et les poids ont été fixés avant la notation & \OUI{} & Procès-verbal du 9 avril \\
L'ordre de grandeur de chaque scénario est justifié & \NON{} & Scénario C estimé par analogie seule \\
La cible d'organisation est décrite & \OUI{} & Section 5 du dossier \\
Aucun produit n'est nommé & \OUI{} & Relecture du dossier \\
\end{tableaucours}

\medskip
\textbf{Décompte.} Huit \OUI{}, un \NON{}, zéro \NEK{} sur neuf.
\hfill\lecture{0{,}889}{0}{0{,}50}

\medskip
\textbf{Lecture.} Huit \OUI{} sur neuf. Le \NON{} porte sur le scénario écarté. Il affaiblit peu la décision, puisque C
était aussi dernier sur la capacité à conduire. Le dossier respecte la règle du Chapitre~\ref{chap:droit}~:
aucune marque n'y figure.
\end{boxlecture}

\begin{boxmetier}{Architecte d'entreprise}
Tenir le plan fonctionnel. Refuser qu'une fonction appartienne à deux blocs.
\end{boxmetier}

\begin{boxmetier}{Directeur de métier}
Formuler des capacités, non des logiciels. Accepter par écrit les conséquences des principes.
\end{boxmetier}

\begin{boxmetier}{Consultant}
Présenter des scénarios tous éligibles. Faire fixer les critères avant de noter.
\end{boxmetier}

\begin{boxresume}
\begin{itemize}
    \item La cible se déduit~: objectifs, capacités, principes, architecture.
    \item Une capacité se formule par un verbe et un résultat.
    \item Un principe directeur vaut par ses conséquences.
    \item Quatre plans~: métier, fonctionnel, applicatif, technique.
    \item Le rapprochement du plan fonctionnel et de l'inventaire révèle redondances et manques.
    \item Deux ou trois scénarios éligibles, notés sur une grille fixée d'avance.
    \item Le dossier de cible ne nomme aucun produit.
\end{itemize}
\end{boxresume}

\begin{boxinfo}
\textbf{Points de contrôle.} Vous avez compris ce chapitre si vous savez~:
\begin{itemize}
    \item traduire un objectif en capacités~;
    \item écrire un principe en quatre parties~;
    \item placer des applications sur un plan fonctionnel~;
    \item calculer la note pondérée d'un scénario.
\end{itemize}
\end{boxinfo}

\begin{boxexo}[title={\ding{45}\quad Exercice~10.1. Capacité ou solution \niveaubase}]
Dites si chaque formule est une capacité ou une solution, puis reformulez les solutions~: « avoir un entrepôt de données », « connaître chaque jour les paiements en attente », « passer au nuage », « payer sur compte mobile », « identifier un bénéficiaire en moins d'une minute ».
\end{boxexo}

\begin{boxexo}[title={\ding{45}\quad Exercice~10.2. Recalculer \niveaumoyen}]
Recalculez les notes pondérées du Tableau~10.6 avec les poids 30~\%, 30~\%, 20~\%, 20~\%. Le classement change-t-il~? Qu'en concluez-vous sur la solidité de la décision~?
\end{boxexo}

\begin{boxexo}[title={\ding{45}\quad Exercice~10.3. Un principe \niveaumoyen}]
Rédigez en quatre parties un principe directeur sur l'hébergement des données de l'ANPP, en tenant compte de la loi n°~2024/017. Écrivez deux conséquences qui gêneront quelqu'un.
\end{boxexo}

\begin{boxexo}[title={\ding{45}\quad Exercice~10.4. La cible de LiZbiZ \niveauexpert}]
LiZbiZ veut répondre à la jeune pousse qui lui a pris 9~\% des livraisons urbaines. Déduisez trois capacités de cette situation, proposez deux scénarios éligibles et les critères pour les départager. Justifiez les poids.
\end{boxexo}

\bigskip

\begin{boxinfo}
\textbf{Pour aller plus loin.}
\begin{itemize}
    \item TOGAF, phases A à D~\cite{togaf10}.
    \item Longépé sur l'urbanisation~\cite{longepe2001}.
    \item Ross, Weill et Robertson sur le modèle opérationnel~\cite{ross2006}.
\end{itemize}
\end{boxinfo}

\begin{boxinfo}
\textbf{Transition.} La cible est choisie. Elle ne dit pas encore par quels projets l'atteindre ni dans quel ordre. C'est l'objet du Chapitre~\ref{chap:portefeuille}.
\end{boxinfo}

        % ============================================================================
% Fichier  : partie4/P04_C11_portefeuille.tex
% Titre    : Portefeuille et trajectoire
% Auteur   : MINKA MI NGUIDJOÏ Thierry Emmanuel
% Version  : 1.0 -- novembre 2026
% Corrigés : corriges/Corr_C11.tex
% ============================================================================

\chapter{Portefeuille et trajectoire}
\label{chap:portefeuille}
\chaptermark{Portefeuille}

\epigraph{%
    « Gouverner, c'est choisir. »%
}{Pierre Mendès France, discours à l'Assemblée nationale, 3~juin~1953}

\niveauChapitre{Intermédiaire}{%
    Chapitre~\ref{chap:cible}
}

\begin{boxobjectifs}
\begin{itemize}
    \item Rédiger une fiche de projet.
    \item Distinguer un projet obligatoire d'un projet arbitrable.
    \item Classer des projets par valeur rapportée au coût et corriger le classement.
    \item Vérifier l'équilibre d'un portefeuille.
    \item Ordonner des projets selon trois contraintes.
    \item Tracer un profil annuel de dépense.
\end{itemize}
\end{boxobjectifs}

\bigskip

% ============================================================================
\section*{Avant de commencer : « quarante projets, tous prioritaires »}
% ============================================================================

Yaoundé, lundi 28~avril~2025. Les groupes de travail remettent leurs fiches. Il y en a quarante. Chaque
direction a écrit « priorité~1 » en tête des siennes. Le total fait 6,8~milliards. M.~OWONA a dit 4,2.

Le directeur des prestations propose de tout garder et d'étaler sur huit ans. M.~ABDOULAYE propose de réduire
chaque projet de 40~\%. Mme~EKANI fait remarquer qu'un site de secours à 60~\% ne secourt rien. Il faudra
retirer des projets entiers. Personne ne veut retirer les siens.

\begin{boxalert}
\textbf{Violation typique.} Tenir dans l'enveloppe en rabotant tous les projets du même pourcentage.

\smallskip
\textbf{Règle~:} un portefeuille se réduit en retirant des projets, non en les amputant. Un projet sous-financé
coûte son budget et ne livre pas son résultat.
\end{boxalert}

% ============================================================================
\section{Du chantier au projet}
\label{sec:c11-projet}
% ============================================================================

La cible a donné des chantiers. Le portefeuille les découpe en projets. Un projet a un résultat livrable, un
responsable, un coût et une durée. La norme ISO~21504 appelle \termecle{portefeuille} l'ensemble des projets et
des programmes réunis pour atteindre des objectifs stratégiques~\cite{iso21504}. Gérer un portefeuille, c'est
choisir, ordonner et réviser.

Chaque projet candidat tient sur une \termecle{fiche de projet} d'une page. L'Annexe~B en donne le modèle.

\begin{tableaucours}{L{3.4cm}L{8.6cm}}{%
    \tableauHeader{Rubrique} &
    \tableauHeader{Contenu}
}
Résultat & Ce qui existera à la fin, en une phrase vérifiable \\
Objectif servi & L'objectif du plan stratégique et la capacité visée \\
Responsable & Un directeur de métier, nommé \\
Coût & Investissement et coût annuel après mise en service \\
Durée & En mois, passation du marché comprise \\
Dépendances & Les projets qui doivent être finis avant \\
Ce qu'il remplace & Les applications retirées \\
Risque principal & Un seul, avec sa parade \\
\end{tableaucours}
\vspace{-0.8em}
\begin{center}{\small\itshape Tableau~11.1. La fiche de projet}\end{center}

À l'ANPP, les groupes de travail ont produit 40~fiches. Leur coût cumulé atteint 6,8~milliards de FCFA.
L'enveloppe indiquée par le directeur général après le choix du scénario~B est de 4,2~milliards sur cinq ans.
Il faut choisir.

% ============================================================================
\section{Trier}
\label{sec:c11-trier}
% ============================================================================

Le tri sépare d'abord les projets \termecle{obligatoires}. Un projet est obligatoire s'il répond à un texte, à
une injonction ou à un risque que l'organisation ne peut pas porter. Ils ne se notent pas, ils s'inscrivent. On
vérifie seulement que le mot n'est pas employé par commodité.

L'ANPP en compte cinq~: le site de secours et le plan de continuité, l'audit de sécurité, la mise en conformité
à la loi sur les données personnelles, le remplacement des onduleurs, la sortie d'un système dont l'éditeur
cesse la maintenance en 2026. Ils pèsent 1~150~millions.

Les 35~autres se notent. La note de valeur, sur 100, reprend les critères du choix de scénario. Le rapport de
la valeur au coût donne un rang.

\begin{tableaucours}{L{0.9cm}L{4.4cm}ccc}{%
    \tableauHeader{N°} &
    \tableauHeader{Projet} &
    \tableauHeader{Valeur} &
    \tableauHeader{Coût (millions)} &
    \tableauHeader{Valeur par million}
}
P15 & Tableau de bord des paiements & 55 & 60 & 0,92 \\
P38 & Refonte du site d'information & 20 & 45 & 0,44 \\
P21 & Comptabilité analytique & 60 & 150 & 0,40 \\
P12 & Guichet unique des demandes & 70 & 180 & 0,39 \\
P07 & Référentiel des bénéficiaires & 90 & 340 & 0,26 \\
P33 & Application mobile des agents & 30 & 120 & 0,25 \\
\end{tableaucours}
\vspace{-0.8em}
\begin{center}{\small\itshape Tableau~11.2. Six projets classés par valeur rapportée au coût}\end{center}

Le classement ne suffit pas. Le référentiel des bénéficiaires arrive cinquième, pourtant la refonte de la
liquidation en dépend. Le site d'information arrive deuxième avec une valeur de 20. Le ratio favorise les
petits projets. On le corrige par deux règles~: un seuil de valeur, fixé à 40 à l'ANPP, en dessous duquel un
projet n'est pas retenu, puis l'examen des dépendances.

\begin{boxalert}
\textbf{Violation typique.} Qualifier d'obligatoire un projet que l'on veut soustraire à l'arbitrage.

\smallskip
\textbf{Règle~:} un projet obligatoire cite le texte, l'injonction ou le risque qui l'impose. Sans cette
référence, il se note comme les autres.
\end{boxalert}

% ============================================================================
\section{Équilibrer}
\label{sec:c11-equilibrer}
% ============================================================================

Un portefeuille se juge aussi dans son ensemble. La grille de McFarlan, vue au Chapitre~\ref{chap:alignement},
sert de contrôle~: un portefeuille tout entier consacré au soutien ne prépare rien, un portefeuille tout entier
stratégique néglige ce qui tourne. On regarde aussi la répartition par objectif, par direction et par niveau de
risque.

\begin{tableaucours}{L{4.6cm}cr}{%
    \tableauHeader{Domaine} &
    \tableauHeader{Projets} &
    \tableauHeader{Montant (millions)}
}
Obligatoires~: secours, sécurité, conformité & 5 & 1 150 \\
Refonte de la production & 6 & 1 820 \\
Référentiels et données & 3 & 520 \\
Relation avec les bénéficiaires & 4 & 410 \\
Soutien et pilotage & 4 & 280 \\
Total retenu & 22 & 4 180 \\
\end{tableaucours}
\vspace{-0.8em}
\begin{center}{\small\itshape Tableau~11.3. Le portefeuille retenu par l'ANPP}\end{center}

Vingt-deux projets sont retenus pour 4~180~millions. Sur les dix-huit autres, onze sont différés après 2029
avec leur fiche conservée, sept sont rejetés. La liste des projets écartés figure au schéma directeur avec le
motif. Elle évite que les mêmes demandes reviennent chaque année par une autre porte.

% ============================================================================
\section{Ordonner}
\label{sec:c11-ordonner}
% ============================================================================

Trois contraintes fixent l'ordre. Les \termecle{dépendances}~: on ne refond pas la liquidation avant d'avoir le
référentiel des bénéficiaires, on ne migre rien avant d'avoir un secours. La \termecle{capacité à faire}~:
l'ANPP ne peut pas conduire plus de six projets à la fois avec ses chefs de projet et ses directions. Les
\termecle{délais de droit public}, vus au Chapitre~\ref{chap:droit}~: inscription budgétaire, programmation,
passation.

La refonte de la production est découpée en trois lots. Chaque lot livre un résultat utilisable seul.

\begin{tableaucours}{L{4.2cm}ccccc}{%
    \tableauHeader{Projet ou lot} &
    \tableauHeader{2025} &
    \tableauHeader{2026} &
    \tableauHeader{2027} &
    \tableauHeader{2028} &
    \tableauHeader{2029}
}
Onduleurs, audit de sécurité & X &  &  &  &  \\
Site de secours, continuité & X & X &  &  &  \\
Référentiel des bénéficiaires & X & X &  &  &  \\
Lot 1. Immatriculation &  & X & X &  &  \\
Lot 2. Liquidation &  &  & X & X &  \\
Lot 3. Paiement, recouvrement &  &  &  & X & X \\
Guichet unique &  & X & X &  &  \\
Comptabilité analytique &  &  & X & X &  \\
Retrait de neuf applications & X & X &  &  &  \\
\end{tableaucours}
\vspace{-0.8em}
\begin{center}{\small\itshape Tableau~11.4. La trajectoire 2025-2029 (extrait)}\end{center}

\begin{boxdef}
\textbf{Trajectoire.} Ordre daté des projets retenus, avec pour chaque année ce qui démarre, ce qui est livré
et ce qui est retiré. Elle se lit année par année~: à la fin de chacune, l'organisation dispose d'un système
qui fonctionne.
\end{boxdef}

% ============================================================================
\section{Montrer des résultats tôt}
\label{sec:c11-tot}
% ============================================================================

Un schéma directeur qui ne livre rien la première année perd ses soutiens. Le portefeuille réserve donc
quelques projets courts, peu coûteux, visibles. À l'ANPP, le tableau de bord des paiements coûte 60~millions
et sera livré en six mois. Le retrait de neuf applications ne coûte presque rien et libère 85~millions par an de
contrats.

Ces résultats rapides ne doivent pas contredire la cible. Un tableau de bord branché sur l'application de 2004
sera rebranché en 2027. Le coût de ce rebranchement est écrit dans la fiche.

% ============================================================================
\section{Les paiements par année}
\label{sec:c11-annees}
% ============================================================================

La trajectoire donne un profil de dépense. Il se compare à ce que l'organisation peut payer chaque année.

\begin{tableaucours}{L{3.4cm}rrrrrr}{%
    \tableauHeader{} &
    \tableauHeader{2025} &
    \tableauHeader{2026} &
    \tableauHeader{2027} &
    \tableauHeader{2028} &
    \tableauHeader{2029} &
    \tableauHeader{Total}
}
Paiements (millions) & 380 & 1 050 & 1 180 & 940 & 630 & 4 180 \\
Projets en cours & 5 & 6 & 6 & 5 & 4 &  \\
\end{tableaucours}
\vspace{-0.8em}
\begin{center}{\small\itshape Tableau~11.5. Le profil de dépense de la trajectoire}\end{center}

L'année 2025 est faible. Les marchés ne seront notifiés qu'au second semestre. Le pic de 2027 dépasse d'un
facteur huit la moyenne annuelle des projets de la période passée, qui était de 136~millions. La trajectoire
n'est donc crédible que si le financement change d'échelle. C'est la question du chapitre suivant.

\begin{boxprocedure}
\textbf{Construire un portefeuille et sa trajectoire.}
\begin{enumerate}
    \item Réunir une fiche par projet candidat.
    \item Isoler les obligatoires, référence à l'appui.
    \item Noter les autres, appliquer le seuil de valeur, classer.
    \item Retenir dans la limite de l'enveloppe. Écrire le motif de chaque rejet.
    \item Vérifier l'équilibre par objectif, par domaine, par risque.
    \item Ordonner selon les dépendances, la capacité à faire et les délais de droit.
    \item Découper les grands projets en lots utilisables seuls.
    \item Tracer le profil annuel et le comparer aux moyens.
\end{enumerate}
\end{boxprocedure}

\begin{boxlecture}{Le portefeuille de l'ANPP, 23 mai 2025}
\begin{tableaucours}{L{7.6cm}cL{3.4cm}}{%
    \tableauHeader{Exigence} &
    \tableauHeader{Cote} &
    \tableauHeader{Preuve}
}
Chaque projet retenu a une fiche complète & \OUI{} & 22 fiches \\
Chaque projet a un responsable côté métier & \OUI{} & Rubrique renseignée sur 22 fiches \\
Chaque projet obligatoire cite sa référence & \OUI{} & Cinq références \\
Le total retenu tient dans l'enveloppe & \OUI{} & 4 180 pour 4 200 \\
Les projets écartés sont listés avec leur motif & \OUI{} & Annexe du portefeuille \\
Les dépendances sont respectées par la trajectoire & \OUI{} & Tableau 11.4 \\
Le nombre de projets simultanés reste dans la capacité & \OUI{} & Six au plus \\
Chaque projet nomme ce qu'il remplace & \NON{} & Sept fiches muettes \\
Le coût après mise en service est chiffré & \NON{} & Neuf fiches sans coût annuel \\
Une provision pour aléas est prévue & \NON{} & Marge de 20 millions seulement \\
\end{tableaucours}

\medskip
\textbf{Décompte.} Sept \OUI{}, trois \NON{}, zéro \NEK{} sur dix.
\hfill\lecture{0{,}70}{0}{0{,}88}

\medskip
\textbf{Lecture.} Sept \OUI{} sur dix. Les trois \NON{} vont dans le même sens~: le portefeuille est bien choisi et mal
bordé. Sans coût annuel, le fonctionnement de 2029 est inconnu. Avec 20~millions de marge sur 4,2~milliards, le
premier dépassement oblige à retirer un projet. Le comité de pilotage a approuvé le portefeuille sous réserve
que la phase de chiffrage lève ces trois points.
\end{boxlecture}

\begin{boxmetier}{Directeur de métier}
Signer comme responsable de ses projets. Accepter par écrit le rejet des autres.
\end{boxmetier}

\begin{boxmetier}{Responsable de portefeuille}
Tenir les fiches à jour. Présenter chaque trimestre l'état des 22 projets sur une page.
\end{boxmetier}

\begin{boxmetier}{Directeur financier}
Rapprocher le profil annuel du cadre de dépenses. Exiger le coût après mise en service.
\end{boxmetier}

\begin{boxmetier}{Consultant}
Faire écrire les motifs de rejet. Refuser les obligatoires sans référence.
\end{boxmetier}

\begin{boxresume}
\begin{itemize}
    \item Un projet tient sur une fiche d'une page.
    \item Les obligatoires s'inscrivent d'abord, référence à l'appui.
    \item Les autres se notent. Le ratio se corrige par un seuil de valeur et par les dépendances.
    \item Un portefeuille se réduit en retirant des projets entiers.
    \item Trois contraintes d'ordre~: dépendances, capacité à faire, délais de droit.
    \item Les grands projets se découpent en lots utilisables seuls.
    \item Le profil annuel se compare aux moyens.
\end{itemize}
\end{boxresume}

\begin{boxinfo}
\textbf{Points de contrôle.} Vous avez compris ce chapitre si vous savez~:
\begin{itemize}
    \item remplir une fiche de projet~;
    \item classer dix projets et justifier les corrections~;
    \item tracer une trajectoire qui respecte des dépendances~;
    \item dresser le profil annuel de dépense.
\end{itemize}
\end{boxinfo}

\begin{boxexo}[title={\ding{45}\quad Exercice~11.1. Classer \niveaubase}]
Calculez la valeur par million des projets suivants et classez-les~: valeur 80 pour 400~millions, valeur 45 pour 90, valeur 35 pour 50, valeur 65 pour 260. Appliquez ensuite le seuil de valeur de 40.
\end{boxexo}

\begin{boxexo}[title={\ding{45}\quad Exercice~11.2. L'enveloppe \niveaumoyen}]
Avec les six projets du Tableau~11.2 et une enveloppe de 600~millions, composez le meilleur portefeuille sachant que P12 dépend de P07. Donnez le total de valeur obtenu et comparez-le à celui d'un choix par ratio seul.
\end{boxexo}

\begin{boxexo}[title={\ding{45}\quad Exercice~11.3. Obligatoire ou non \niveaumoyen}]
Dites si chaque projet de l'ANPP peut être qualifié d'obligatoire et sur quelle référence~: le plan de continuité, le renouvellement des postes de travail, l'audit de sécurité, la refonte du site d'information, le registre des traitements de données personnelles.
\end{boxexo}

\begin{boxexo}[title={\ding{45}\quad Exercice~11.4. La trajectoire contrainte \niveauexpert}]
Le marché du site de secours est déclaré infructueux en 2025 et relancé. Il glisse d'un an. Reprenez le Tableau~11.4, déterminez les projets qui glissent avec lui et ceux qui peuvent avancer. Proposez une trajectoire révisée qui tient dans six projets simultanés.
\end{boxexo}

\bigskip

\begin{boxinfo}
\textbf{Pour aller plus loin.}
\begin{itemize}
    \item ISO~21504~\cite{iso21504}.
    \item McFarlan sur le portefeuille d'applications~\cite{mcfarlan1984}.
    \item Weill et Ross sur les décisions d'investissement~\cite{weillross2004}.
\end{itemize}
\end{boxinfo}

\begin{boxinfo}
\textbf{Transition.} Le portefeuille est choisi et ordonné. Il reste à le chiffrer complètement, à le financer, à l'acheter et à trouver les personnes. C'est l'objet du Chapitre~\ref{chap:financer}.
\end{boxinfo}

        % ============================================================================
% Fichier  : partie4/P04_C12_financer.tex
% Titre    : Chiffrer, Financer, Acheter, Doter
% Auteur   : MINKA MI NGUIDJOÏ Thierry Emmanuel
% Version  : 1.0 -- novembre 2026
% Corrigés : corriges/Corr_C12.tex
% ============================================================================

\chapter{Chiffrer, Financer, Acheter, Doter}
\label{chap:financer}
\chaptermark{Financer}

\epigraph{%
    « Un plan sans ressources est un vœu. »%
}{Principe du cours}

\niveauChapitre{Intermédiaire}{%
    Chapitres~\ref{chap:droit} et \ref{chap:portefeuille}
}

\begin{boxobjectifs}
\begin{itemize}
    \item Calculer le coût complet d'un projet.
    \item Donner une estimation avec sa fourchette et sa provision.
    \item Calculer une valeur actuelle nette et en dire les limites.
    \item Présenter un projet en engagements et en paiements.
    \item Établir un plan de financement avec un scénario de repli.
    \item Choisir entre faire, acheter et louer.
    \item Établir un plan de dotation.
\end{itemize}
\end{boxobjectifs}

\bigskip

% ============================================================================
\section*{Avant de commencer : « qui paiera ensuite ? »}
% ============================================================================

Yaoundé, mardi 25~mai~2027. M.~ABDOULAYE relit les 22~fiches. Il additionne les investissements~: 4~180~millions.
Il cherche ensuite la colonne des coûts annuels. Neuf fiches n'en ont pas. Il appelle Mme~EKANI. « Vous me
demandez d'approuver l'achat de vingt-deux voitures sans me dire ce qu'elles consomment. »

Le cabinet reprend les neuf fiches en une semaine. Le résultat surprend le comité~: en 2031, l'Agence paiera
chaque année 225~millions de plus qu'aujourd'hui pour faire tourner ce qu'elle aura construit. Personne ne
l'avait vu. Le schéma directeur de 2021 ne le disait pas non plus.

\begin{boxalert}
\textbf{Violation typique.} Présenter au conseil un total d'investissement sans la charge récurrente qu'il
crée.

\smallskip
\textbf{Règle~:} l'organe qui approuve un schéma directeur approuve deux chiffres~: l'enveloppe de la période
et le fonctionnement annuel à son terme.
\end{boxalert}

% ============================================================================
\section{Le Coût Complet}
\label{sec:c12-cout}
% ============================================================================

Le coût d'un projet ne s'arrête pas à sa mise en service. Le \termecle{coût complet de possession} additionne
l'investissement et tout ce que l'organisation paiera ensuite pour faire vivre le résultat.

\begin{tableaucours}{L{3.0cm}L{9.0cm}}{%
    \tableauHeader{Poste} &
    \tableauHeader{Contenu}
}
Investissement & Logiciels, matériels, prestations d'intégration, reprise des données, formation, conduite du changement \\
Charges internes & Temps des agents du métier et de l'informatique affectés au projet \\
Fonctionnement & Maintenance, licences ou abonnements, hébergement, liaisons, énergie, exploitation \\
Renouvellement & Remplacement des matériels, montées de version \\
Sortie & Reprise des données et retrait en fin de vie \\
\end{tableaucours}
\vspace{-0.8em}
\begin{center}{\small\itshape Tableau~12.1. Les cinq postes du coût complet}\end{center}

Le référentiel des bénéficiaires de l'ANPP demande 340~millions d'investissement, payés 60 en 2027 et 280 en
2028. Il coûte ensuite 48~millions par an à partir de 2029. Sur la période du schéma directeur, son coût
complet est de $340 + 3 \times 48 = 484$~millions. La fiche de projet n'affichait que 340.

\begin{boxalert}
\textbf{Violation typique.} Chiffrer un schéma directeur en investissement seul.

\smallskip
\textbf{Règle~:} chaque projet affiche deux montants~: l'investissement et le coût annuel après mise en
service. Le schéma directeur affiche la dépense de fonctionnement de sa dernière année.
\end{boxalert}

% ============================================================================
\section{Estimer}
\label{sec:c12-estimer}
% ============================================================================

Trois méthodes d'estimation se combinent. L'\termecle{analogie} part d'un projet comparable déjà réalisé. Le
\termecle{détail} additionne les éléments un à un. Le \termecle{prix de marché} interroge des fournisseurs sans
les engager. Aucune n'est exacte. Chaque estimation se donne donc avec une fourchette.

La fourchette se resserre avec l'avancement. En phase de schéma directeur, une estimation à plus ou moins
30~\% est honnête. Elle se couvre par une \termecle{provision pour aléas}. L'ANPP avait retenu 20~millions de
marge sur 4,2~milliards, soit moins de 0,5~\%. Le cabinet a proposé 8~\% des projets non obligatoires.

% ============================================================================
\section{La Valeur d'un Projet}
\label{sec:c12-valeur}
% ============================================================================

Quand un projet rapporte ou économise de l'argent, on compare ce qu'il coûte à ce qu'il rapporte, en tenant
compte du temps. Un franc reçu dans cinq ans vaut moins qu'un franc reçu aujourd'hui. La
\termecle{valeur actuelle nette} ramène tous les flux à la date de départ.

\begin{equation}
    \mathrm{VAN} = -I_0 + \sum_{t=1}^{n} \frac{F_t}{(1+r)^t}
\end{equation}

$I_0$ est l'investissement, $F_t$ le flux net de l'année $t$, $r$ le taux d'actualisation, $n$ la durée.

Le guichet unique des demandes coûte 180~millions. Il réduit les déplacements d'agents et les doubles saisies
pour un gain net estimé à 55~millions par an pendant cinq ans. Au taux de 8~\%, la somme actualisée des gains
vaut $55 \times 3{,}993 = 219{,}6$~millions. La valeur actuelle nette est de $219{,}6 - 180 = 39{,}6$~millions.
Le projet se paie.

\begin{boxinfo}
\textbf{Limite dans le secteur public.} La plupart des bénéfices d'une administration ne sont pas des recettes.
Un paiement à date fixe ne rapporte rien à l'Agence. Il évite à un retraité de s'endetter. Ward et Daniel
proposent de classer les bénéfices selon qu'ils sont financiers, quantifiables, mesurables ou seulement
observables~\cite{warddaniel2006}. Un bénéfice observable reste un bénéfice, à condition de dire qui
l'observera et comment.
\end{boxinfo}

% ============================================================================
\section{Engager et Payer}
\label{sec:c12-aecp}
% ============================================================================

Le Chapitre~\ref{chap:droit} a distingué autorisations d'engagement et crédits de paiement. Chaque projet se
présente donc sur deux lignes.

\begin{tableaucours}{L{4.6cm}rrrr}{%
    \tableauHeader{Lot 2. Liquidation (millions)} &
    \tableauHeader{2028} &
    \tableauHeader{2029} &
    \tableauHeader{2030} &
    \tableauHeader{Total}
}
Autorisations d'engagement & 780 &  &  & 780 \\
Crédits de paiement & 150 & 420 & 210 & 780 \\
\end{tableaucours}
\vspace{-0.8em}
\begin{center}{\small\itshape Tableau~12.2. Engagements et paiements d'un lot}\end{center}

L'engagement est pris en 2028 pour la totalité, sans quoi le marché ne peut pas être signé. Les paiements
suivent les livraisons. Si les crédits de 2029 sont réduits, le fournisseur n'est pas payé à l'échéance, le
projet ralentit et des intérêts moratoires courent.

% ============================================================================
\section{Financer}
\label{sec:c12-financer}
% ============================================================================

Le plan de financement met en face de chaque année une ressource. Il en existe quatre sortes pour un
établissement public~: les ressources propres, la subvention de l'État, le concours d'un partenaire technique et
financier, l'économie dégagée par le schéma directeur lui-même.

\begin{tableaucours}{L{5.0cm}rL{4.4cm}}{%
    \tableauHeader{Ressource} &
    \tableauHeader{Montant (millions)} &
    \tableauHeader{Condition}
}
Ressources propres de l'Agence & 2 300 & Délibération annuelle du conseil \\
Subvention d'investissement de l'État & 1 200 & Inscription en loi de finances \\
Partenaire technique et financier & 680 & Convention à signer, fléchée sur le secours et le référentiel \\
Total & 4 180 &  \\
\end{tableaucours}
\vspace{-0.8em}
\begin{center}{\small\itshape Tableau~12.3. Le plan de financement de l'ANPP}\end{center}

Une ressource n'est acquise que lorsqu'un acte l'établit. Le plan distingue donc trois états~: acquise,
demandée, espérée. Au 11~juin~2027, seules les ressources propres de 2027 sont acquises. Le schéma directeur
l'écrit et prévoit ce qui serait retiré si la subvention ou la convention manquait. C'est le
\termecle{scénario de repli}. Celui de l'ANPP ramène le portefeuille à 2~300~millions en conservant les cinq
projets obligatoires, le référentiel et le premier lot de la refonte.

Le fonctionnement change aussi. En 2031, les nouveaux systèmes ajoutent 310~millions de charges annuelles. Le
retrait de neuf applications en retire 85. Le fonctionnement annuel augmente donc de 225~millions. Le conseil
d'administration doit le savoir le jour où il approuve.

% ============================================================================
\section{Faire, Acheter, Louer}
\label{sec:c12-acheter}
% ============================================================================

Pour chaque besoin, trois voies s'ouvrent. Coase puis Williamson ont montré que le choix dépend des coûts de
transaction~\cite{coase1937,williamson1975}~: on garde en interne ce qui est spécifique, incertain et fréquent,
on achète le reste.

\begin{tableaucours}{L{2.2cm}L{4.6cm}L{5.2cm}}{%
    \tableauHeader{Voie} &
    \tableauHeader{Convient quand} &
    \tableauHeader{Point de vigilance}
}
Faire & Le besoin est propre à l'organisation et au cœur du métier & Dépendance à quelques agents \\
Acheter & Un produit du marché couvre le besoin commun & Coût des adaptations, dépendance à l'éditeur \\
Louer & Le besoin varie ou l'investissement est impossible & Localisation des données, réversibilité, charge récurrente \\
\end{tableaucours}
\vspace{-0.8em}
\begin{center}{\small\itshape Tableau~12.4. Trois voies d'acquisition}\end{center}

Louer un service en nuage transforme un investissement en charge de fonctionnement. Pour une administration,
cela déplace la dépense d'une ligne budgétaire à une autre, ce qui n'est pas toujours possible. Cela pose aussi
la question du lieu des données, que la loi n°~2024/017 encadre pour les transferts hors du
Cameroun~\cite{loi2024017}.

Tout contrat prévoit sa fin. La \termecle{réversibilité} oblige le prestataire à restituer les données dans un
format utilisable et à assister la reprise. Elle se négocie à la signature.

Le plan d'achat du schéma directeur liste les marchés, leur type, leur montant estimé et leur date de
lancement. L'ANPP en prévoit quatorze pour ses vingt-deux projets. Le code des marchés publics en fixe la
procédure~\cite{decret2018366}.

% ============================================================================
\section{Doter}
\label{sec:c12-doter}
% ============================================================================

Un portefeuille de 4,2~milliards ne se conduit pas avec une division conçue pour en dépenser 136~millions par
an. Le plan de dotation répond à trois questions~: quelles compétences, combien, par quelle voie.

\begin{tableaucours}{L{4.6cm}cL{5.4cm}}{%
    \tableauHeader{Compétence} &
    \tableauHeader{Postes} &
    \tableauHeader{Voie}
}
Responsable de la sécurité & 1 & Recrutement en 2027 \\
Chefs de projet & 2 & Un recrutement, une mobilité interne \\
Architecte & 1 & Recrutement en 2028 \\
Données & 3 & Deux recrutements, une reconversion \\
Acheteur informatique & 1 & Mobilité depuis les marchés \\
\end{tableaucours}
\vspace{-0.8em}
\begin{center}{\small\itshape Tableau~12.5. Le plan de dotation de l'ANPP~: huit postes}\end{center}

La division passe de 19 à 27~agents. Les postes sont décrits avec le référentiel européen des
compétences~\cite{ecf}. Le plan prévoit aussi la formation des agents en place et le transfert de compétences
dans chaque marché. Les deux agents qui connaissent la liquidation sont affectés à la refonte. Leur
remplacement à la maintenance est organisé avant.

\begin{boxprocedure}
\textbf{Établir le plan de financement.}
\begin{enumerate}
    \item Compléter chaque fiche~: investissement, coût annuel, fourchette.
    \item Ajouter la provision pour aléas.
    \item Présenter chaque projet en engagements et en paiements.
    \item Additionner par année. Comparer au profil des ressources.
    \item Qualifier chaque ressource~: acquise, demandée, espérée.
    \item Écrire le scénario de repli.
    \item Calculer le fonctionnement de la dernière année.
    \item Joindre le plan d'achat et le plan de dotation.
\end{enumerate}
\end{boxprocedure}

\begin{boxlecture}{Le plan de financement de l'ANPP, 11 juin 2027}
\begin{tableaucours}{L{7.6cm}cL{3.4cm}}{%
    \tableauHeader{Exigence} &
    \tableauHeader{Cote} &
    \tableauHeader{Preuve}
}
Chaque projet affiche investissement et coût annuel & \OUI{} & 22 fiches complétées \\
Une provision pour aléas est chiffrée & \OUI{} & 8~\% des projets non obligatoires \\
Chaque projet est présenté en engagements et en paiements & \OUI{} & Tableaux du plan \\
Chaque ressource est qualifiée & \OUI{} & Tableau 12.3 et colonne d'état \\
Un scénario de repli est écrit & \OUI{} & Portefeuille à 2 300 millions \\
Le fonctionnement de 2031 est calculé & \OUI{} & Hausse de 225 millions par an \\
Le total avec provision tient dans l'enveloppe & \NON{} & 4 422 pour 4 200 \\
Le plan d'achat est rapproché de la programmation & \OUI{} & Quatorze marchés \\
Les huit postes sont autorisés & \NEK{} & Demande adressée à la tutelle, réponse attendue \\
\end{tableaucours}

\medskip
\textbf{Décompte.} Sept \OUI{}, un \NON{}, un \NEK{} sur neuf.
\hfill\lecture{0{,}778}{0{,}111}{0{,}99}

\medskip
\textbf{Lecture.} Sept \OUI{} sur neuf. Le \NON{} est le prix de l'honnêteté. La provision de 8~\% porte sur 3~030~millions
de projets non obligatoires, soit 242~millions. Le total atteint 4~422~millions. Le comité de pilotage a deux
voies~: relever l'enveloppe ou retirer un projet. Il a retiré la refonte du site d'information et un projet de
soutien. Le \NEK{} est de preuve et il est grave~: sans les postes, la trajectoire ne tient pas.
\end{boxlecture}

\begin{boxmetier}{Directeur financier}
Qualifier chaque ressource. Refuser un plan dont tout le financement est espéré.
\end{boxmetier}

\begin{boxmetier}{Acheteur}
Traduire le portefeuille en marchés. Caler leur lancement sur la programmation. Écrire les clauses de réversibilité.
\end{boxmetier}

\begin{boxmetier}{Responsable des ressources humaines}
Faire autoriser les postes avant le lancement des projets qui en dépendent.
\end{boxmetier}

\begin{boxmetier}{Consultant}
Donner des fourchettes. Chiffrer la provision. Écrire le scénario de repli sans attendre qu'on le demande.
\end{boxmetier}

\begin{boxresume}
\begin{itemize}
    \item Le coût complet ajoute à l'investissement le fonctionnement, le renouvellement et la sortie.
    \item Une estimation se donne avec une fourchette et se couvre par une provision.
    \item La valeur actuelle nette compare des flux ramenés à la date de départ. Elle ne saisit pas les bénéfices non financiers.
    \item Un projet public se présente en engagements et en paiements.
    \item Une ressource est acquise, demandée ou espérée.
    \item Le scénario de repli dit ce qui reste si une ressource manque.
    \item Le conseil approuve l'enveloppe et le fonctionnement annuel à terme.
    \item Le plan de dotation précède les projets.
\end{itemize}
\end{boxresume}

\begin{boxinfo}
\textbf{Points de contrôle.} Vous avez compris ce chapitre si vous savez~:
\begin{itemize}
    \item calculer le coût complet d'un projet sur cinq ans~;
    \item calculer une valeur actuelle nette~;
    \item dresser le tableau des engagements et des paiements~;
    \item écrire un scénario de repli.
\end{itemize}
\end{boxinfo}

\begin{boxexo}[title={\ding{45}\quad Exercice~12.1. Coût complet \niveaubase}]
Un projet demande 200~millions d'investissement en 2027 et 30~millions par an à partir de 2028. Calculez son coût complet sur 2027-2031.
\end{boxexo}

\begin{boxexo}[title={\ding{45}\quad Exercice~12.2. Valeur actuelle nette \niveaumoyen}]
Un projet coûte 120~millions et procure un gain net de 40~millions par an pendant quatre ans. Calculez sa valeur actuelle nette au taux de 8~\%, puis au taux de 15~\%. Commentez.
\end{boxexo}

\begin{boxexo}[title={\ding{45}\quad Exercice~12.3. La provision \niveaumoyen}]
Vérifiez les chiffres de la lecture du plan de financement~: montant des projets non obligatoires, provision, total. De combien faut-il réduire le portefeuille pour tenir dans 4~200~millions, provision comprise~?
\end{boxexo}

\begin{boxexo}[title={\ding{45}\quad Exercice~12.4. Le repli \niveauexpert}]
La convention avec le partenaire technique et financier n'est pas signée à la fin de 2027. Rédigez la note au comité de pilotage~: projets touchés, conséquences sur la trajectoire, options, recommandation.
\end{boxexo}

\begin{boxexo}[title={\ding{45}\quad Exercice~12.5. Louer ou acheter \niveauexpert}]
La BANQUE SANAGA hésite entre acheter ses serveurs de secours et louer un service d'hébergement hors de la CEMAC. Comparez les deux voies sur cinq critères, dont le regard du superviseur et la réversibilité.
\end{boxexo}

\bigskip

\begin{boxinfo}
\textbf{Pour aller plus loin.}
\begin{itemize}
    \item Ward et Daniel sur la gestion des bénéfices~\cite{warddaniel2006}.
    \item La loi n°~2018/012~\cite{loi2018012}.
    \item Williamson sur les coûts de transaction~\cite{williamson1975}.
\end{itemize}
\end{boxinfo}

\begin{boxinfo}
\textbf{Transition.} Le schéma directeur est complet. Le conseil d'administration l'approuve le 24~juin~2027. Le plus difficile commence~: l'exécuter. C'est l'objet du Chapitre~\ref{chap:piloter}.
\end{boxinfo}


    \part{Mise en œuvre : piloter et maîtriser}
    \label{part:pilotage}

        % ============================================================================
% Fichier  : partie5/P05_C13_piloter.tex
% Titre    : Piloter l'exécution
% Auteur   : MINKA MI NGUIDJOÏ Thierry Emmanuel
% Version  : 1.0 -- novembre 2026
% Corrigés : corriges/Corr_C13.tex
% ============================================================================

\chapter{Piloter l'exécution}
\label{chap:piloter}
\chaptermark{Piloter}

\epigraph{%
    « Aucun plan ne résiste au premier contact. »%
}{D'après Helmuth von Moltke, 1871}

\niveauChapitre{Intermédiaire}{%
    Chapitres~\ref{chap:portefeuille} et \ref{chap:financer}
}

\begin{boxobjectifs}
\begin{itemize}
    \item Composer les instances d'exécution d'un schéma directeur.
    \item Construire un tableau de bord sur quatre axes.
    \item Calculer les indices de délai et de coût par la valeur acquise.
    \item Tenir un registre des bénéfices.
    \item Instruire une demande hors portefeuille.
    \item Conduire une révision annuelle.
\end{itemize}
\end{boxobjectifs}

\bigskip

% ============================================================================
\section*{Avant de commencer : « tout avance normalement »}
% ============================================================================

Yaoundé, mardi 30~juin~2026. Première revue annuelle. Les chefs de projet présentent leurs fiches. Toutes
portent la mention « avancement conforme ». M.~OWONA s'en réjouit.

La cellule de portefeuille distribue ensuite une page. Elle a compté ce qui est réellement achevé. Sur
905~millions de travaux prévus à cette date, 540 sont faits. Le site de secours n'a pas de fournisseur. Le
référentiel a quatre mois de retard. Le directeur général demande pourquoi personne ne l'a dit. « Chaque projet
avance, répond l'un des agents de la cellule. C'est le portefeuille qui est en retard. »

\begin{boxalert}
\textbf{Violation typique.} Piloter un portefeuille par l'addition des déclarations des chefs de projet.

\smallskip
\textbf{Règle~:} l'avancement se mesure par une instance qui ne conduit aucun projet, avec une méthode
identique pour tous.
\end{boxalert}

% ============================================================================
\section{Qui suit l'exécution}
\label{sec:c13-instances}
% ============================================================================

Les instances de la mission se dissolvent le jour de l'approbation. D'autres les remplacent pour cinq ans. La
norme ISO/IEC~38500 confie à l'organe de gouvernance la tâche de surveiller~\cite{iso38500}. Weill et Ross
ajoutent que chaque type de décision doit avoir un titulaire connu~\cite{weillross2004}.

\begin{tableaucours}{L{3.2cm}L{3.6cm}L{2.0cm}L{3.2cm}}{%
    \tableauHeader{Instance} &
    \tableauHeader{Composition} &
    \tableauHeader{Rythme} &
    \tableauHeader{Décide de}
}
Conseil d'administration & Administrateurs & Chaque année & Révision, budget \\
Comité de portefeuille & Directeur général, directeurs & Chaque trimestre & Lancement, arrêt, arbitrage entre projets \\
Comité de projet & Responsable métier, chef de projet, fournisseur & Chaque mois & Conduite d'un projet \\
Cellule de portefeuille & Deux agents auprès du directeur général & En continu & Rien. Elle instruit et alerte \\
\end{tableaucours}
\vspace{-0.8em}
\begin{center}{\small\itshape Tableau~13.1. Les instances d'exécution à l'ANPP}\end{center}

La \termecle{cellule de portefeuille} est l'instance la plus modeste et la plus utile. Elle tient les fiches à
jour, prépare le comité, vérifie les chiffres des chefs de projet. Elle n'est pas rattachée à la division
informatique. À l'ANPP, elle est créée par une décision du 4~juillet~2025.

% ============================================================================
\section{Le tableau de bord}
\label{sec:c13-tableau}
% ============================================================================

Kaplan et Norton ont proposé en 1992 de suivre une organisation sur quatre axes plutôt que sur le seul résultat
financier~\cite{kaplannorton1992}. Transposés au schéma directeur, ils donnent une douzaine d'indicateurs au
plus.

\begin{tableaucours}{L{2.8cm}L{4.8cm}L{4.4cm}}{%
    \tableauHeader{Axe} &
    \tableauHeader{Indicateur} &
    \tableauHeader{Cible de l'ANPP pour 2029}
}
Bénéficiaires & Prestations payées à la date prévue & 99~\% \\
Bénéficiaires & Délai de traitement d'une demande & 15~jours \\
Finances & Part des projets dans la dépense & 35~\% \\
Finances & Écart de coût du portefeuille & Moins de 10~\% \\
Processus & Projets livrés à la date prévue & 70~\% \\
Processus & Applications sans responsable & 0 \\
Capacités & Postes du plan de dotation pourvus & 8 sur 8 \\
Capacités & Applications clés avec reprise éprouvée & Toutes \\
\end{tableaucours}
\vspace{-0.8em}
\begin{center}{\small\itshape Tableau~13.2. Le tableau de bord du schéma directeur}\end{center}

Un indicateur se définit par cinq éléments~: une formule, une source, un responsable de la mesure, une
fréquence, une cible. Un indicateur sans source ne sera pas produit. Un indicateur sans cible ne dit rien.

\begin{boxalert}
\textbf{Violation typique.} Suivre l'avancement des projets et jamais leurs effets.

\smallskip
\textbf{Règle~:} le tableau de bord contient au moins un indicateur de résultat pour le bénéficiaire du
service. Vingt-deux projets livrés à l'heure sans que les paiements arrivent à date fixe sont un échec.
\end{boxalert}

% ============================================================================
\section{Mesurer l'avancement}
\label{sec:c13-avancement}
% ============================================================================

Le pourcentage d'avancement déclaré par un chef de projet est peu fiable. La méthode de la valeur acquise
compare trois montants à une date donnée~: ce qui devait être fait, ce qui est fait, ce qui a été dépensé. La
norme ISO~21502 la range parmi les pratiques de maîtrise~\cite{iso21502}.

\begin{equation}
    \mathrm{IP_d} = \frac{\mathrm{VA}}{\mathrm{VP}} \qquad\qquad \mathrm{IP_c} = \frac{\mathrm{VA}}{\mathrm{CR}}
\end{equation}

VP est la valeur planifiée, budget des travaux prévus à la date. VA est la valeur acquise, budget des travaux
réellement achevés. CR est le coût réel. $\mathrm{IP_d}$ mesure le respect des délais, $\mathrm{IP_c}$ le
respect des coûts. Au-dessous de 1, le portefeuille est en retard ou trop cher.

\begin{boxscenario}{ANPP. La première revue, 30 juin 2026}
Le profil du Tableau~11.5 prévoyait 380~millions en 2025 et la moitié de 1~050 au premier semestre 2026. La
valeur planifiée au 30~juin~2026 est donc de $380 + 525 = 905$~millions. La cellule de portefeuille établit que
les travaux achevés représentent 540~millions de budget et que 610~millions ont été payés.

\[ \mathrm{IP_d} = \frac{540}{905} = 0{,}60 \qquad \mathrm{IP_c} = \frac{540}{610} = 0{,}89 \]

Le portefeuille a réalisé 60~\% de ce qui était prévu. Chaque franc dépensé a produit 89~centimes de travaux
prévus. Le retard vient de deux marchés~: le site de secours, relancé après une première procédure
infructueuse, le référentiel, notifié avec quatre mois de retard.
\end{boxscenario}

% ============================================================================
\section{Suivre les bénéfices}
\label{sec:c13-benefices}
% ============================================================================

Un projet livré n'a encore rien rapporté. Ward et Daniel demandent de désigner pour chaque bénéfice un
\termecle{propriétaire}, de le mesurer avant le projet, puis de le mesurer après~\cite{warddaniel2006}. Sans
mesure de départ, aucun progrès ne se prouve.

\begin{tableaucours}{L{3.6cm}L{2.8cm}L{2.6cm}L{2.6cm}}{%
    \tableauHeader{Bénéfice attendu} &
    \tableauHeader{Propriétaire} &
    \tableauHeader{Mesure de départ} &
    \tableauHeader{Cible}
}
Paiement à date fixe & Directeur des prestations & 91~\% en 2024 & 99~\% en 2029 \\
Fin des doubles paiements & Directeur financier & 2~340 cas en 2024 & Moins de 200 \\
Délai de traitement & Chefs d'antenne & 47~jours & 15~jours \\
Économie de contrats & Cheffe de division & 0 & 85~millions par an \\
\end{tableaucours}
\vspace{-0.8em}
\begin{center}{\small\itshape Tableau~13.3. Le registre des bénéfices (extrait)}\end{center}

Le bénéfice appartient au métier. La fonction informatique livre un outil. Le directeur des prestations
obtient, ou n'obtient pas, le paiement à date fixe. Cela suppose qu'il change l'organisation du travail autour
de l'outil. Le Chapitre~\ref{chap:alignement} l'a montré~: la valeur vient de l'usage.

% ============================================================================
\section{Faire adopter}
\label{sec:c13-adopter}
% ============================================================================

Rogers a décrit comment une nouveauté se répand dans un groupe~: quelques pionniers, une majorité qui suit, des
retardataires~\cite{rogers1962}. Davis a montré que l'usage d'un système dépend de deux perceptions, son
utilité et sa facilité~\cite{davis1989}. Ces deux travaux fondent trois gestes simples.

\begin{enumerate}
    \item Choisir des antennes pilotes parmi les volontaires, non parmi les plus proches du siège.
    \item Montrer aux agents ce que l'outil leur épargne avant de montrer ce qu'il permet de contrôler.
    \item Former au poste de travail, sur des dossiers réels, la semaine de la mise en service.
\end{enumerate}

Le budget de formation et d'accompagnement figure dans chaque fiche de projet. À l'ANPP il représente 9~\% de
l'investissement. Une fiche sans ce budget est renvoyée à son auteur.

% ============================================================================
\section{Accueillir l'imprévu}
\label{sec:c13-imprevu}
% ============================================================================

Entre 2019 et 2023, onze projets de l'ANPP se sont faits hors du schéma directeur. Mintzberg appelle cela la
stratégie émergente~\cite{mintzberg1985}. L'interdire est vain. Il faut une porte d'entrée.

\begin{boxprocedure}
\textbf{Instruire une demande hors portefeuille.}
\begin{enumerate}
    \item Le demandeur remplit une fiche de projet, comme pour les autres.
    \item La cellule de portefeuille la note sur la grille de valeur et vérifie le plan fonctionnel.
    \item Le comité de portefeuille décide~: rejet, attente de la révision annuelle ou entrée immédiate.
    \item Une entrée immédiate nomme le projet qu'elle décale ou la ressource nouvelle qui la finance.
    \item La décision est inscrite au registre du portefeuille.
\end{enumerate}
\end{boxprocedure}

Une réserve pour projets non prévus facilite la règle. Les 242~millions de provision pour aléas ne servent pas
à cela. L'ANPP a choisi de ne pas créer de réserve et d'appliquer la règle du décalage.

% ============================================================================
\section{Réviser}
\label{sec:c13-reviser}
% ============================================================================

Un schéma directeur de cinq ans rédigé une fois est faux au bout de dix-huit mois. Il se révise chaque année, à
date fixe, avant la préparation du budget. La révision répond à quatre questions.

\begin{enumerate}
    \item Les objectifs de l'organisation ont-ils changé~?
    \item Qu'a-t-on livré, dépensé, obtenu~?
    \item Les hypothèses de financement tiennent-elles~?
    \item Que faut-il avancer, décaler, retirer, ajouter~?
\end{enumerate}

La révision produit un document court, approuvé par le conseil. Elle peut glisser l'horizon d'un an. On parle
alors de \termecle{schéma directeur glissant}. La cible ne se révise pas chaque année. Si elle doit changer,
c'est qu'un fait majeur est survenu~: une réforme, une fusion, une rupture technique. On rouvre alors une
démarche complète.

\begin{boxinfo}
\textbf{Arrêter un projet.} Le comité de portefeuille a le pouvoir d'arrêter. C'est sa décision la plus
difficile, parce que l'argent déjà dépensé pèse sur le jugement. La question à poser est simple~: sachant ce
que nous savons, lancerions-nous ce projet aujourd'hui pour ce qu'il reste à payer~? Entre 2019 et 2023, quatre
projets de l'ANPP ont été arrêtés sans qu'aucune instance ne l'ait décidé. Ils se sont éteints.
\end{boxinfo}

\begin{boxlecture}{Le dispositif de pilotage de l'ANPP, revue du 30 juin 2026}
\begin{tableaucours}{L{7.6cm}cL{3.4cm}}{%
    \tableauHeader{Exigence} &
    \tableauHeader{Cote} &
    \tableauHeader{Preuve}
}
Un comité de portefeuille se réunit chaque trimestre & \OUI{} & Quatre procès-verbaux \\
La cellule de portefeuille est indépendante de l'informatique & \OUI{} & Décision du 4 juillet 2025 \\
Chaque indicateur a une source et une cible & \OUI{} & Fiches d'indicateur \\
Le tableau de bord contient un indicateur de résultat & \OUI{} & Paiements à date \\
L'avancement est mesuré par la valeur acquise & \OUI{} & Note de la cellule \\
Chaque bénéfice a un propriétaire et une mesure de départ & \NON{} & Deux bénéfices sans mesure de départ \\
Les demandes hors portefeuille passent par la procédure & \NON{} & Un projet lancé par une direction en mars 2026 \\
La révision annuelle est approuvée par le conseil & \NEK{} & Conseil prévu en septembre \\
Les projets en retard ont un plan de rattrapage & \OUI{} & Deux plans joints \\
\end{tableaucours}

\medskip
\textbf{Décompte.} Six \OUI{}, deux \NON{}, un \NEK{} sur neuf.
\hfill\lecture{0{,}667}{0{,}111}{1{,}22}

\medskip
\textbf{Lecture.} Six \OUI{} sur neuf. Le dispositif existe et fonctionne. Les deux \NON{} montrent que les habitudes
résistent~: une direction a encore lancé son projet seule. Le comité l'a régularisé au lieu de l'instruire. Le
\NEK{} est de preuve et il est normal à cette date.
\end{boxlecture}

\begin{boxmetier}{Directeur général}
Présider le comité de portefeuille. Arrêter les projets qui doivent l'être.
\end{boxmetier}

\begin{boxmetier}{Responsable de la cellule de portefeuille}
Mesurer sans conduire. Alerter par écrit. Tenir le registre des décisions.
\end{boxmetier}

\begin{boxmetier}{Directeur de métier}
Mesurer le bénéfice avant le projet. En répondre après.
\end{boxmetier}

\begin{boxmetier}{Auditeur interne}
Vérifier chaque année un échantillon d'indicateurs en remontant à leur source.
\end{boxmetier}

\begin{boxresume}
\begin{itemize}
    \item Quatre instances~: conseil, comité de portefeuille, comités de projet, cellule de portefeuille.
    \item Un indicateur a une formule, une source, un responsable, une fréquence, une cible.
    \item La valeur acquise donne deux indices~: délai et coût.
    \item Un bénéfice a un propriétaire dans le métier et une mesure de départ.
    \item Une demande hors portefeuille entre par une procédure qui nomme ce qu'elle décale.
    \item Le schéma directeur se révise chaque année avant le budget.
    \item Le comité de portefeuille peut arrêter un projet.
\end{itemize}
\end{boxresume}

\begin{boxinfo}
\textbf{Points de contrôle.} Vous avez compris ce chapitre si vous savez~:
\begin{itemize}
    \item définir un indicateur en cinq éléments~;
    \item calculer les deux indices de la valeur acquise~;
    \item remplir une ligne du registre des bénéfices~;
    \item préparer une révision annuelle.
\end{itemize}
\end{boxinfo}

\begin{boxexo}[title={\ding{45}\quad Exercice~13.1. Deux indices \niveaubase}]
À une date donnée, un portefeuille présente une valeur planifiée de 1~200~millions, une valeur acquise de 900 et un coût réel de 1~000. Calculez les deux indices et commentez.
\end{boxexo}

\begin{boxexo}[title={\ding{45}\quad Exercice~13.2. Un indicateur \niveaumoyen}]
Définissez en cinq éléments l'indicateur « applications clés avec reprise éprouvée ». Dites ce que signifie « éprouvée » et quelle pièce le prouve.
\end{boxexo}

\begin{boxexo}[title={\ding{45}\quad Exercice~13.3. Le projet de mars \niveaumoyen}]
Une direction de l'ANPP a lancé seule un projet en mars 2026. Rédigez la décision que le comité de portefeuille aurait dû prendre, selon la procédure de ce chapitre.
\end{boxexo}

\begin{boxexo}[title={\ding{45}\quad Exercice~13.4. Arrêter \niveauexpert}]
Supposez qu'à une revue ultérieure le lot~1 de la refonte ait consommé 80~\% de son budget pour 45~\% de ses travaux. Calculez l'indice de coût, estimez le coût final si la tendance se prolonge et rédigez la note d'options au comité~: poursuivre, redéfinir, arrêter.
\end{boxexo}

\bigskip

\begin{boxinfo}
\textbf{Pour aller plus loin.}
\begin{itemize}
    \item Kaplan et Norton~\cite{kaplannorton1992}.
    \item Ward et Daniel~\cite{warddaniel2006}.
    \item ISO~21502~\cite{iso21502} et ISO~21504~\cite{iso21504}.
    \item Weill et Ross~\cite{weillross2004}.
\end{itemize}
\end{boxinfo}

\begin{boxinfo}
\textbf{Transition.} Piloter suppose que rien de grave n'arrive. Le Chapitre~\ref{chap:maitriser} traite de ce qui peut arriver.
\end{boxinfo}

        % ============================================================================
% Fichier  : partie5/P05_C14_maitriser.tex
% Titre    : Risques, Sécurité, Continuité, Conformité
% Auteur   : MINKA MI NGUIDJOÏ Thierry Emmanuel
% Version  : 1.0 -- novembre 2026
% Corrigés : corriges/Corr_C14.tex
% ============================================================================

\chapter{Risques, Sécurité, Continuité, Conformité}
\label{chap:maitriser}
\chaptermark{Maîtriser}

\epigraph{%
    « Tout ce qui est susceptible d'aller mal\\
    finira par aller mal. »%
}{Loi dite de Murphy}

\niveauChapitre{Intermédiaire}{%
    Chapitres~\ref{chap:droit} et \ref{chap:piloter}
}

\begin{boxobjectifs}
\begin{itemize}
    \item Distinguer les risques du schéma directeur des risques qu'il traite.
    \item Écrire un risque en trois parties et le coter.
    \item Dire ce que le schéma directeur apporte à la sécurité.
    \item Conduire un bilan d'impact sur l'activité.
    \item Tirer un registre de conformité d'une carte d'applicabilité.
    \item Citer quatre parades à l'échec des projets.
\end{itemize}
\end{boxobjectifs}

\bigskip

% ============================================================================
\section*{Avant de commencer : « quatre jours »}
% ============================================================================

Yaoundé, mars~2026. Un onduleur de la salle des serveurs cède un mardi à 14~heures. Les serveurs s'arrêtent sans
précaution. Au redémarrage, la base des paiements est corrompue. La dernière sauvegarde utilisable date de onze
jours. Il faut la restaurer, puis ressaisir onze jours d'opérations à partir des bordereaux des antennes.

Les paiements du mois partent avec quatre jours de retard. Pour 186~000~bénéficiaires, quatre jours sont longs.
Le ministre de tutelle convoque le directeur général. Le schéma directeur 2021-2025 prévoyait bien un « plan de
secours ». C'était le projet n°~17. Il n'avait jamais démarré, faute de budget. Personne n'avait écrit ce que
coûterait son absence.

\begin{boxalert}
\textbf{Violation typique.} Mettre la continuité et la sécurité en concurrence avec les autres projets lors de
l'arbitrage.

\smallskip
\textbf{Règle~:} les projets qui répondent à un risque que l'organisation ne peut pas porter sont
obligatoires. Ils se financent avant les autres.
\end{boxalert}

% ============================================================================
\section{Deux Sortes de Risques}
\label{sec:c14-deux}
% ============================================================================

Le schéma directeur rencontre le risque deux fois. Il y a les risques \emph{du} schéma directeur~: ce qui peut
empêcher de l'exécuter. Il y a les risques que le schéma directeur \emph{traite}~: ceux qui pèsent sur le
système d'information et que les projets doivent réduire. Les premiers se gèrent dans le pilotage. Les seconds
commandent une partie du portefeuille.

La norme ISO~31000 définit le risque comme l'effet de l'incertitude sur les objectifs et décrit un processus en
cinq temps~: établir le contexte, identifier, analyser, évaluer, traiter~\cite{iso31000}. S'y ajoutent deux
activités continues, communiquer et surveiller.

% ============================================================================
\section{Le Registre des Risques du Schéma Directeur}
\label{sec:c14-registre}
% ============================================================================

Chaque risque s'écrit en trois parties~: une cause, un événement, une conséquence sur un objectif. Il reçoit
une probabilité et un impact notés de 1 à 5. Leur produit donne la criticité. Un propriétaire en répond.

\begin{tableaucours}{L{0.7cm}L{4.6cm}cccL{3.2cm}}{%
    \tableauHeader{N°} &
    \tableauHeader{Risque} &
    \tableauHeader{P} &
    \tableauHeader{I} &
    \tableauHeader{C} &
    \tableauHeader{Traitement}
}
R1 & Subvention non inscrite en loi de finances & 4 & 4 & 16 & Scénario de repli \\
R2 & Reprise des onze fichiers de bénéficiaires défaillante & 4 & 4 & 16 & Fiabilisation avant reprise \\
R3 & Départ des deux agents de la liquidation & 3 & 5 & 15 & Documentation, binômes \\
R4 & Marché du site de secours infructueux & 3 & 4 & 12 & Consultation préalable du marché \\
R5 & Postes du plan de dotation non autorisés & 3 & 4 & 12 & Saisine de la tutelle, assistance externe \\
R6 & Refus du fichier unique par les antennes & 3 & 3 & 9 & Antennes pilotes, accompagnement \\
\end{tableaucours}
\vspace{-0.8em}
\begin{center}{\small\itshape Tableau~14.1. Le registre des risques du schéma directeur de l'ANPP (extrait)}\end{center}

Le seuil d'attention de l'ANPP est fixé à 12. Cinq risques l'atteignent. Chacun a un traitement, un
propriétaire et une date de revue. Quatre traitements sont possibles~: éviter, réduire, transférer, accepter.
Accepter est une décision. Elle se signe.

\begin{boxalert}
\textbf{Violation typique.} Un registre des risques rempli une fois, annexé au document, jamais rouvert.

\smallskip
\textbf{Règle~:} le registre est relu à chaque comité de portefeuille. Un risque dont la note n'a pas bougé en
un an n'a pas été regardé.
\end{boxalert}

% ============================================================================
\section{La Sécurité}
\label{sec:c14-securite}
% ============================================================================

La norme ISO/IEC~27001 demande un système de management de la sécurité de l'information~: une politique, une
appréciation des risques, des mesures, des contrôles, une amélioration~\cite{iso27001}. La loi n°~2010/012
impose de protéger les systèmes et de les faire auditer~\cite{loi2010012,decret20121643}. La lettre circulaire
de la COBAC demande aux banques une politique fondée sur les normes internationales~\cite{cobac2022lc}.

Le schéma directeur ne remplace pas la politique de sécurité. Il lui donne quatre choses.

\begin{enumerate}
    \item \textbf{Une organisation.} Un responsable de la sécurité qui ne dépend pas de celui qu'il contrôle.
          À l'ANPP, il est rattaché au directeur général.
    \item \textbf{Un budget.} Une ligne propre, distincte des projets. L'ANPP y consacre 6~\% de sa dépense
          informatique annuelle à partir de 2028.
    \item \textbf{Une règle de projet.} Chaque fiche comporte une analyse de risque de sécurité et son coût.
    \item \textbf{Un calendrier.} L'audit obligatoire, les tests d'intrusion, les revues d'habilitations sont
          programmés sur cinq ans.
\end{enumerate}

% ============================================================================
\section{La Continuité}
\label{sec:c14-continuite}
% ============================================================================

La norme ISO~22301 traite de la continuité d'activité~\cite{iso22301}. Sa démarche part du métier. Un
\termecle{bilan d'impact sur l'activité} demande à chaque direction combien de temps elle peut travailler sans
son système et combien de données elle peut accepter de perdre.

\begin{boxdef}
\textbf{Durée maximale d'interruption admissible.} Délai au-delà duquel l'arrêt d'une activité cause un
dommage inacceptable. Elle fixe le temps laissé pour redémarrer.

\smallskip
\textbf{Perte de données maximale admissible.} Ancienneté maximale des données que l'on accepte de retrouver
après un sinistre. Elle fixe la fréquence des sauvegardes ou de la réplication.
\end{boxdef}

\begin{tableaucours}{L{3.4cm}L{2.6cm}L{2.6cm}L{3.4cm}}{%
    \tableauHeader{Activité de l'ANPP} &
    \tableauHeader{Interruption admissible} &
    \tableauHeader{Perte admissible} &
    \tableauHeader{Constat de mars 2026}
}
Paiement des prestations & 4 heures & 1 heure & 4 jours, 11 jours de données \\
Encaissement (e-RECETTES) & 8 heures & 1 heure & Aucun secours \\
Liquidation des droits & 2 jours & 1 jour & 4 jours \\
Comptabilité & 5 jours & 1 jour & 4 jours \\
\end{tableaucours}
\vspace{-0.8em}
\begin{center}{\small\itshape Tableau~14.2. Le bilan d'impact sur l'activité de l'ANPP}\end{center}

L'écart entre les besoins et le constat dimensionne le projet de secours. Une interruption admissible de quatre
heures impose un second site prêt à prendre le relais et une réplication des données. Elle exclut une simple
sauvegarde sur bande. Voilà pourquoi le site de secours coûte 620~millions.

Un plan de continuité qui n'a pas été essayé n'existe pas. Le schéma directeur programme un exercice par an et
en fait un indicateur du tableau de bord. Pour une banque, le règlement COBAC~R-2008/01 fait du plan de
continuité une obligation~\cite{cobac200801}.

% ============================================================================
\section{La Conformité}
\label{sec:c14-conformite}
% ============================================================================

La carte d'applicabilité du Chapitre~\ref{chap:droit} devient un \termecle{registre de conformité}. Chaque
exigence y reçoit un responsable, un état et une pièce.

\begin{tableaucours}{L{3.6cm}L{3.2cm}L{2.2cm}L{3.0cm}}{%
    \tableauHeader{Exigence} &
    \tableauHeader{Source} &
    \tableauHeader{État en 2027} &
    \tableauHeader{Projet qui la traite}
}
Audit de sécurité & Loi n°~2010/012 & Non réalisé depuis 2022 & Obligatoire n°~2 \\
Formalités préalables des traitements & Loi n°~2024/017 & Non engagées & Obligatoire n°~3 \\
Engagements et paiements & Loi n°~2018/012 & Conforme & Plan de financement \\
Programmation des marchés & Décret n°~2018/366 & Conforme pour 2027 & Plan d'achat \\
Approbation par le conseil & Loi n°~2017/010 & Conforme & Délibération du 24~juin \\
\end{tableaucours}
\vspace{-0.8em}
\begin{center}{\small\itshape Tableau~14.3. Le registre de conformité du schéma directeur (extrait)}\end{center}

La conformité se construit dans les projets. Le registre impose que la fiche de tout projet touchant des
données personnelles porte une tâche d'analyse et de formalité, avec son délai. La même logique vaut pour la
sécurité et la continuité. On parle de prise en compte \termecle{dès la conception}.

% ============================================================================
\section{Les Projets Eux-Mêmes}
\label{sec:c14-projets}
% ============================================================================

Les projets informatiques échouent d'une manière connue~: ils dépassent leur budget, glissent ou livrent moins
que prévu. Quatre parades relèvent du schéma directeur et non du chef de projet.

\begin{enumerate}
    \item \textbf{La taille.} Aucun lot ne dépasse deux ans ni 800~millions à l'ANPP. Le lot~2 de la refonte,
          à 780~millions, est à la limite.
    \item \textbf{Le jalon d'arrêt.} Chaque projet a un point où le comité décide de poursuivre ou non.
    \item \textbf{La coexistence.} L'ancien système reste en service jusqu'à la recette du nouveau. Son coût
          pendant la période est chiffré.
    \item \textbf{Le regard extérieur.} Les trois projets les plus lourds font l'objet d'une revue
          indépendante chaque année.
\end{enumerate}

\begin{boxprocedure}
\textbf{Intégrer la maîtrise des risques au schéma directeur.}
\begin{enumerate}
    \item Dresser le registre des risques du schéma directeur. Désigner les propriétaires.
    \item Faire le bilan d'impact sur l'activité avec les directions.
    \item En déduire les projets obligatoires de sécurité et de continuité.
    \item Tirer le registre de conformité de la carte d'applicabilité.
    \item Ajouter à la fiche de projet les rubriques sécurité, continuité, données personnelles.
    \item Programmer les audits, les exercices et les revues sur toute la période.
\end{enumerate}
\end{boxprocedure}

\begin{boxlecture}{La maîtrise des risques dans le schéma directeur 2027-2031 de l'ANPP}
\begin{tableaucours}{L{7.6cm}cL{3.4cm}}{%
    \tableauHeader{Exigence} &
    \tableauHeader{Cote} &
    \tableauHeader{Preuve}
}
Un registre des risques existe avec des propriétaires & \OUI{} & Registre du 11 juin 2027 \\
Chaque risque au-dessus du seuil a un traitement & \OUI{} & Cinq traitements \\
Un bilan d'impact sur l'activité a été conduit & \OUI{} & Tableau 14.2 \\
Le secours est dimensionné sur ce bilan & \OUI{} & Fiche du projet de secours \\
Un exercice de continuité annuel est programmé & \OUI{} & Calendrier 2028-2031 \\
Le responsable de la sécurité est indépendant de l'informatique & \OUI{} & Cible d'organisation \\
La sécurité a une ligne budgétaire propre & \OUI{} & 6~\% à partir de 2028 \\
Chaque fiche de projet porte les rubriques de risque & \NON{} & Rubriques absentes de huit fiches \\
Le registre de conformité est tenu & \OUI{} & Tableau 14.3 \\
Le registre des risques a été relu depuis l'approbation & \NEK{} & Aucun procès-verbal fourni \\
\end{tableaucours}

\medskip
\textbf{Décompte.} Huit \OUI{}, un \NON{}, un \NEK{} sur dix.
\hfill\lecture{0{,}80}{0{,}10}{0{,}92}

\medskip
\textbf{Lecture.} Huit \OUI{} sur dix. Le document de 2027 fait ce que celui de 2021 ignorait. Le \NON{} est facile à
lever. Le \NEK{} est de preuve et il est le plus instructif~: un registre bien fait à la date d'approbation ne
prouve rien sur sa vie ensuite. La cote dépendra des procès-verbaux du comité de portefeuille.
\end{boxlecture}

\begin{boxmetier}{Responsable de la sécurité}
Tenir la politique, programmer les audits, viser les fiches de projet.
\end{boxmetier}

\begin{boxmetier}{Responsable de la continuité}
Conduire le bilan d'impact avec les métiers. Organiser l'exercice annuel et en rendre compte.
\end{boxmetier}

\begin{boxmetier}{Directeur de métier}
Dire combien de temps il peut tenir sans son système. Signer sa réponse.
\end{boxmetier}

\begin{boxmetier}{Auditeur}
Demander le dernier exercice de continuité et les procès-verbaux de relecture du registre.
\end{boxmetier}

\begin{boxresume}
\begin{itemize}
    \item Deux sortes de risques~: ceux du schéma directeur, ceux qu'il traite.
    \item Un risque s'écrit~: cause, événement, conséquence. Il a un propriétaire.
    \item Quatre traitements~: éviter, réduire, transférer, accepter.
    \item Le schéma directeur donne à la sécurité une organisation, un budget, une règle de projet, un calendrier.
    \item Le bilan d'impact fixe l'interruption et la perte de données admissibles.
    \item Un plan de continuité non essayé n'existe pas.
    \item La carte d'applicabilité devient un registre de conformité.
\end{itemize}
\end{boxresume}

\begin{boxinfo}
\textbf{Points de contrôle.} Vous avez compris ce chapitre si vous savez~:
\begin{itemize}
    \item écrire et coter six risques~;
    \item conduire un bilan d'impact pour une activité~;
    \item dresser un registre de conformité~;
    \item ajouter les rubriques de risque à une fiche de projet.
\end{itemize}
\end{boxinfo}

\begin{boxexo}[title={\ding{45}\quad Exercice~14.1. Coter \niveaubase}]
Écrivez en trois parties puis cotez le risque suivant pour LiZbiZ~: son client principal, qui pèse 38~\% du chiffre d'affaires, exige un suivi des colis en temps réel que le système actuel ne fournit pas.
\end{boxexo}

\begin{boxexo}[title={\ding{45}\quad Exercice~14.2. Le bilan d'impact \niveaumoyen}]
Pour l'activité « paiement des prestations », expliquez pourquoi une interruption admissible de quatre heures et une perte admissible d'une heure excluent la sauvegarde quotidienne. Quelles solutions restent~?
\end{boxexo}

\begin{boxexo}[title={\ding{45}\quad Exercice~14.3. Accepter \niveaumoyen}]
Le comité de l'ANPP veut accepter le risque~R6. Rédigez la décision d'acceptation~: motif, limite, signataire, date de revue.
\end{boxexo}

\begin{boxexo}[title={\ding{45}\quad Exercice~14.4. La banque \niveauexpert}]
Dressez le registre de conformité du schéma directeur de la BANQUE SANAGA pour la sécurité et la continuité, à partir des textes du Chapitre~\ref{chap:droit}. Cinq lignes au moins, avec la pièce attendue par le superviseur.
\end{boxexo}

\bigskip

\begin{boxinfo}
\textbf{Pour aller plus loin.}
\begin{itemize}
    \item ISO~31000~\cite{iso31000}.
    \item ISO/IEC~27001~\cite{iso27001} et ISO~22301~\cite{iso22301}.
    \item Le règlement COBAC~R-2008/01~\cite{cobac200801}.
\end{itemize}
\end{boxinfo}

\begin{boxinfo}
\textbf{Transition.} La méthode est complète, de la commande à l'exécution. La dernière partie regarde ce qui la bouscule aujourd'hui, puis la manière de la contrôler. Le Chapitre~\ref{chap:donnees} commence par les données.
\end{boxinfo}


    \part{Défis et recherche}
    \label{part:defis}

        % ============================================================================
% Fichier  : partie6/P06_C15_donnees_ia.tex
% Titre    : Données, intelligence artificielle, plateformes, sobriété
% Auteur   : MINKA MI NGUIDJOÏ Thierry Emmanuel
% Version  : 1.0 -- novembre 2026
% Corrigés : corriges/Corr_C15.tex
% ============================================================================

\chapter{Données, intelligence artificielle, plateformes, sobriété}
\label{chap:donnees}
\chaptermark{Défis}

\epigraph{%
    « Une donnée sans propriétaire\\
    n'est tenue par personne. »%
}{Principe du cours}

\niveauChapitre{Avancé}{%
    Chapitres~\ref{chap:strategie}, \ref{chap:doctrines} et \ref{chap:cible}
}

\begin{boxobjectifs}
\begin{itemize}
    \item Répondre aux cinq questions de gouvernance d'une donnée.
    \item Juger un cas d'usage d'intelligence artificielle sur cinq conditions.
    \item Lister les échanges d'une organisation avec l'extérieur et leurs dépendances.
    \item Formuler une doctrine d'hébergement par classe de données.
    \item Chiffrer un levier de sobriété.
    \item Dire ce que ces forces ajoutent au schéma directeur.
\end{itemize}
\end{boxobjectifs}

\bigskip

% ============================================================================
\section*{Avant de commencer : « mettons de l'intelligence artificielle »}
% ============================================================================

Yaoundé, mars~2026. Un administrateur revient d'un colloque. Au conseil, il demande pourquoi le schéma directeur
voté neuf mois plus tôt « ne contient pas d'intelligence artificielle ». Il veut un projet avant la fin de
l'année. M.~OWONA transmet à la cellule de portefeuille.

La cellule ne répond ni oui ni non. Elle demande quelle décision de l'Agence on veut améliorer. La direction
financière a une réponse~: repérer les paiements indus. La cellule demande alors sur quelles données. La réponse
est moins assurée. Le fichier compte encore 7,4~\% de doublons probables.

\begin{boxalert}
\textbf{Violation typique.} Ajouter une technique au schéma directeur parce qu'elle est à la mode.

\smallskip
\textbf{Règle~:} une technique entre au portefeuille par un cas d'usage, instruit comme toute demande
nouvelle.
\end{boxalert}

% ============================================================================
\section{Quatre forces qui bousculent la méthode}
\label{sec:c15-forces}
% ============================================================================

La méthode des chapitres précédents a été conçue pour des applications et des projets. Quatre forces
l'obligent à s'élargir. Les données deviennent un actif à part. L'intelligence artificielle change ce qu'un
système peut faire et ce qu'il peut rater. Les plateformes déplacent la frontière de l'organisation. La
sobriété impose de compter autre chose que des francs.

Aucune ne rend la méthode caduque. Chacune ajoute un chapitre au schéma directeur.

% ============================================================================
\section{Les données}
\label{sec:c15-donnees}
% ============================================================================

Le Cadre stratégique de l'Union africaine en matière de données, adopté en février~2022, demande aux États de
traiter les données comme une ressource à gouverner et à faire circuler~\cite{ua2022donnees}. Pour une
organisation, gouverner ses données revient à répondre à cinq questions pour chaque donnée importante.

\begin{tableaucours}{L{3.2cm}L{4.4cm}L{4.4cm}}{%
    \tableauHeader{Question} &
    \tableauHeader{Réponse attendue} &
    \tableauHeader{À l'ANPP, pour le bénéficiaire}
}
À qui appartient-elle~? & Un propriétaire dans le métier & Directeur des prestations \\
Où est sa référence~? & Un seul système maître & Référentiel des bénéficiaires \\
Que vaut-elle~? & Une mesure de qualité & Taux de doublons, taux d'identifiants vérifiés \\
Qui y accède~? & Des règles d'habilitation & Par antenne, par fonction \\
Combien de temps~? & Une durée et un sort final & Durée fixée par le registre des traitements \\
\end{tableaucours}
\vspace{-0.8em}
\begin{center}{\small\itshape Tableau~15.1. Cinq questions de gouvernance pour une donnée}\end{center}

Le diagnostic a trouvé onze fichiers de bénéficiaires à l'ANPP. Le rapprochement d'un échantillon de
5~000~dossiers a donné 7,4~\% de doublons probables. Sur 186~000~bénéficiaires, cela représente près de
13~800~dossiers à examiner. Le projet de référentiel commence donc par un chantier de fiabilisation, qui n'est
pas un chantier informatique.

\begin{boxdef}
\textbf{Donnée de référence.} Donnée partagée par plusieurs processus et qui décrit un objet stable de
l'activité~: une personne, un employeur, un barème. Elle a un propriétaire, un système maître et une règle de
mise à jour.
\end{boxdef}

Le schéma directeur comporte désormais un volet données~: la liste des données de référence, leurs
propriétaires, les indicateurs de qualité, la cellule qui les suit. La loi n°~2024/017 y ajoute ses propres
exigences pour les données personnelles~\cite{loi2024017}.

% ============================================================================
\section{L'intelligence artificielle}
\label{sec:c15-ia}
% ============================================================================

La Stratégie continentale de l'Union africaine sur l'intelligence artificielle date de
juillet~2024~\cite{ua2024ia}. Le Cameroun a présenté sa stratégie nationale en juillet~2025~\cite{snia2025}. La
norme ISO/IEC~42001 décrit un système de management de l'intelligence artificielle~\cite{iso42001}.

Pour le schéma directeur, l'intelligence artificielle n'est pas un projet. C'est une famille de techniques qui
peut servir un objectif. La question reste celle du Chapitre~\ref{chap:cible}~: quelle capacité~? Un cas
d'usage se juge sur cinq conditions avant d'entrer au portefeuille.

\begin{tableaucours}{L{3.0cm}L{9.0cm}}{%
    \tableauHeader{Condition} &
    \tableauHeader{Question}
}
Objectif & Quel objectif du plan stratégique ce cas sert-il, avec quelle mesure~? \\
Données & Les données existent-elles, en quantité et en qualité suffisantes et a-t-on le droit de les employer~? \\
Décision & Qui décide à la fin~: la machine ou un agent~? Qui répond d'une erreur~? \\
Preuve & Peut-on expliquer un résultat à la personne concernée~? \\
Exploitation & Qui surveillera la dérive du modèle, avec quel budget~? \\
\end{tableaucours}
\vspace{-0.8em}
\begin{center}{\small\itshape Tableau~15.2. Cinq conditions pour un cas d'usage d'intelligence artificielle}\end{center}

\begin{boxscenario}{ANPP. Le cas des paiements indus, révision de 2026}
La direction financière propose un modèle qui signale les dossiers à risque de paiement indu~: bénéficiaire
décédé, double immatriculation, droit éteint. Le comité de portefeuille applique les cinq conditions.

L'objectif est clair~: ne payer que les ayants droit. Les données ne sont pas prêtes~: le référentiel ne sera
fiabilisé qu'à la fin de 2026. La décision est fixée~: le modèle signale, un agent instruit, aucune suspension
automatique. La preuve est exigée~: chaque signalement indique ses motifs. L'exploitation est confiée à la
cellule des données.

Le projet entre au portefeuille pour 2027, après le référentiel. Il coûte 95~millions. Il décale un projet de
soutien, selon la règle du Chapitre~\ref{chap:piloter}.
\end{boxscenario}

\begin{boxalert}
\textbf{Violation typique.} Inscrire au schéma directeur « un projet d'intelligence artificielle » sans dire
pour quoi faire.

\smallskip
\textbf{Règle~:} un cas d'usage se nomme par la décision qu'il aide à prendre. Il n'entre au portefeuille
qu'après le projet qui lui fournit des données fiables.
\end{boxalert}

Les assistants conversationnels du commerce posent une autre question. Les agents les emploient déjà, sans
attendre le schéma directeur. C'est une forme nouvelle d'informatique hors contrôle. Une règle d'usage, courte,
dit quelles données ne doivent jamais y être saisies. À l'ANPP, aucune donnée de bénéficiaire.

% ============================================================================
\section{Les plateformes}
\label{sec:c15-plateformes}
% ============================================================================

Le Chapitre~\ref{chap:strategie} a présenté les marchés bifaces~\cite{rochettirole2003} et la logique de
plateforme~\cite{parker2016}. Une plateforme crée de la valeur en mettant en relation des tiers. Elle oblige le
schéma directeur à regarder hors des murs.

Pour LiZbiZ, la question est offensive. L'entreprise possède dix ans de données de tournées et un réseau de
quarante correspondants ruraux. Elle peut rester un transporteur, ou ouvrir son réseau à d'autres expéditeurs
et à d'autres livreurs. Dans le second cas, son système d'information cesse d'être un outil interne. Il devient
le produit.

Pour une administration, la question est celle de l'\termecle{interopérabilité}. L'ANPP a besoin de l'état
civil pour connaître les décès, des employeurs pour les carrières, des opérateurs de paiement pour verser. Les
textes qui régissent l'ANTIC lui confient de favoriser l'interopérabilité entre les systèmes
d'information~\cite{decret2012180,decret2019150}. Le schéma directeur liste donc les échanges attendus, leurs partenaires et leur
état.

\begin{tableaucours}{L{3.2cm}L{3.2cm}L{2.6cm}L{3.0cm}}{%
    \tableauHeader{Échange} &
    \tableauHeader{Partenaire} &
    \tableauHeader{État en 2025} &
    \tableauHeader{Dépend de l'ANPP~?}
}
Décès & État civil & Courrier mensuel & Non \\
Carrières & Employeurs & Bordereaux papier & En partie \\
Paiements & Banques, opérateurs de paiement & Fichiers déposés & Oui \\
Identité & Fichier national d'identité & Aucun & Non \\
\end{tableaucours}
\vspace{-0.8em}
\begin{center}{\small\itshape Tableau~15.3. Les échanges de l'ANPP avec l'extérieur}\end{center}

Deux échanges sur quatre ne dépendent pas de l'Agence. Le schéma directeur les inscrit comme des hypothèses,
avec une solution d'attente. Un projet dont le succès dépend d'un tiers qui ne s'est engagé à rien est un
risque à porter au registre.

% ============================================================================
\section{Le nuage et la maîtrise}
\label{sec:c15-nuage}
% ============================================================================

Louer des capacités chez un prestataire réduit l'investissement et accélère les projets. La dépendance grandit
en proportion. Trois questions reviennent dans tout schéma directeur.

\begin{enumerate}
    \item \textbf{Où sont les données~?} La loi n°~2024/017 encadre les transferts hors du Cameroun.
    \item \textbf{Qui peut y accéder~?} Le droit du pays du prestataire s'applique aussi.
    \item \textbf{Comment en sortir~?} La réversibilité, vue au Chapitre~\ref{chap:financer}.
\end{enumerate}

Le schéma directeur fixe une doctrine par classe de données. L'ANPP en retient trois~: les données de
bénéficiaires restent sur des sites situés au Cameroun, les données de gestion peuvent être hébergées dans la
sous-région, les données publiques n'ont pas de contrainte. La doctrine est un principe directeur. Elle a des
conséquences de coût, que le plan de financement porte.

% ============================================================================
\section{La sobriété}
\label{sec:c15-sobriete}
% ============================================================================

Un système d'information consomme de l'énergie, des matériels et de l'attention. Le principe « retirer avant
d'ajouter », adopté au Chapitre~\ref{chap:cible}, en est la traduction. Trois leviers relèvent du schéma
directeur.

\begin{enumerate}
    \item \textbf{La durée de vie.} Allonger d'un an l'usage des postes de travail réduit les achats. L'ANPP
          passe de quatre à six ans.
    \item \textbf{Le retrait.} Une application retirée libère des serveurs, des licences, de l'énergie.
    \item \textbf{L'énergie.} Dans un pays où le réseau électrique connaît des coupures, la salle des serveurs
          tourne une partie du temps sur groupe électrogène. La mesure de sa consommation est un indicateur
          financier autant qu'écologique.
\end{enumerate}

Avec 380~postes de travail à 450~000~FCFA, le renouvellement tous les quatre ans coûte 42,75~millions par an.
Tous les six ans, il coûte 28,5~millions. L'économie annuelle est de 14,25~millions.

% ============================================================================
\section{Ce qui change dans le document}
\label{sec:c15-document}
% ============================================================================

\begin{tableaucours}{L{3.0cm}L{9.0cm}}{%
    \tableauHeader{Force} &
    \tableauHeader{Ajout au schéma directeur}
}
Données & Liste des données de référence, propriétaires, indicateurs de qualité, cellule des données \\
Intelligence artificielle & Grille des cinq conditions, règle d'usage des assistants, registre des cas d'usage \\
Plateformes & Tableau des échanges avec l'extérieur, hypothèses et solutions d'attente \\
Nuage & Doctrine d'hébergement par classe de données \\
Sobriété & Durées de vie, plan de retrait, indicateur de consommation \\
\end{tableaucours}
\vspace{-0.8em}
\begin{center}{\small\itshape Tableau~15.4. Cinq ajouts au schéma directeur}\end{center}

\begin{boxlecture}{Les défis nouveaux dans le schéma directeur de l'ANPP, révision de 2026}
\begin{tableaucours}{L{7.6cm}cL{3.4cm}}{%
    \tableauHeader{Exigence} &
    \tableauHeader{Cote} &
    \tableauHeader{Preuve}
}
Les données de référence sont listées avec leur propriétaire & \OUI{} & Quatre données, quatre propriétaires \\
Un indicateur de qualité est mesuré pour chacune & \NON{} & Mesure pour les bénéficiaires seuls \\
Chaque cas d'usage d'IA passe la grille des cinq conditions & \OUI{} & Fiche des paiements indus \\
Une règle d'usage des assistants est diffusée & \OUI{} & Note du 3 mars 2026 \\
Les échanges avec l'extérieur sont listés avec leur état & \OUI{} & Tableau 15.3 \\
Chaque échange dépendant d'un tiers a une solution d'attente & \NON{} & Rien pour le fichier d'identité \\
Une doctrine d'hébergement par classe de données existe & \OUI{} & Principe directeur n° 7 \\
La consommation de la salle des serveurs est mesurée & \NEK{} & Compteur posé, aucun relevé transmis \\
\end{tableaucours}

\medskip
\textbf{Décompte.} Cinq \OUI{}, deux \NON{}, un \NEK{} sur huit.
\hfill\lecture{0{,}625}{0{,}125}{1{,}30}

\medskip
\textbf{Lecture.} Cinq \OUI{} sur huit. L'Agence a pris les sujets dans le bon ordre~: les données avant l'intelligence
artificielle. Les deux \NON{} sont des travaux à finir. Le \NEK{} est de preuve.
\end{boxlecture}

\begin{boxmetier}{Responsable des données}
Tenir la liste des données de référence. Mesurer leur qualité. Refuser un modèle sur des données non fiabilisées.
\end{boxmetier}

\begin{boxmetier}{Directeur de métier}
Rester propriétaire de ses données et des décisions prises avec l'aide d'un modèle.
\end{boxmetier}

\begin{boxmetier}{Délégué à la protection des données}
Viser chaque cas d'usage et la doctrine d'hébergement.
\end{boxmetier}

\begin{boxmetier}{Consultant}
Ramener toute demande de technique à une capacité et à une mesure.
\end{boxmetier}

\begin{boxresume}
\begin{itemize}
    \item Quatre forces élargissent la méthode~: données, intelligence artificielle, plateformes, sobriété.
    \item Une donnée de référence a un propriétaire, un système maître, une mesure de qualité.
    \item Un cas d'usage d'intelligence artificielle se juge sur cinq conditions~: objectif, données, décision, preuve, exploitation.
    \item Les données viennent avant les modèles.
    \item Le schéma directeur liste les échanges avec l'extérieur et les traite comme des hypothèses quand ils dépendent d'un tiers.
    \item La doctrine d'hébergement se fixe par classe de données.
    \item Retirer avant d'ajouter.
\end{itemize}
\end{boxresume}

\begin{boxinfo}
\textbf{Points de contrôle.} Vous avez compris ce chapitre si vous savez~:
\begin{itemize}
    \item appliquer les cinq questions à une donnée~;
    \item instruire un cas d'usage d'intelligence artificielle~;
    \item dresser le tableau des échanges d'une organisation~;
    \item chiffrer l'effet d'un allongement de durée de vie.
\end{itemize}
\end{boxinfo}

\begin{boxexo}[title={\ding{45}\quad Exercice~15.1. Les doublons \niveaubase}]
Vérifiez le calcul des dossiers à examiner à l'ANPP. Si un agent en examine 25 par jour, combien de jours de travail faut-il~? Combien d'agents pour finir en six mois de 22~jours ouvrés~?
\end{boxexo}

\begin{boxexo}[title={\ding{45}\quad Exercice~15.2. Cinq conditions \niveaumoyen}]
La BANQUE SANAGA veut un modèle qui note les demandes de crédit des particuliers. Appliquez les cinq conditions. Dites laquelle vous paraît la plus difficile à remplir.
\end{boxexo}

\begin{boxexo}[title={\ding{45}\quad Exercice~15.3. La sobriété \niveaumoyen}]
Refaites le calcul du renouvellement des postes pour 520~postes à 480~000~FCFA, en passant de trois à cinq ans. Citez un coût que cet allongement fait apparaître.
\end{boxexo}

\begin{boxexo}[title={\ding{45}\quad Exercice~15.4. LiZbiZ plateforme \niveauexpert}]
Rédigez la note de deux pages au conseil de LiZbiZ sur l'ouverture de son réseau à d'autres expéditeurs. Appuyez-vous sur la théorie des marchés bifaces~: qui sont les deux faces, laquelle subventionner, quel rôle pour les données de tournées, quels risques.
\end{boxexo}

\begin{boxexo}[title={\ding{45}\quad Exercice~15.5. Question de recherche \niveauexpert}]
Formulez une question de recherche sur la qualité des données de référence dans les administrations d'Afrique centrale. Proposez une méthode de mesure et dites quelles données il faudrait collecter.
\end{boxexo}

\bigskip

\begin{boxinfo}
\textbf{Pour aller plus loin.}
\begin{itemize}
    \item Le Cadre de l'Union africaine sur les données~\cite{ua2022donnees}.
    \item La Stratégie continentale sur l'intelligence artificielle~\cite{ua2024ia}.
    \item ISO/IEC~42001~\cite{iso42001}.
    \item Parker, Van Alstyne et Choudary~\cite{parker2016}.
\end{itemize}
\end{boxinfo}

\begin{boxinfo}
\textbf{Transition.} Le schéma directeur est écrit, exécuté, élargi. Reste une question~: comment savoir s'il est bon~? Le Chapitre~\ref{chap:audit} y répond.
\end{boxinfo}

        % ============================================================================
% Fichier  : partie6/P06_C16_auditer.tex
% Titre    : Auditer un schéma directeur, mesurer l'alignement
% Auteur   : MINKA MI NGUIDJOÏ Thierry Emmanuel
% Version  : 1.0 -- novembre 2026
% Corrigés : corriges/Corr_C16.tex
% ============================================================================

\chapter{Auditer un schéma directeur, mesurer l'alignement}
\label{chap:audit}
\chaptermark{Auditer}

\epigraph{%
    « Ce qui est affirmé sans preuve\\
    peut être nié sans preuve. »%
}{Adage latin, \textit{quod gratis asseritur, gratis negatur}}

\niveauChapitre{Avancé}{%
    Tous les chapitres précédents
}

\begin{boxobjectifs}
\begin{itemize}
    \item Distinguer les trois moments d'audit d'un schéma directeur.
    \item Écrire une exigence vérifiable par une pièce.
    \item Coter OUI, NON ou NEK et qualifier un NEK.
    \item Calculer et lire ensemble $\mu$, $U$ et $H$.
    \item Lire des résultats par domaine.
    \item Calculer un indice d'alignement prouvé.
    \item Structurer un rapport d'audit.
\end{itemize}
\end{boxobjectifs}

\bigskip

% ============================================================================
\section*{Avant de commencer : « donnez-moi une note sur vingt »}
% ============================================================================

Yaoundé, vendredi 20~juin~2025. Six jours avant le conseil d'administration. M.~TCHUENTE, son président, reçoit
l'auditeur interne. Il a lu la synthèse de quatre pages. Il veut savoir si le document est bon. « Donnez-moi une
note sur vingt. »

L'auditeur ne donne pas de note sur vingt. Il pose trois nombres sur la table. Trente-sept exigences sur
quarante-huit sont prouvées. Sept ne sont pas satisfaites, dont trois sur le chiffrage. Quatre n'ont pas pu
être vérifiées. « Je peux vous dire ce qui est établi, ce qui manque et ce que j'ignore. Une note unique
mélangerait les trois. »

\begin{boxalert}
\textbf{Violation typique.} Résumer un audit par une note unique.

\smallskip
\textbf{Règle~:} un audit rend trois informations~: ce qui est conforme, ce qui ne l'est pas, ce qui n'a pas
pu être établi. Les fondre en un nombre détruit la troisième.
\end{boxalert}

% ============================================================================
\section{Trois moments, trois questions}
\label{sec:c16-moments}
% ============================================================================

Un schéma directeur s'audite à trois moments. Chacun pose une question différente et s'appuie sur des pièces
différentes.

\begin{tableaucours}{L{2.6cm}L{4.4cm}L{5.0cm}}{%
    \tableauHeader{Moment} &
    \tableauHeader{Question} &
    \tableauHeader{Pièces}
}
Avant l'approbation & Le document est-il exécutable~? & Document, note de cadrage, diagnostic, fiches, plan de financement \\
Pendant l'exécution & Est-il exécuté et piloté~? & Procès-verbaux, tableau de bord, registres, marchés \\
À l'échéance & A-t-il produit ce qu'il promettait~? & Mesures de départ et d'arrivée des bénéfices, comptes \\
\end{tableaucours}
\vspace{-0.8em}
\begin{center}{\small\itshape Tableau~16.1. Les trois audits d'un schéma directeur}\end{center}

Le prologue a pratiqué le troisième sur le document de 2019. Les chapitres de méthode ont pratiqué le premier
sur chaque livrable de 2025. Ce chapitre donne la méthode complète.

% ============================================================================
\section{L'objet, le critère, la preuve}
\label{sec:c16-objet}
% ============================================================================

L'Auditologie, que l'auteur développe dans un autre ouvrage~\cite{minka2026auditologie}, pose qu'un audit
compare un objet à un critère au moyen d'une preuve et qu'il rend compte de ce qu'il n'a pas pu établir. Trois
exigences en découlent.

\begin{enumerate}
    \item \textbf{L'objet est délimité.} On audite un document daté, une période, un périmètre.
    \item \textbf{Le critère préexiste.} Il vient d'un texte, d'une norme, d'un engagement de l'organisation
          ou d'une règle de l'art reconnue. L'auditeur ne l'invente pas en cours de route.
    \item \textbf{La preuve est une pièce.} Un document, un enregistrement, une observation consignée. Une
          déclaration orale n'est une preuve que de ce qui a été déclaré.
\end{enumerate}

Chaque exigence reçoit l'une de trois cotes. \OUI{}~: la preuve montre que l'exigence est satisfaite.
\NON{}~: la preuve montre qu'elle ne l'est pas. \NEK{}~: l'auditeur ne peut pas conclure. Le \NEK{} est de
preuve quand la pièce manque. Il est de critère quand l'exigence ne s'applique pas ou ne peut pas être
formulée pour cet objet.

\begin{boxalert}
\textbf{Violation typique.} Coter \NON{} faute de pièce, ou \OUI{} sur la foi d'une déclaration.

\smallskip
\textbf{Règle~:} l'absence de preuve n'est pas la preuve de l'absence. Sans pièce, la cote est \NEK{}.
\end{boxalert}

% ============================================================================
\section{Le référentiel d'exigences}
\label{sec:c16-referentiel}
% ============================================================================

Le référentiel d'audit d'un schéma directeur rassemble des exigences de cinq sources. L'Annexe~D donne celui de
la mission ANPP. Il compte 48~exigences réparties en huit domaines, un par chapitre de méthode.

\begin{tableaucours}{L{3.4cm}L{5.4cm}c}{%
    \tableauHeader{Source} &
    \tableauHeader{Exemple d'exigence} &
    \tableauHeader{Exigences}
}
Textes & Le document est approuvé par l'organe compétent & 9 \\
Normes de gouvernance & L'organe de gouvernance surveille l'exécution & 8 \\
Engagements de l'organisation & Chaque projet sert un objectif du plan stratégique & 7 \\
Règles de l'art & Chaque projet affiche son coût annuel après mise en service & 19 \\
Cohérence interne & Le total des fiches égale le total du plan de financement & 5 \\
\end{tableaucours}
\vspace{-0.8em}
\begin{center}{\small\itshape Tableau~16.2. Les cinq sources du référentiel de l'ANPP}\end{center}

Une exigence bien écrite est vérifiable par une pièce. « Le schéma directeur est aligné sur la stratégie » ne
se vérifie pas. « Chaque projet retenu cite un objectif du plan stratégique » se vérifie fiche par fiche.

% ============================================================================
\section{Trois mesures}
\label{sec:c16-mesures}
% ============================================================================

Les cotes se résument par trois nombres. Soit $n$ le nombre d'exigences, $n_O$, $n_N$ et $n_K$ les effectifs
des trois cotes.

\begin{equation}
    \mu = \frac{n_O}{n} \qquad\qquad U = \frac{n_K}{n} \qquad\qquad
    H = -\sum_{c \in \{O,N,K\}} \frac{n_c}{n} \log_2 \frac{n_c}{n}
\end{equation}

$\mu$ est la part des exigences dont la satisfaction est prouvée. $U$ est la part de ce que l'auditeur ignore.
$H$ est l'entropie de Shannon de la répartition~\cite{shannon1948}. Elle vaut 0 quand toutes les cotes sont
identiques et atteint son maximum, $\log_2 3 \approx 1{,}58$, quand les trois cotes sont à égalité. Un terme
dont l'effectif est nul ne compte pas.

Les trois nombres se lisent ensemble. Un $\mu$ de 0,60 avec un $U$ de 0,05 décrit un objet connu et
moyennement conforme. Le même $\mu$ avec un $U$ de 0,35 décrit un objet mal connu~: il peut être très bon ou
médiocre. La conclusion n'est pas la même. Dans le second cas, la première recommandation est d'obtenir les
pièces.

\begin{boxscenario}{Audit du schéma directeur 2025-2029 de l'ANPP avant approbation, 20 juin 2025}
Le directeur général demande à l'auditeur interne de l'Agence un avis avant le conseil du 26~juin. L'auditeur
applique les 48~exigences de l'Annexe~D. Il obtient 37~\OUI{}, 7~\NON{} et 4~\NEK{}.

\[ \mu = \frac{37}{48} = 0{,}771 \qquad U = \frac{4}{48} = 0{,}083 \]
\[ H = -(0{,}771 \log_2 0{,}771 + 0{,}146 \log_2 0{,}146 + 0{,}083 \log_2 0{,}083) = 0{,}99 \]

Le document est connu à 92~\%. Plus des trois quarts des exigences sont prouvées. Les sept \NON{} se
concentrent dans deux domaines~: le chiffrage, avec trois exigences et la maîtrise des risques dans les
fiches, avec deux. Les quatre \NEK{} sont tous de preuve.
\end{boxscenario}

% ============================================================================
\section{Lire par domaine}
\label{sec:c16-domaines}
% ============================================================================

Une moyenne cache les écarts. La lecture se fait donc aussi par domaine.

\begin{tableaucours}{L{4.0cm}ccccc}{%
    \tableauHeader{Domaine} &
    \tableauHeader{Exigences} &
    \tableauHeader{\OUI{}} &
    \tableauHeader{\NON{}} &
    \tableauHeader{\NEK{}} &
    \tableauHeader{$\mu$}
}
Cadrage & 6 & 5 & 0 & 1 & 0,83 \\
Diagnostic & 6 & 5 & 1 & 0 & 0,83 \\
Cible & 6 & 6 & 0 & 0 & 1,00 \\
Portefeuille & 6 & 5 & 1 & 0 & 0,83 \\
Chiffrage et financement & 7 & 3 & 3 & 1 & 0,43 \\
Pilotage & 6 & 5 & 0 & 1 & 0,83 \\
Maîtrise des risques & 6 & 4 & 2 & 0 & 0,67 \\
Cadre juridique & 5 & 4 & 0 & 1 & 0,80 \\
Ensemble & 48 & 37 & 7 & 4 & 0,77 \\
\end{tableaucours}
\vspace{-0.8em}
\begin{center}{\small\itshape Tableau~16.3. Les résultats de l'audit par domaine}\end{center}

Le chiffrage est le point faible. C'est aussi là que le document de 2019 avait failli. L'avis de l'auditeur
tient en trois phrases~: le document peut être soumis au conseil, sous réserve que le total des fiches et celui du plan
de financement soient mis en accord et que le conseil soit informé de la hausse du fonctionnement. Les quatre pièces manquantes
doivent être produites avant la première revue.

% ============================================================================
\section{Mesurer l'alignement}
\label{sec:c16-alignement}
% ============================================================================

Le Chapitre~\ref{chap:alignement} a présenté le modèle de Henderson et Venkatraman~\cite{henderson1993} et la
mesure de Luftman~\cite{luftman2000}. Luftman note six critères sur cinq niveaux, du processus initial au
processus optimisé. Sa mesure repose sur l'opinion de dirigeants interrogés. Elle capte une perception.

L'audit par exigences capte autre chose~: des faits prouvés. Les deux se complètent. On peut dériver du
référentiel un \termecle{indice d'alignement prouvé} en ne gardant que les exigences qui relient le système
d'information à la stratégie.

\begin{tableaucours}{L{8.4cm}cL{2.6cm}}{%
    \tableauHeader{Exigence d'alignement} &
    \tableauHeader{Cote} &
    \tableauHeader{Preuve}
}
Chaque capacité se rattache à un objectif du plan stratégique & \OUI{} & Dossier de cible \\
Chaque projet retenu cite un objectif & \OUI{} & 22 fiches \\
Chaque objectif est servi par un projet au moins & \OUI{} & Table de correspondance \\
Chaque projet a un responsable dans le métier & \OUI{} & 22 fiches \\
Le comité qui arbitre est présidé par la direction générale & \OUI{} & Note de cadrage \\
Chaque bénéfice a un propriétaire et une mesure de départ & \NON{} & Deux bénéfices sans mesure \\
Le tableau de bord contient un indicateur de résultat métier & \OUI{} & Tableau de bord \\
La révision annuelle précède la préparation du budget & \NEK{} & Calendrier non arrêté \\
\end{tableaucours}
\vspace{-0.8em}
\begin{center}{\small\itshape Tableau~16.4. L'indice d'alignement prouvé de l'ANPP, juin 2025}\end{center}

Six \OUI{} sur huit~: l'indice vaut 0,75. Appliquée au document de 2019, la même grille donnait un \OUI{} sur
huit, soit 0,125. La mesure ne dit pas que le système d'information de l'ANPP est aligné. Elle dit que le
dispositif qui doit produire l'alignement existe et qu'il est prouvé.

% ============================================================================
\section{Le rapport}
\label{sec:c16-rapport}
% ============================================================================

Le rapport d'audit d'un schéma directeur tient en dix pages. Il suit un ordre constant.

\begin{enumerate}
    \item L'objet, le périmètre, la date, le référentiel employé.
    \item Les trois mesures, pour l'ensemble et par domaine.
    \item Les \NON{}, chacun avec le fait, la pièce, la conséquence et une recommandation.
    \item Les \NEK{}, chacun avec la pièce attendue et la personne qui doit la produire.
    \item L'avis.
    \item En annexe, la table complète des cotes.
\end{enumerate}

Les \NEK{} sont publiés. C'est la marque de la méthode. Un rapport qui les tait laisse croire que l'auditeur a
tout vu.

\begin{boxprocedure}
\textbf{Auditer un schéma directeur.}
\begin{enumerate}
    \item Délimiter l'objet et le moment.
    \item Arrêter le référentiel avec le commanditaire avant de voir les pièces.
    \item Demander les pièces par écrit. Tenir un registre de ce qui est reçu.
    \item Coter chaque exigence. Écrire la référence de la pièce.
    \item Calculer $\mu$, $U$ et $H$, pour l'ensemble et par domaine.
    \item Faire valider les faits par l'audité. Les cotes ne se négocient pas.
    \item Rédiger le rapport. Publier les \NEK{}.
\end{enumerate}
\end{boxprocedure}

% ============================================================================
\section{Pistes de recherche}
\label{sec:c16-recherche}
% ============================================================================

Le cours ouvre plus de questions qu'il n'en ferme. Quatre d'entre elles peuvent nourrir un mémoire.

\begin{enumerate}
    \item \textbf{Le devenir des schémas directeurs.} Combien de documents adoptés par des organismes publics
          de la sous-région sont exécutés à plus de moitié~? Aucune mesure publiée n'existe. La grille du
          prologue permet de la produire.
    \item \textbf{Perception et preuve.} L'indice d'alignement prouvé et la maturité de Luftman
          classent-ils les organisations dans le même ordre~?
    \item \textbf{Le paradoxe de la productivité.} La question de Solow~\cite{solow1987} et de
          Brynjolfsson~\cite{brynjolfsson1993} a été posée pour des économies industrielles. Elle reste
          ouverte pour les administrations d'Afrique centrale.
    \item \textbf{L'informatique hors contrôle.} Quelle part de l'activité réelle des administrations repose
          sur des outils non recensés~? Seize applications sur 47 à l'ANPP est un cas, non une statistique.
\end{enumerate}

\begin{boxlecture}{L'audit du 20 juin 2025 jugé selon sa propre méthode}
\begin{tableaucours}{L{7.6cm}cL{3.4cm}}{%
    \tableauHeader{Exigence} &
    \tableauHeader{Cote} &
    \tableauHeader{Preuve}
}
L'objet et la date sont délimités & \OUI{} & Version du 13 juin 2025 \\
Le référentiel a été arrêté avant l'examen des pièces & \OUI{} & Courriel du 4 juin \\
Chaque cote renvoie à une pièce & \OUI{} & Table des cotes \\
Les mesures sont données par domaine & \OUI{} & Tableau 16.3 \\
Les NEK sont publiés avec la pièce attendue & \OUI{} & Section 4 du rapport \\
Les faits ont été validés par l'audité & \OUI{} & Réunion du 19 juin \\
L'auditeur est indépendant de l'objet audité & \NON{} & Il a siégé au groupe de travail finances \\
L'avis distingue les réserves des recommandations & \OUI{} & Section 5 \\
\end{tableaucours}

\medskip
\textbf{Décompte.} Sept \OUI{}, un \NON{}, zéro \NEK{} sur huit.
\hfill\lecture{0{,}875}{0}{0{,}54}

\medskip
\textbf{Lecture.} Sept \OUI{} sur huit. Le \NON{} est sérieux. L'auditeur interne a participé au groupe de travail dont il
contrôle ensuite le produit. Le domaine du chiffrage, où il siégeait, est aussi le plus mal coté, ce qui plaide
pour sa bonne foi. La règle demeure~: on n'audite pas son propre travail. Le rapport le signale en première
page.
\end{boxlecture}

\begin{boxmetier}{Auditeur interne}
Arrêter le référentiel avant les pièces. Publier les NEK. Se récuser s'il a participé aux travaux.
\end{boxmetier}

\begin{boxmetier}{Président du conseil d'administration}
Demander un avis indépendant avant d'approuver. Lire les NEK avant les NON.
\end{boxmetier}

\begin{boxmetier}{Consultant}
Livrer avec le document la table des pièces qui prouvent chaque exigence.
\end{boxmetier}

\begin{boxmetier}{Chercheur}
Employer la grille sur un échantillon d'organisations pour produire les mesures qui manquent.
\end{boxmetier}

\begin{boxresume}
\begin{itemize}
    \item Trois moments d'audit~: avant l'approbation, pendant l'exécution, à l'échéance.
    \item Un audit compare un objet à un critère au moyen d'une preuve.
    \item Trois cotes~: OUI, NON, NEK. Le NEK est de preuve ou de critère.
    \item $\mu$ mesure le prouvé, $U$ l'ignoré, $H$ la dispersion des cotes.
    \item Les résultats se lisent par domaine.
    \item L'indice d'alignement prouvé complète la mesure de perception de Luftman.
    \item Le rapport publie les NEK.
    \item On n'audite pas son propre travail.
\end{itemize}
\end{boxresume}

\begin{boxinfo}
\textbf{Points de contrôle.} Vous avez compris ce chapitre si vous savez~:
\begin{itemize}
    \item écrire dix exigences vérifiables~;
    \item coter un document et calculer les trois mesures~;
    \item rédiger un constat et sa recommandation~;
    \item donner un avis en trois phrases.
\end{itemize}
\end{boxinfo}

\begin{boxexo}[title={\ding{45}\quad Exercice~16.1. Trois mesures \niveaubase}]
Un audit de 20~exigences donne 12~OUI, 5~NON et 3~NEK. Calculez $\mu$, $U$ et $H$.
\end{boxexo}

\begin{boxexo}[title={\ding{45}\quad Exercice~16.2. Réécrire \niveaumoyen}]
Réécrivez chaque exigence pour qu'une pièce puisse la vérifier~: « la gouvernance est efficace », « les utilisateurs sont impliqués », « le budget est réaliste », « la sécurité est prise en compte ».
\end{boxexo}

\begin{boxexo}[title={\ding{45}\quad Exercice~16.3. Deux objets \niveaumoyen}]
Deux schémas directeurs obtiennent $\mu = 0{,}60$. Le premier a $U = 0{,}05$, le second $U = 0{,}35$. Rédigez pour chacun l'avis en trois phrases.
\end{boxexo}

\begin{boxexo}[title={\ding{45}\quad Exercice~16.4. Vérifier le tableau \niveaumoyen}]
Recalculez le $\mu$ de chaque domaine du Tableau~16.3 et l'entropie de l'ensemble. Quel domaine a l'entropie la plus forte~? Qu'est-ce que cela signifie~?
\end{boxexo}

\begin{boxexo}[title={\ding{45}\quad Exercice~16.5. L'audit à l'échéance \niveauexpert}]
Nous sommes en 2030. Concevez l'audit à l'échéance du schéma directeur 2025-2029 de l'ANPP~: douze exigences, les pièces à demander, les mesures de départ dont vous aurez besoin et que vous devez donc faire conserver dès 2025.
\end{boxexo}

\bigskip

\begin{boxinfo}
\textbf{Pour aller plus loin.}
\begin{itemize}
    \item Minka, \textit{Auditologie}~\cite{minka2026auditologie}.
    \item Luftman~\cite{luftman2000}.
    \item Henderson et Venkatraman~\cite{henderson1993}.
    \item Shannon~\cite{shannon1948}.
\end{itemize}
\end{boxinfo}

\begin{boxinfo}
\textbf{Transition.} Le cours s'achève où il a commencé~: devant un document de schéma directeur, avec une grille et des pièces. Les annexes rassemblent les outils~: glossaire, modèles, textes, référentiel d'exigences.
\end{boxinfo}


\appendix

    \chapter{Glossaire}
\label{ann:glossaire}

Ce glossaire reprend sous une forme abrégée les définitions des chapitres. Le renvoi indique le chapitre où le terme est défini et employé.

\begin{description}
    \item[Alignement stratégique] Accord entre la stratégie de l'organisation, celle de son système d'information et leurs organisations respectives (Ch.~\ref{chap:alignement}).
    \item[Analyse des écarts] Comparaison, capacité par capacité, de l'existant et de la cible (Ch.~\ref{chap:cible}).
    \item[Autorisation d'engagement] Limite supérieure des dépenses pouvant être engagées, y compris sur plusieurs années (Ch.~\ref{chap:droit}).
    \item[Bénéfice] Avantage mesurable ou observable obtenu par le métier grâce à un projet. Il a un propriétaire et une mesure de départ (Ch.~\ref{chap:piloter}).
    \item[Bilan d'impact sur l'activité] Étude qui fixe, pour chaque activité, l'interruption et la perte de données admissibles (Ch.~\ref{chap:maitriser}).
    \item[Cadrage] Phase qui répond par écrit aux six questions préalables et se clôt par une note signée (Ch.~\ref{chap:cadrer}).
    \item[Capacité] Aptitude attendue du système d'information, formulée par un verbe et un résultat, sans nommer de solution (Ch.~\ref{chap:cible}).
    \item[Carte d'applicabilité] Tableau daté qui dit, texte par texte, s'il s'applique à l'organisation et sous quelle condition (Ch.~\ref{chap:droit}).
    \item[Cellule de portefeuille] Petite équipe indépendante des projets qui tient les fiches, mesure l'avancement et alerte (Ch.~\ref{chap:piloter}).
    \item[Cible] Description du système d'information voulu à la fin de la période (Ch.~\ref{chap:cible}).
    \item[Comité de pilotage] Instance qui décide pendant la démarche, présidée par celui qui arbitre les budgets (Ch.~\ref{chap:cadrer}).
    \item[Comité de portefeuille] Instance qui lance, arrête et arbitre les projets pendant l'exécution (Ch.~\ref{chap:piloter}).
    \item[Constat] Énoncé en trois temps~: le fait, la pièce, la conséquence pour le métier (Ch.~\ref{chap:diagnostic}).
    \item[Coût complet de possession] Somme de l'investissement, du fonctionnement, du renouvellement et de la sortie sur une période (Ch.~\ref{chap:financer}).
    \item[Crédit de paiement] Limite supérieure des dépenses pouvant être payées dans l'année (Ch.~\ref{chap:droit}).
    \item[Donnée de référence] Donnée partagée par plusieurs processus, dotée d'un propriétaire et d'un système maître (Ch.~\ref{chap:donnees}).
    \item[Durée maximale d'interruption admissible] Délai au-delà duquel l'arrêt d'une activité cause un dommage inacceptable (Ch.~\ref{chap:maitriser}).
    \item[Exigence] Énoncé vérifiable par une pièce, tiré d'un texte, d'une norme, d'un engagement ou d'une règle de l'art (Ch.~\ref{chap:audit}).
    \item[Facteur clé de succès] Domaine où les résultats doivent être bons pour que l'organisation réussisse (Ch.~\ref{chap:alignement}).
    \item[Fiche d'application] Description d'une application en dix rubriques (Ch.~\ref{chap:diagnostic}).
    \item[Fiche de projet] Description d'un projet candidat en une page et huit rubriques (Ch.~\ref{chap:portefeuille}).
    \item[Gouvernance du système d'information] Système par lequel l'usage du numérique est dirigé et contrôlé. Trois tâches~: évaluer, diriger, surveiller (Ch.~\ref{chap:normes}).
    \item[Indice d'alignement prouvé] Part des exigences d'alignement dont la satisfaction est prouvée par une pièce (Ch.~\ref{chap:audit}).
    \item[Informatique hors contrôle] Outils employés sans que la fonction informatique les ait fournis ni recensés (Ch.~\ref{chap:diagnostic}).
    \item[Interopérabilité] Aptitude de systèmes distincts à échanger des données et à les employer (Ch.~\ref{chap:donnees}).
    \item[Lot] Partie d'un grand projet qui livre un résultat utilisable seul (Ch.~\ref{chap:portefeuille}).
    \item[NEK] Cote donnée quand l'auditeur ne peut pas conclure. De preuve si la pièce manque, de critère si l'exigence ne peut pas s'appliquer (Ch.~\ref{chap:audit}).
    \item[Note de cadrage] Document de dix pages, signé du commanditaire, qui fixe la démarche (Ch.~\ref{chap:cadrer}).
    \item[Perte de données maximale admissible] Ancienneté maximale des données que l'on accepte de retrouver après un sinistre (Ch.~\ref{chap:maitriser}).
    \item[Plan fonctionnel] Découpage du système en zones, quartiers et blocs, indépendant des logiciels (Ch.~\ref{chap:cible}).
    \item[Portefeuille] Ensemble des projets réunis pour atteindre des objectifs stratégiques (Ch.~\ref{chap:portefeuille}).
    \item[Principe directeur] Règle durable qui tranche d'avance une famille de décisions. Elle vaut par ses conséquences (Ch.~\ref{chap:cible}).
    \item[Projet obligatoire] Projet imposé par un texte, une injonction ou un risque que l'organisation ne peut pas porter (Ch.~\ref{chap:portefeuille}).
    \item[Provision pour aléas] Montant réservé pour couvrir l'incertitude des estimations (Ch.~\ref{chap:financer}).
    \item[Registre de conformité] Liste des exigences applicables avec leur responsable, leur état et leur pièce (Ch.~\ref{chap:maitriser}).
    \item[Réversibilité] Obligation faite à un prestataire de restituer les données et d'assister la reprise en fin de contrat (Ch.~\ref{chap:financer}).
    \item[Scénario de repli] Portefeuille réduit prévu d'avance pour le cas où une ressource manquerait (Ch.~\ref{chap:financer}).
    \item[Schéma directeur] Document approuvé par l'organe de gouvernance qui fixe pour plusieurs années la cible du système d'information, les projets, leur ordre, leur coût et leur pilotage (Ch.~\ref{chap:vocabulaire}).
    \item[Schéma directeur glissant] Schéma directeur dont l'horizon est repoussé d'un an à chaque révision annuelle (Ch.~\ref{chap:piloter}).
    \item[Stratégie numérique] Choix par lesquels une organisation emploie le numérique pour atteindre ses objectifs ou les redéfinir (Ch.~\ref{chap:vocabulaire}).
    \item[Trajectoire] Ordre daté des projets retenus, lisible année par année (Ch.~\ref{chap:portefeuille}).
    \item[Urbanisation] Démarche qui organise le système d'information en blocs aux frontières stables, à la manière d'un plan de ville (Ch.~\ref{chap:cible}).
    \item[Valeur acquise] Budget des travaux réellement achevés à une date. Elle se compare à la valeur planifiée et au coût réel (Ch.~\ref{chap:piloter}).
    \item[Valeur actuelle nette] Somme des flux d'un projet ramenés à la date de départ, investissement déduit (Ch.~\ref{chap:financer}).
\end{description}

    \providecommand{\gardetitre}{\par\vskip 0pt plus 9\baselineskip\penalty-300\vskip 0pt plus -9\baselineskip}
\chapter{Modèles}
\label{ann:modeles}

Ces modèles sont ceux de la mission ANPP. Ils se recopient et s'adaptent. Aucun ne dépasse une page une fois
rempli, sauf la note de cadrage et le schéma directeur lui-même.

\gardetitre
\section*{B1. Note de cadrage (Chapitre~\ref{chap:cadrer})}

\begin{tableaucours}{cL{4.0cm}L{7.4cm}}{%
    \tableauHeader{N°} &
    \tableauHeader{Section} &
    \tableauHeader{Contenu}
}
1 & Fait déclencheur & Trois lignes~: ce qui s'est passé, à quelle date \\
2 & Commanditaire et approbation & Qui commande, quel organe approuve, par quel acte \\
3 & Périmètre & Entités, métiers, objets. Liste des exclusions datées et motivées \\
4 & Horizon et ambition & Années couvertes. Remettre à niveau, aligner ou transformer \\
5 & Instances & Composition nominative, présidence, rythme \\
6 & Parties prenantes & Grille pouvoir-intérêt, conduite pour chaque groupe \\
7 & Responsabilités & Matrice par livrable, un seul approbateur \\
8 & Calendrier et livrables & Phases, dates, comités, plafond de pages \\
9 & Communication & À qui, quoi, quand, par quel moyen \\
10 & Signature & Commanditaire, date \\
\end{tableaucours}

\gardetitre
\section*{B2. Fiche d'application (Chapitre~\ref{chap:diagnostic})}

\begin{tableaucours}{L{4.4cm}L{7.6cm}}{%
    \tableauHeader{Rubrique} &
    \tableauHeader{Exemple~: liquidation des droits}
}
Nom et métier servi & LIQUID, direction des prestations \\
Responsable côté métier & Directeur des prestations \\
Utilisateurs & 64 agents, siège et antennes \\
Éditeur ou auteur & Développement interne \\
Mise en service & 2004 \\
Technologie & Langage et base qui ne sont plus maintenus \\
Maintenance & Deux agents de la division \\
Coût annuel & 38~millions, salaires compris \\
Données personnelles & Oui~: identité, carrière, droits \\
Notes & Valeur 4,5. Santé 1,5. Décision~: refondre \\
\end{tableaucours}

\gardetitre
\section*{B3. Fiche de projet (Chapitre~\ref{chap:portefeuille})}

\begin{tableaucours}{L{4.4cm}L{7.6cm}}{%
    \tableauHeader{Rubrique} &
    \tableauHeader{Exemple~: P07, référentiel des bénéficiaires}
}
Résultat & Un fichier unique des 186~000 bénéficiaires, maître pour toutes les applications \\
Objectif servi & Ne payer que les ayants droit \\
Responsable & Directeur des prestations \\
Investissement & 340~millions (60 en 2027, 280 en 2028) \\
Coût annuel ensuite & 48~millions à partir de 2029 \\
Durée & 20~mois, passation comprise \\
Dépendances & Aucune en amont. Les trois lots de la refonte en dépendent \\
Ce qu'il remplace & Onze fichiers de bénéficiaires \\
Risque principal & Reprise des données défaillante. Parade~: fiabilisation préalable \\
Sécurité, continuité, données & Analyse de risque, reprise sur le site de secours, formalités de la loi n°~2024/017 \\
Jalon d'arrêt & Fin de la fiabilisation~: moins de 1~\% de doublons, sinon report \\
\end{tableaucours}

\gardetitre
\section*{B4. Fiche d'indicateur (Chapitre~\ref{chap:piloter})}

\begin{tableaucours}{L{4.4cm}L{7.6cm}}{%
    \tableauHeader{Rubrique} &
    \tableauHeader{Exemple~: paiements à date}
}
Formule & Prestations payées au plus tard à la date prévue, divisées par les prestations dues \\
Source & Journal des paiements \\
Responsable de la mesure & Cellule de portefeuille \\
Fréquence & Chaque mois \\
Mesure de départ & 91~\% en 2026 \\
Cible & 99~\% en 2031 \\
\end{tableaucours}

\gardetitre
\section*{B5. Sommaire type d'un schéma directeur}

Soixante pages au plus, hors annexes. Une synthèse de quatre pages l'accompagne pour l'organe qui approuve.

\begin{tableaucours}{cL{5.0cm}cL{4.6cm}}{%
    \tableauHeader{N°} &
    \tableauHeader{Partie} &
    \tableauHeader{Pages} &
    \tableauHeader{Chapitre du cours}
}
1 & Synthèse pour décision & 4 & Tous \\
2 & Mandat, périmètre, gouvernance de la démarche & 4 & \ref{chap:cadrer} \\
3 & Stratégie de l'organisation et cadre de référence & 5 & \ref{chap:cameroun} et \ref{chap:droit} \\
4 & Diagnostic~: dix constats & 8 & \ref{chap:diagnostic} \\
5 & Cible~: capacités, principes, plan fonctionnel, scénario retenu & 10 & \ref{chap:cible} \\
6 & Portefeuille et trajectoire & 10 & \ref{chap:portefeuille} \\
7 & Financement, achats, dotation & 8 & \ref{chap:financer} \\
8 & Pilotage et révision & 5 & \ref{chap:piloter} \\
9 & Risques, sécurité, continuité, conformité & 6 & \ref{chap:maitriser} \\
 & Annexes~: fiches, inventaire, registres, projets écartés &  &  \\
\end{tableaucours}

\gardetitre
\section*{B6. Page de décision pour un comité}

Chaque comité reçoit trois jours avant une page par décision. Elle contient cinq lignes~: la décision demandée,
les faits et leurs pièces, les options avec leur coût, la recommandation, ce qui se passe si l'on ne décide pas.

    \providecommand{\gardetitre}{\par\vskip 0pt plus 9\baselineskip\penalty-300\vskip 0pt plus -9\baselineskip}
\chapter{Textes et Documents de Référence}
\label{ann:textes}

Cette annexe liste les textes cités dans le cours, par famille. Chaque ligne renvoie à la bibliographie, qui
donne la référence complète. La liste est arrêtée au 13~novembre~2026. Les textes changent~: la carte
d'applicabilité d'une mission se date et se relit (Chapitre~\ref{chap:droit}).

\gardetitre
\section*{C1. Lois camerounaises}

\begin{tableaucours}{L{6.4cm}L{5.6cm}}{%
    \tableauHeader{Texte} &
    \tableauHeader{Intérêt pour le schéma directeur}
}
Loi n°~2010/012 du 21~décembre~2010, cybersécurité et cybercriminalité~\cite{loi2010012} & Protection des systèmes, audit de sécurité obligatoire \\
Loi n°~2010/013 du 21~décembre~2010, communications électroniques~\cite{loi2010013} & Réseaux et services \\
Loi n°~2010/021 du 21~décembre~2010, commerce électronique~\cite{loi2010021} & Activités commerciales en ligne \\
Loi-cadre n°~2011/012 du 6~mai~2011, protection du consommateur~\cite{loi2011012} & Services aux particuliers \\
Loi n°~2017/010 du 12~juillet~2017, établissements publics~\cite{loi2017010} & Organe qui approuve \\
Loi n°~2017/011 du 12~juillet~2017, entreprises publiques~\cite{loi2017011} & Organe qui approuve \\
Loi n°~2018/011 du 11~juillet~2018, transparence et bonne gouvernance~\cite{loi2018011} & Reddition des comptes \\
Loi n°~2018/012 du 11~juillet~2018, régime financier de l'État~\cite{loi2018012} & Programmes, engagements, paiements \\
Loi n°~2023/009 du 25~juillet~2023, protection des enfants en ligne~\cite{loi2023009} & Services ouverts au public \\
Loi n°~2024/017 du 23~décembre~2024, données personnelles~\cite{loi2024017} & Formalités, sécurité, transferts \\
\end{tableaucours}

\gardetitre
\section*{C2. Décrets camerounais}

\begin{tableaucours}{L{6.4cm}L{5.6cm}}{%
    \tableauHeader{Texte} &
    \tableauHeader{Intérêt pour le schéma directeur}
}
Décret n°~2011/1521/PM du 15~juin~2011~\cite{decret20111521} & Application de la loi sur le commerce électronique \\
Décret n°~2012/180 du 10~avril~2012~\cite{decret2012180} & Ancien texte de l'ANTIC \\
Décret n°~2012/1643/PM~\cite{decret20121643} & Audit de sécurité obligatoire \\
Décret n°~2018/355 du 12~juin~2018~\cite{decret2018355} & Marchés des entreprises publiques \\
Décret n°~2018/366 du 20~juin~2018~\cite{decret2018366} & Code des marchés publics \\
Décret n°~2019/150 du 22~mars~2019~\cite{decret2019150} & Organisation de l'ANTIC \\
\end{tableaucours}

\gardetitre
\section*{C3. Documents d'orientation du Cameroun}

\begin{tableaucours}{L{6.4cm}L{5.6cm}}{%
    \tableauHeader{Texte} &
    \tableauHeader{Intérêt pour le schéma directeur}
}
Vision 2035~\cite{vision2035} & Cap de long terme \\
Stratégie nationale de développement 2020-2030~\cite{snd30} & Document de rattachement des administrations \\
Plan stratégique Cameroun numérique 2020~\cite{camnum2020} & Échu. Référence historique \\
Stratégie nationale d'intelligence artificielle~\cite{snia2025} & Sept piliers, présentée en juillet 2025 \\
Plan stratégique national de santé numérique 2026-2030~\cite{psnsn2026} & Exemple de stratégie sectorielle \\
Schéma directeur informatique 2020-2024 du MINEPDED~\cite{minepded2020sdi} & Exemple de schéma directeur ministériel \\
\end{tableaucours}

\begin{boxinfo}
\textbf{Le plan 2026-2030.} Le Plan de développement de l'économie numérique au Cameroun 2026-2030, édition
2025~\cite{pden2026}, est finalisé et déposé. Il est en cours de validation dans les services du Premier
ministre et n'est donc pas en cours de mise en œuvre. Il ne figure pas dans la liste des documents en vigueur.
Un schéma directeur le mentionne comme attendu (Chapitre~\ref{chap:cameroun}, section~\ref{sec:c06-projet}).
\end{boxinfo}

\gardetitre
\section*{C4. Textes de la CEMAC, de la COBAC et de l'OHADA}

\begin{tableaucours}{L{6.4cm}L{5.6cm}}{%
    \tableauHeader{Texte} &
    \tableauHeader{Intérêt pour le schéma directeur}
}
Règlement n°~04/08/CEMAC/UMAC/COBAC du 6~octobre~2008~\cite{cemac200804} & Gouvernement d'entreprise des établissements de crédit \\
Règlement COBAC R-2008/01~\cite{cobac200801} & Plan de continuité des activités \\
Règlement COBAC R-2016/04~\cite{cobac201604} & Contrôle interne \\
Règlement n°~03/16/CEMAC/UMAC/CM du 21~décembre~2016~\cite{cemac2016paiement} & Systèmes, moyens et incidents de paiement \\
Règlement n°~04/18/CEMAC/UMAC/COBAC du 21~décembre~2018~\cite{cemac2018paiement} & Services de paiement \\
Lettre circulaire LC-COB/04 du 21~janvier~2022~\cite{cobac2022lc} & Audit et politique de sécurité \\
Rapport annuel 2024 de la COBAC~\cite{cobac2024rapport} & Règlement sur le risque informatique \\
Acte uniforme OHADA du 26~janvier~2017~\cite{ohada2017audcif} & Droit comptable, conservation dix ans \\
\end{tableaucours}

\gardetitre
\section*{C5. Textes africains et internationaux}

\begin{tableaucours}{L{6.4cm}L{5.6cm}}{%
    \tableauHeader{Texte} &
    \tableauHeader{Intérêt pour le schéma directeur}
}
Convention de Malabo, 27~juin~2014~\cite{malabo2014} & Cybersécurité et données personnelles \\
Stratégie de transformation numérique pour l'Afrique 2020-2030~\cite{ua2020dts} & Cap continental \\
Cadre stratégique de l'Union africaine sur les données, 2022~\cite{ua2022donnees} & Gouvernance des données \\
Stratégie continentale sur l'intelligence artificielle, 2024~\cite{ua2024ia} & Usage responsable \\
Pacte numérique africain, 2024~\cite{ua2024pacte} & Position commune \\
Protocole de la ZLECAf sur le commerce numérique, 2024~\cite{zlecaf2024} & Échanges numériques \\
Pacte numérique mondial, A/RES/79/1, 22~septembre~2024~\cite{onu2024pacte} & Cadre des Nations unies \\
Recommandation de l'OCDE sur les stratégies numériques gouvernementales, 2014~\cite{ocde2014} & Doctrine de l'administration numérique \\
\end{tableaucours}

\gardetitre
\section*{C6. Normes et référentiels}

\begin{tableaucours}{L{6.4cm}L{5.6cm}}{%
    \tableauHeader{Texte} &
    \tableauHeader{Intérêt pour le schéma directeur}
}
ISO/IEC~38500:2024~\cite{iso38500} & Gouvernance du numérique \\
ISO~37000:2021~\cite{iso37000} & Gouvernance des organisations \\
COBIT~2019~\cite{cobit2019} & Objectifs de gouvernance et de gestion \\
TOGAF, 10\textsuperscript{e}~édition~\cite{togaf10} & Architecture d'entreprise \\
ISO/IEC/IEEE~42010~\cite{iso42010} & Description d'architecture \\
ITIL~4~\cite{itil4} & Gestion des services \\
ISO~21502 et ISO~21504~\cite{iso21502,iso21504} & Projets et portefeuilles \\
ISO~31000~\cite{iso31000} & Management du risque \\
ISO/IEC~27001~\cite{iso27001} & Sécurité de l'information \\
ISO~22301~\cite{iso22301} & Continuité d'activité \\
ISO/IEC~42001~\cite{iso42001} & Management de l'intelligence artificielle \\
ISO~56002~\cite{iso56002} & Management de l'innovation \\
EN~16234-1~\cite{ecf} & Compétences numériques \\
\end{tableaucours}


    \providecommand{\gardetitre}{\par\vskip 0pt plus 9\baselineskip\penalty-300\vskip 0pt plus -9\baselineskip}
\chapter{Les quarante-huit exigences de la mission ANPP}
\label{ann:exigences}

Cette annexe réunit le référentiel employé le 18~juin~2027 pour l'audit du schéma directeur 2027-2031 de l'ANPP
avant son approbation (Chapitre~\ref{chap:audit}). Elle permet de refaire chaque calcul~: 6 + 6 + 6 + 6 + 7 +
6 + 6 + 5 = 48 exigences, 37~\OUI{}, 7~\NON{}, 4~\NEK{}. La colonne Source indique d'où vient l'exigence~: un
texte, une norme, un engagement de l'Agence, une règle de l'art ou la cohérence interne du document. Les
quatre \NEK{} sont de preuve.

Les cotes sont celles du 18~juin~2027. Elles diffèrent parfois de celles des encadrés de lecture des chapitres,
établies à d'autres dates sur des livrables intermédiaires.

\gardetitre
\section*{D1. Cadrage (Chapitre~\ref{chap:cadrer})}

\begin{tableaucours}{cL{4.9cm}L{1.8cm}cL{3.3cm}}{%
    \tableauHeader{N°} &
    \tableauHeader{Exigence} &
    \tableauHeader{Source} &
    \tableauHeader{Cote} &
    \tableauHeader{Preuve, nature du NEK}
}
1 & La note de cadrage est signée par le commanditaire & Règle de l'art & \OUI{} & Signature du 29~janvier \\
2 & Le périmètre du diagnostic est celui de la note de cadrage & Cohérence & \OUI{} & Rapprochement des deux documents \\
3 & L'organe d'approbation est le conseil d'administration & Texte & \OUI{} & Section 2 de la note \\
4 & Le comité de pilotage est présidé par la direction générale & Norme & \OUI{} & Composition du comité \\
5 & L'horizon couvre le plan stratégique de l'Agence & Engagement & \OUI{} & 2027-2031 pour un plan 2026-2030 \\
6 & La tutelle a été informée du lancement & Règle de l'art & \NEK{} & Copie du courrier non remise. \textit{Preuve} \\
\end{tableaucours}

\gardetitre
\section*{D2. Diagnostic (Chapitre~\ref{chap:diagnostic})}

\begin{tableaucours}{cL{4.9cm}L{1.8cm}cL{3.3cm}}{%
    \tableauHeader{N°} &
    \tableauHeader{Exigence} &
    \tableauHeader{Source} &
    \tableauHeader{Cote} &
    \tableauHeader{Preuve, nature du NEK}
}
7 & L'inventaire est construit depuis les métiers & Règle de l'art & \OUI{} & 47~fiches, dont 16 hors liste officielle \\
8 & Chaque application a un responsable côté métier & Règle de l'art & \NON{} & 19~fiches sans responsable \\
9 & Chaque note renvoie à une pièce & Règle de l'art & \OUI{} & Registre des pièces \\
10 & La dépense de cinq exercices est reconstituée & Règle de l'art & \OUI{} & Tableau 9.5 \\
11 & Le total de la dépense égale la somme par nature & Cohérence & \OUI{} & 2~418 + 682 = 3~100 \\
12 & Le rapport est approuvé par le comité de pilotage & Règle de l'art & \OUI{} & Procès-verbal du 19~mars \\
\end{tableaucours}

\gardetitre
\section*{D3. Cible (Chapitre~\ref{chap:cible})}

\begin{tableaucours}{cL{4.9cm}L{1.8cm}cL{3.3cm}}{%
    \tableauHeader{N°} &
    \tableauHeader{Exigence} &
    \tableauHeader{Source} &
    \tableauHeader{Cote} &
    \tableauHeader{Preuve, nature du NEK}
}
13 & Chaque capacité se rattache à un objectif du plan stratégique & Engagement & \OUI{} & Tableau 10.1 \\
14 & Chaque objectif est servi par une capacité et par un projet au moins & Engagement & \OUI{} & Table de correspondance \\
15 & Les principes directeurs sont approuvés avec leurs conséquences & Norme & \OUI{} & Procès-verbal du 9~avril \\
16 & L'architecture est décrite par des vues adaptées aux lecteurs & Norme & \OUI{} & Quatre plans \\
17 & Toutes les applications de l'inventaire sont situées sur le plan fonctionnel & Cohérence & \OUI{} & 47 sur 47 \\
18 & Deux scénarios éligibles ont été comparés sur des critères fixés d'avance & Règle de l'art & \OUI{} & Tableau 10.6 \\
\end{tableaucours}

\gardetitre
\section*{D4. Portefeuille (Chapitre~\ref{chap:portefeuille})}

\begin{tableaucours}{cL{4.9cm}L{1.8cm}cL{3.3cm}}{%
    \tableauHeader{N°} &
    \tableauHeader{Exigence} &
    \tableauHeader{Source} &
    \tableauHeader{Cote} &
    \tableauHeader{Preuve, nature du NEK}
}
19 & Chaque projet a une fiche et un responsable dans le métier & Règle de l'art & \OUI{} & 22~fiches \\
20 & Chaque projet retenu cite un objectif du plan stratégique & Engagement & \OUI{} & 22~fiches \\
21 & La sélection suit des critères écrits et les projets écartés sont motivés & Norme & \OUI{} & Annexe du portefeuille \\
22 & Le total des projets retenus tient dans l'enveloppe fixée & Engagement & \OUI{} & 4~180 pour 4~200 \\
23 & La trajectoire respecte les dépendances & Cohérence & \OUI{} & Tableau 11.4 \\
24 & Chaque projet nomme ce qu'il remplace & Règle de l'art & \NON{} & Sept fiches muettes \\
\end{tableaucours}

\gardetitre
\section*{D5. Chiffrage et financement (Chapitre~\ref{chap:financer})}

\begin{tableaucours}{cL{4.9cm}L{1.8cm}cL{3.3cm}}{%
    \tableauHeader{N°} &
    \tableauHeader{Exigence} &
    \tableauHeader{Source} &
    \tableauHeader{Cote} &
    \tableauHeader{Preuve, nature du NEK}
}
25 & Chaque projet affiche son investissement et son coût annuel & Règle de l'art & \OUI{} & 22~fiches complétées \\
26 & Une provision pour aléas est chiffrée & Règle de l'art & \OUI{} & 242~millions \\
27 & Chaque projet est présenté en engagements et en paiements & Texte & \OUI{} & Tableaux du plan \\
28 & Le total des fiches égale le total du plan de financement & Cohérence & \NON{} & 4~180 d'un côté, 4~422 de l'autre \\
29 & Chaque bénéfice a un propriétaire et une mesure de départ & Règle de l'art & \NON{} & Deux bénéfices sans mesure \\
30 & Le fonctionnement à terme est porté à la connaissance du conseil & Engagement & \NON{} & Absent de la synthèse de quatre pages \\
31 & Chaque ressource est qualifiée par un acte & Règle de l'art & \NEK{} & Projet de convention non transmis. \textit{Preuve} \\
\end{tableaucours}

\gardetitre
\section*{D6. Pilotage (Chapitre~\ref{chap:piloter})}

\begin{tableaucours}{cL{4.9cm}L{1.8cm}cL{3.3cm}}{%
    \tableauHeader{N°} &
    \tableauHeader{Exigence} &
    \tableauHeader{Source} &
    \tableauHeader{Cote} &
    \tableauHeader{Preuve, nature du NEK}
}
32 & L'organe de gouvernance surveille l'exécution chaque année & Norme & \OUI{} & Revue annuelle inscrite au document \\
33 & Un comité de portefeuille trimestriel est institué & Règle de l'art & \OUI{} & Section 8 \\
34 & Une méthode unique de mesure de l'avancement est fixée & Norme & \OUI{} & Valeur acquise \\
35 & Le tableau de bord contient un indicateur de résultat du plan stratégique & Engagement & \OUI{} & Paiements à date \\
36 & Une procédure traite les demandes hors portefeuille & Règle de l'art & \OUI{} & Section 8 \\
37 & La révision annuelle précède la préparation du budget & Règle de l'art & \NEK{} & Calendrier non arrêté. \textit{Preuve} \\
\end{tableaucours}

\gardetitre
\section*{D7. Maîtrise des risques (Chapitre~\ref{chap:maitriser})}

\begin{tableaucours}{cL{4.9cm}L{1.8cm}cL{3.3cm}}{%
    \tableauHeader{N°} &
    \tableauHeader{Exigence} &
    \tableauHeader{Source} &
    \tableauHeader{Cote} &
    \tableauHeader{Preuve, nature du NEK}
}
38 & Un registre des risques existe avec des propriétaires & Norme & \OUI{} & Registre du 11~juin \\
39 & Le secours est dimensionné par un bilan d'impact sur l'activité & Norme & \OUI{} & Tableau 14.2 \\
40 & L'audit de sécurité obligatoire est programmé & Texte & \OUI{} & Projet obligatoire n°~2 \\
41 & Les traitements de données personnelles sont recensés et leurs formalités programmées & Texte & \OUI{} & Projet obligatoire n°~3 \\
42 & Chaque fiche porte les rubriques sécurité, continuité, données & Règle de l'art & \NON{} & Rubriques absentes de huit fiches \\
43 & Chaque projet a un jalon d'arrêt & Règle de l'art & \NON{} & Neuf fiches sans jalon \\
\end{tableaucours}

\gardetitre
\section*{D8. Cadre juridique (Chapitre~\ref{chap:droit})}

\begin{tableaucours}{cL{4.9cm}L{1.8cm}cL{3.3cm}}{%
    \tableauHeader{N°} &
    \tableauHeader{Exigence} &
    \tableauHeader{Source} &
    \tableauHeader{Cote} &
    \tableauHeader{Preuve, nature du NEK}
}
44 & Une carte d'applicabilité datée est jointe & Texte & \OUI{} & Annexe du document \\
45 & Aucun produit ni aucune marque n'est nommé & Texte & \OUI{} & Relecture \\
46 & Les délais de passation des marchés sont comptés dans la trajectoire & Texte & \OUI{} & Plan d'achat \\
47 & Chaque projet est rattaché à un programme budgétaire & Texte & \OUI{} & 22~fiches \\
48 & Les transferts de données hors du Cameroun sont identifiés & Texte & \NEK{} & Contrats d'hébergement non transmis. \textit{Preuve} \\
\end{tableaucours}

\gardetitre
\section*{Récapitulatif}

\begin{tableaucours}{L{4.4cm}ccccc}{%
    \tableauHeader{Domaine} &
    \tableauHeader{Exigences} &
    \tableauHeader{\OUI{}} &
    \tableauHeader{\NON{}} &
    \tableauHeader{\NEK{}} &
    \tableauHeader{$\mu$}
}
Cadrage & 6 & 5 & 0 & 1 & 0{,}83 \\
Diagnostic & 6 & 5 & 1 & 0 & 0{,}83 \\
Cible & 6 & 6 & 0 & 0 & 1{,}00 \\
Portefeuille & 6 & 5 & 1 & 0 & 0{,}83 \\
Chiffrage et financement & 7 & 3 & 3 & 1 & 0{,}43 \\
Pilotage & 6 & 5 & 0 & 1 & 0{,}83 \\
Maîtrise des risques & 6 & 4 & 2 & 0 & 0{,}67 \\
Cadre juridique & 5 & 4 & 0 & 1 & 0{,}80 \\
Ensemble & 48 & 37 & 7 & 4 & 0,77 \\
\end{tableaucours}

\hfill\lecture{0{,}771}{0{,}083}{0{,}99}


\backmatter

    \nocite{*}
    \printbibliography[heading=bibintoc, title={Références bibliographiques}]

    \printindex

% ==============================================================================
% 4E DE COUVERTURE
% ==============================================================================
\clearpage
\thispagestyle{empty}
\pagecolor{CouleurPrimaire}

\begin{tikzpicture}[remember picture, overlay]
    \fill[CouleurAccent]
        (current page.south west) rectangle
        ([yshift=1.8cm] current page.south east);
    \fill[CouleurAccent]
        ([yshift=-1.4cm] current page.north west) rectangle
        ([yshift=-1.7cm] current page.north east);
\end{tikzpicture}

\vspace*{1.8cm}

\begin{center}
    {\color{CouleurAccent}\rule{0.6\linewidth}{1pt}}\\[1.5ex]
    {\color{white}\Large\bfseries À propos de ce cours}\\[1.5ex]
    {\color{CouleurAccent}\rule{0.6\linewidth}{1pt}}
\end{center}

\vspace{1cm}

\begin{center}
\begin{minipage}{0.75\textwidth}
\color{white!85!gray}\setstretch{1.4}
Beaucoup d'organisations possèdent un schéma directeur. Peu l'exécutent. Ce cours apprend à en
produire un qui tienne~: partir de la stratégie de l'organisation, regarder l'existant sans complaisance,
fixer une cible, choisir des projets, les ordonner, les financer, puis piloter et réviser.

\medskip

Il progresse en six parties, précédées d'un prologue~: les fondements théoriques, le contexte international
(ISO/IEC~38500, COBIT, TOGAF, stratégies de l'Union africaine), le contexte national (orientations, finances
publiques, marchés, textes de la CEMAC et de la COBAC), la méthode en cinq temps, la mise en œuvre, les défis
actuels et l'audit. Une mission suivie de bout en bout sert de fil conducteur.

\medskip

Il s'adresse aux étudiants du Master Professionnel en Management des Systèmes d'Information et aux
professionnels qui exercent comme directeurs des systèmes d'information, consultants, responsables de
portefeuille, contrôleurs de gestion ou auditeurs.
\end{minipage}
\end{center}

\vfill

\begin{center}
    {\color{CouleurAccent}\rule{0.7\linewidth}{1.2pt}}\\[1.2ex]
    {\color{white}\large\bfseries MINKA MI NGUIDJOÏ Thierry Emmanuel}\\[0.4ex]
    {\color{white!80!gray}\small
        CISA~·~CISM~·~CGEIT~·~CRISC~·~CDPSE~·~ISO 27001~·~ISO 22301~·~ISO 9001}\\[0.3ex]
    {\color{white!70!gray}\footnotesize\itshape
        Chercheur au laboratoire d'Ingénierie Mathématiques et des Systèmes d'Information}\\[0.3ex]
    {\color{white!70!gray}\footnotesize\itshape
        Chercheur~|~Cryptologue~|~Auditeur Senior~|~Expert Judiciaire~|~GRC \& Sécurité de l'Information}\\[1.2ex]
    {\color{CouleurAccent}\rule{0.7\linewidth}{1.2pt}}
\end{center}

\vspace{1cm}

\afterpage{\nopagecolor}
\end{document}
